\documentclass[11pt,a4paper]{report}
\usepackage[T1]{fontenc}
\usepackage[utf8]{inputenc}
\usepackage[english]{babel}
\usepackage{amsmath,amssymb,mathtools,bm}
\usepackage{array,longtable,booktabs}
\usepackage{xcolor}
\usepackage[margin=2.5cm]{geometry}
\usepackage[colorlinks=true,linkcolor=blue!50!black,citecolor=green!40!black,urlcolor=blue!50!black]{hyperref}
\usepackage[numbers,sort&compress]{natbib}
\usepackage{xurl}
\newcommand{\script}[1]{\path{#1}}   % breakable file names

% ------------------------------------------------------------------ macros
\newcommand{\G}{\mathcal{G}}                 % gravitational constant
\newcommand{\I}{\mathrm{i}}                  % imaginary unit
\newcommand{\M}{\mathcal{M}}                 % mean anomaly
\newcommand{\dd}{\mathrm{d}}
\newcommand{\thetap}{\theta_{\mathrm p}}     % rotation angle of the planet
\newcommand{\Op}{\Omega_{\mathrm p}}         % spin rate of the planet
\newcommand{\PK}{P^{\mathrm K}}              % Kaula's associated Legendre function (no CS phase)
\newcommand{\kb}{\bar k}                     % complex Love number
\newcommand{\er}{\mathbf e_r}
\newcommand{\et}{\mathbf e_\theta}
\newcommand{\ep}{\mathbf e_\varphi}
\newcommand{\vr}{\mathbf r}
\newcommand{\vF}{\mathbf F}
% (\hat kept as the LaTeX default)
\renewcommand{\Re}{\operatorname{Re}}
\renewcommand{\Im}{\operatorname{Im}}
\newcommand{\sgn}{\operatorname{sgn}}
\newcommand{\Fsp}{\mathbf F_{\star\to\mathrm p}}
\newcommand{\Fps}{\mathbf F_{\mathrm p\to\star}}

\numberwithin{equation}{chapter}

\title{Tides in Kaula's formalism for an $N$-body code\\[1ex]
\large Derivation, verification, and implementation in Posidonius\\ (tide raised in a planet and in its star, degree two, arbitrary rheology)}
\author{}
\date{}

\begin{document}
\maketitle

\begin{abstract}
We derive, step by step, the tidal force exerted by the tidal bulge of a planet on its host star, and its reaction on the planet, for a planet of arbitrary rheology described by a frequency-dependent complex Love number. The tidal-raising potential of the star is expanded in spherical harmonics and transformed into Keplerian orbital elements with Kaula's transformation. The response of the planet is applied term by term. The gradient of the resulting tidal potential is computed in the local polar basis and in a body-fixed Cartesian (cyclic) basis. A second Kaula transformation at the field point, followed by the identification of the elements of the field point with those of the star, gives the force in orbital elements only. The coplanar case is derived as a special case.

The equations are then written in the form used by an $N$-body integrator: from the state vectors, in inertial axes, with the spin torque that conserves the total angular momentum, tidal frequencies built with mean rates, a rheology-independent truncation table, and the dissipated power. The tide raised in the star uses the same equations with the roles of the two bodies exchanged, and a Love number that depends on the sign of the frequency. The last chapter describes the implementation in the $N$-body code Posidonius.

Every formula of this document has been verified numerically, independently of its derivation, against the exact linear response of the deformed body computed without Kaula's expansion. The conventions are stated once in Appendix~\ref{app:conventions} and never changed afterwards.
\end{abstract}

\tableofcontents

\chapter{Introduction}\label{ch:intro}

\paragraph{Problem.} A star of mass $M^\star$ raises a tide on a planet of radius $R$. The deformed planet acts back on the star with the force $\Fps$, and the star acts on the planet with $\Fsp=-\Fps$. We derive this force, and the torque on the planet's spin, for an arbitrary rheology described by a frequency-dependent complex Love number $\kb_l(\sigma)$ \citep{EfroimskyMakarov2013,EfroimskyMakarov2014}. The same equations give the tide raised in the star by the planet.

\paragraph{Route.} The derivation follows Kaula's approach \citep{Kaula1964,Kaula2000}:
\begin{enumerate}
  \item expand the potential of the star in spherical harmonics (Chapter~\ref{ch:potential});
  \item express each term in the star's orbital elements with Kaula's transformation (Chapter~\ref{ch:kaula});
  \item apply the planet's response term by term to obtain the tidal potential $T$ (Chapter~\ref{ch:tidal});
  \item take its gradient (Chapter~\ref{ch:gradient});
  \item express the result in orbital elements only, by a second Kaula transformation at the field point (Chapter~\ref{ch:second});
  \item specialise to a coplanar orbit (Chapter~\ref{ch:coplanar});
  \item give the form of the equations for an $N$-body code (Chapter~\ref{ch:nbody});
  \item apply them to the tide raised in the star (Chapter~\ref{ch:star});
  \item describe their implementation in Posidonius (Chapter~\ref{ch:posidonius}).
\end{enumerate}

\paragraph{How to read this document.}
\begin{itemize}
  \item \textbf{To use the model} in Posidonius: Chapter~\ref{ch:nbody} for the equations and their limits (Sections~\ref{sec:io}, \ref{sec:rates}, \ref{sec:truncation}, \ref{sec:limits}), Chapter~\ref{ch:star} for the tide in the star, and Chapter~\ref{ch:posidonius} for the input spectra, the messages and the known limitations.
  \item \textbf{To check the implementation}: Chapter~\ref{ch:nbody} gives every equation in the order of the algorithm (Section~\ref{sec:algorithm}); the table of Section~\ref{sec:pos-files} gives the Rust function that implements each one. The code comments refer to the equation numbers of this document.
  \item \textbf{To follow the derivation}: Chapters~\ref{ch:potential} to~\ref{ch:coplanar}, in order. Each chapter starts from the result of the previous one; Section~\ref{sec:implementable} links the general gradient of Chapter~\ref{ch:gradient} to the implementable form.
\end{itemize}

\paragraph{Where to find the conventions.} Before reading the derivation, see:
\begin{itemize}
  \item \textbf{Appendix~\ref{app:conventions}} — physical hypotheses H1--H3, frames, orbital elements, Love-number convention, sign of the force, and a table of all symbols;
  \item \textbf{Appendix~\ref{app:harmonics}} — Legendre functions and spherical harmonics, including the phase convention, which is essential for the signs;
  \item \textbf{Appendix~\ref{app:kaula}} — Kaula's inclination and eccentricity functions, with exact tables.
\end{itemize}
The conventions are stated once and never changed.

\paragraph{Verification.} Every formula was checked numerically, independently of its derivation, against the exact linear response of the deformed body. This response is computed directly from the tidal-raising potential and the rheology, without Kaula's expansion. The agreement is at the level of rounding errors ($10^{-16}$ to $10^{-14}$), except where a truncation is stated. The checks of the implementation are listed in Section~\ref{sec:pos-tests}.

\chapter{The tidal-raising potential in spherical harmonics}\label{ch:potential}

\section{Multipole expansion}
The gravitational potential created by the star at the field point $\vr$ is
\begin{equation}
  U(\vr,\vr^\star)=-\frac{\G M^\star}{\|\vr^\star-\vr\|}. \label{eq:U}
\end{equation}
Let $\Theta$ be the angle between $\vr$ and $\vr^\star$. With unit vectors,
\begin{equation}
  \cos\Theta=\hat{\vr}\cdot\hat{\vr}^\star=\cos\theta\cos\theta^\star+\sin\theta\sin\theta^\star\cos(\varphi-\varphi^\star). \label{eq:cosTheta}
\end{equation}
Put $x=\cos\Theta$ and $t=r/r^\star$. Then
\[
  \|\vr^\star-\vr\|^2=r^{\star2}\,(1-2xt+t^2).
\]
The generating function of the Legendre polynomials,
\[
  (1-2xt+t^2)^{-1/2}=\sum_{l\ge0}t^lP_l(x),
\]
converges for $|t|<1$ and $|x|\le1$ \citep[e.g.][\S1.2]{Kaula2000}. Therefore
\begin{equation}
  U=-\frac{\G M^\star}{r^\star}\sum_{l=0}^{\infty}\Big(\frac{r}{r^\star}\Big)^lP_l(\cos\Theta),\qquad r<r^\star. \label{eq:U-legendre}
\end{equation}
The condition $r<r^\star$ always holds where the series is used: inside the planet and at its surface, $r\le R<r^\star$.

\section{Addition theorem}
\paragraph{Spherical harmonics.} We use the spherical harmonics of \citet[eq.~5.2(1)]{Varshalovich1988}:
\begin{equation}
  Y_l^m(\theta,\varphi)=N_{lm}\,P_l^m(\cos\theta)\,e^{\I m\varphi},\qquad N_{lm}=\sqrt{\frac{2l+1}{4\pi}\frac{(l-m)!}{(l+m)!}}, \label{eq:Ylm-main}
\end{equation}
where $P_l^m$ is the associated Legendre function with the Condon--Shortley phase. Both $N_{lm}$ and $P_l^m$ are real. These functions are orthonormal on the unit sphere. Their definitions, conventions and properties are given in Appendix~\ref{app:harmonics}.

\paragraph{Addition theorem.} The addition theorem for spherical harmonics \citep[eq.~5.17(9)]{Varshalovich1988} reads
\begin{equation}
  P_l(\cos\Theta)=\frac{4\pi}{2l+1}\sum_{m=-l}^{l}Y_l^{m}(\theta,\varphi)^*\,Y_l^m(\theta^\star,\varphi^\star). \label{eq:addition-full}
\end{equation}

\paragraph{Folding to $m\ge0$.} For $m>0$, the symmetry $Y_l^{-m}=(-1)^mY_l^{m*}$ \citep[eq.~5.1(11)]{Varshalovich1988} gives
\[
  Y_l^{-m}(\theta,\varphi)^*\,Y_l^{-m}(\theta^\star,\varphi^\star)=(-1)^{2m}\,Y_l^m(\theta,\varphi)\,Y_l^{m}(\theta^\star,\varphi^\star)^*=\Big[Y_l^m(\theta,\varphi)^*Y_l^m(\theta^\star,\varphi^\star)\Big]^*.
\]
The terms $\pm m$ are complex conjugates, so their sum is twice the real part. The $m=0$ term is real. Hence
\begin{equation}
  P_l(\cos\Theta)=\frac{4\pi}{2l+1}\sum_{m=0}^{l}(2-\delta_{0,m})\Re\Big[Y_l^m(\theta,\varphi)^*\,Y_l^m(\theta^\star,\varphi^\star)\Big]. \label{eq:addition-folded}
\end{equation}

\paragraph{Real form.} Insert \eqref{eq:Ylm-main} in the bracket of \eqref{eq:addition-folded}. Since $N_{lm}$ and $P_l^m$ are real, and since $\Re[e^{-\I m\varphi}e^{\I m\varphi^\star}]$ equals $\cos m(\varphi-\varphi^\star)$,
\[
  \Re\big[Y_l^{m*}(\theta,\varphi)Y_l^m(\theta^\star,\varphi^\star)\big]=N_{lm}^2P_l^m(\cos\theta)P_l^m(\cos\theta^\star)\cos m(\varphi-\varphi^\star).
\]
Since $\frac{4\pi}{2l+1}N_{lm}^2=\frac{(l-m)!}{(l+m)!}$, this gives
\begin{equation}
  P_l(\cos\Theta)=\sum_{m=0}^{l}\frac{(l-m)!}{(l+m)!}(2-\delta_{0,m})P_l^m(\cos\theta)P_l^m(\cos\theta^\star)\cos m(\varphi-\varphi^\star). \label{eq:addition-real}
\end{equation}

\begin{quote}\small
\textbf{Remark (phase convention).} Equations \eqref{eq:addition-folded} and \eqref{eq:addition-real} contain only products of two $Y$ or two $P$. A factor $(-1)^m$ attached to each $P_l^m$ therefore cancels: this chapter does not depend on the Condon--Shortley convention. The convention matters only when a \emph{single} harmonic is transformed (Chapter~\ref{ch:kaula}) and in the recursion relations (Chapter~\ref{ch:gradient}).
\end{quote}

\section{The tidal-raising potential}
Insert the addition theorem \eqref{eq:addition-folded} in the Legendre series \eqref{eq:U-legendre}. The prefactor of the degree-$l$ term is
\[
  \frac{\G M^\star}{r^\star}\Big(\frac{r}{r^\star}\Big)^l=\G M^\star\,r^l\times\frac{1}{(r^\star)^{l+1}}.
\]
We split it this way on purpose: $r^l$ depends only on the field point, and $1/(r^\star)^{l+1}$ only on the star. Grouping $1/(r^\star)^{l+1}$ with the harmonic of the star, $Y_l^m(\theta^\star,\varphi^\star)$, gives the combination $Y_l^m(\theta^\star,\varphi^\star)/(r^\star)^{l+1}$. This is a solid harmonic of the star's position, and it is exactly the quantity that Kaula's transformation converts into orbital elements (Chapter~\ref{ch:kaula}, eq.~\eqref{eq:kaula-Ylm}). Hence
\begin{equation}
  U=-\G M^\star\sum_{l\ge0}\frac{4\pi}{2l+1}\,r^l\sum_{m=0}^{l}(2-\delta_{0,m})\Re\bigg[Y_l^m(\theta,\varphi)^*\,\frac{Y_l^m(\theta^\star,\varphi^\star)}{(r^\star)^{l+1}}\bigg]. \label{eq:U-harmonics}
\end{equation}

\paragraph{The terms $l=0$ and $l=1$.}
\begin{itemize}
  \item The $l=0$ term, $-\G M^\star/r^\star$, does not depend on $\vr$. It exerts no force on the planet's matter.
  \item The $l=1$ term, $-\G M^\star(\vr\cdot\hat{\vr}^\star)/r^{\star2}$, gives a uniform acceleration $\G M^\star\hat{\vr}^\star/r^{\star2}$. This equals the orbital acceleration of the planet's centre. In the planet-centred, non-inertial frame it is cancelled by the transport acceleration, so it does not deform the planet.
\end{itemize}
The tidal-raising potential is therefore $U^{\rm tid}=\sum_{l\ge2}U_l$. From now on all sums over $l$ start at $l=2$. The degree $l$ is kept general and set to $l=2$ where stated.

\chapter{Kaula's transformation}\label{ch:kaula}

Kaula's transformation expresses a solid spherical harmonic of the position of an orbiting body in terms of its Keplerian elements. We first state Kaula's real form, then derive its complex and spherical-harmonic forms. These are the forms used throughout this document.

\section{Elements, angle and frequency}
The star's position $\vr^\star$ relative to the planet is described by its osculating elements referred to the planet's equator, $a^\star,e^\star,i^\star,\omega^\star,\Omega^\star,\M^\star$. In this chapter and the next, the $\star$ is kept on the elements; it becomes necessary in Chapter~\ref{ch:second}.

\paragraph{Kaula's angle.} For each term \citep[eq.~3.71]{Kaula2000},
\begin{equation}
  \nu^\star_{lmpq}=(l-2p)\,\omega^\star+(l-2p+q)\,\M^\star+m\,(\Omega^\star-\thetap). \label{eq:nu}
\end{equation}
where:
\begin{itemize}
  \item $l$ is the degree, $m$ the order ($0\le m\le l$), and $p\in[0,l]$, $q\in\mathbb Z$ the indices of the inclination and eccentricity expansions;
  \item $\omega^\star$ is the argument of pericentre of the star's orbit, measured in the orbital plane from the ascending node on the planet's equator;
  \item $\M^\star$ is the mean anomaly;
  \item $\Omega^\star$ is the longitude of the ascending node of the orbit on the equator;
  \item $\thetap$ is the rotation angle of the planet, so that $\Omega^\star-\thetap$ is the body-fixed longitude of the node.
\end{itemize}
Since $\varphi$ is body-fixed, the rotation of the planet enters through $-m\thetap$. Only $\Omega^\star-\thetap$ appears, so the reference direction of the longitudes is irrelevant.

\paragraph{Signed tidal frequency.} We use the angle itself for the phase and define the frequency as its time derivative:
\begin{equation}
  \sigma^\star_{lmpq}\equiv\dot\nu^\star_{lmpq}=(l-2p)\,\dot\omega^\star+(l-2p+q)\,n^\star+m\,(\dot\Omega^\star-\Op),\qquad \chi^\star_{lmpq}=|\sigma^\star_{lmpq}|. \label{eq:sigma}
\end{equation}
where:
\begin{itemize}
  \item $n^\star\simeq\dot\M^\star$ is the mean motion, $n^\star=\sqrt{\G(M^\star+M_{\mathrm p})/a^{\star3}}$;
  \item $\dot\omega^\star$ and $\dot\Omega^\star$ are the precession rates of the pericentre and of the node, measured relative to the equator;
  \item $\Op=\dot\theta_{\mathrm p}$ is the spin rate of the planet;
  \item $\sigma^\star_{lmpq}$ is signed; $\chi^\star_{lmpq}$ is the physical frequency at which the Love number is evaluated.
\end{itemize}
Equivalent forms of this frequency (with the longitude of pericentre $\varpi=\Omega+\omega$, or the mean longitude $\lambda=\Omega+\omega+\M$), the coplanar form, and the form used in an $N$-body code are compared in Appendix~\ref{app:frequency}. Chapter~\ref{ch:nbody} explains how the rates are evaluated in an $N$-body code.

\paragraph{Parity phase.}
\begin{equation}
  f_{lm}=\frac{(-1)^{l-m}-1}{4}\pi=\begin{cases}0,& l-m\ \text{even},\\-\pi/2,& l-m\ \text{odd},\end{cases}\qquad e^{\I f_{lm}}=\begin{cases}1,\\-\I.\end{cases} \label{eq:flm}
\end{equation}

\section{Kaula's real transformation}
Kaula derives the transformation in three stages:
\begin{enumerate}
  \item eq.~(3.53)--(3.61) convert the spherical coordinates into $(r,f,\omega,\Omega,i)$, which introduces the inclination functions $F_{lmp}(i)$ (eq.~3.62);
  \item p.~35 and eq.~(3.67)--(3.69) replace $(r,f)$ by $(a,e,\M)$, which introduces the eccentricity functions $G_{lpq}(e)$;
  \item the result is eq.~(3.70)--(3.71) \citep{Kaula2000,Kaula1964}.
\end{enumerate}
For any constants $C$ and $S$,
\begin{multline}
  \frac{\PK_{lm}(\cos\theta^\star)}{(r^\star)^{l+1}}\big[C\cos m\varphi^\star+S\sin m\varphi^\star\big]\\
  =\frac{1}{(a^\star)^{l+1}}\sum_{p=0}^{l}\sum_{q\in\mathbb Z}F_{lmp}(i^\star)G_{lpq}(e^\star)\times
  \begin{cases}C\cos\nu+S\sin\nu,& l-m\ \text{even},\\ C\sin\nu-S\cos\nu,& l-m\ \text{odd},\end{cases} \label{eq:kaula-real-CS}
\end{multline}
with $\nu=\nu^\star_{lmpq}$. Kaula's latitude $\phi$ satisfies $\sin\phi=\cos\theta^\star$.

\begin{quote}\small
\textbf{Important.} Equation \eqref{eq:kaula-real-CS} holds for Kaula's function $\PK_{lm}$, which has no Condon--Shortley phase. With $P_l^m=(-1)^m\PK_{lm}$, the left-hand side would change sign for odd $m$.
\end{quote}

Choosing $C=\cos m\varphi$ and $S=\sin m\varphi$, where $\varphi$ is an arbitrary longitude (later the field point), gives:
\begin{itemize}
  \item on the left, $C\cos m\varphi^\star+S\sin m\varphi^\star=\cos m(\varphi-\varphi^\star)$;
  \item on the right, $C\cos\nu+S\sin\nu=\cos(\nu-m\varphi)$ and $C\sin\nu-S\cos\nu=\sin(\nu-m\varphi)$.
\end{itemize}
Hence
\begin{equation}
  \frac{\PK_{lm}(\cos\theta^\star)}{(r^\star)^{l+1}}\cos m(\varphi-\varphi^\star)=\frac{1}{(a^\star)^{l+1}}\sum_{p,q}F_{lmp}(i^\star)G_{lpq}(e^\star)\Big\{{\cos\atop\sin}\Big\}^{l-m\ \text{even}}_{l-m\ \text{odd}}\big(\nu^\star_{lmpq}-m\varphi\big). \label{eq:kaula-real}
\end{equation}

\section{Complex form}
The real form \eqref{eq:kaula-real} holds for any longitude $\varphi$. Evaluating it at two particular values of $\varphi$ gives the real and imaginary parts of $\PK_{lm}e^{\I m\varphi^\star}$, which together form the complex version. To shorten the notation, write
\[
  \PK\equiv\PK_{lm}(\cos\theta^\star),\qquad \Sigma[\,X\,]\equiv\frac{1}{(a^\star)^{l+1}}\sum_{p=0}^{l}\sum_{q\in\mathbb Z}F_{lmp}(i^\star)G_{lpq}(e^\star)\,X,
\]
so that \eqref{eq:kaula-real} reads $\PK\cos m(\varphi-\varphi^\star)/(r^\star)^{l+1}=\Sigma[\cos(\nu-m\varphi)]$ for $l-m$ even, and $\Sigma[\sin(\nu-m\varphi)]$ for $l-m$ odd.

\paragraph{Two evaluations ($m\ge1$).}
\begin{itemize}
  \item At $\varphi=0$, the left-hand side is $\PK\cos m\varphi^\star/(r^\star)^{l+1}$ and the phase on the right is $\nu$.
  \item At $\varphi=\pi/(2m)$, the left-hand side is $\cos(\tfrac\pi2-m\varphi^\star)=\sin m\varphi^\star$ times $\PK/(r^\star)^{l+1}$, and the phase on the right is $\nu-\tfrac\pi2$.
\end{itemize}

\paragraph{Case $l-m$ even.} With $\cos(\nu-\tfrac\pi2)=\sin\nu$, the two evaluations give
\begin{align*}
  \frac{\PK\cos m\varphi^\star}{(r^\star)^{l+1}}&=\Sigma[\cos\nu],\\
  \frac{\PK\sin m\varphi^\star}{(r^\star)^{l+1}}&=\Sigma[\sin\nu].
\end{align*}
Multiplying the second line by $\I$ and adding it to the first,
\[
  \frac{\PK\,e^{\I m\varphi^\star}}{(r^\star)^{l+1}}=\Sigma\big[\cos\nu+\I\sin\nu\big]=\Sigma\big[e^{\I\nu}\big].
\]

\paragraph{Case $l-m$ odd.} With $\sin(\nu-\tfrac\pi2)=-\cos\nu$, the two evaluations give
\begin{align*}
  \frac{\PK\cos m\varphi^\star}{(r^\star)^{l+1}}&=\Sigma[\sin\nu],\\
  \frac{\PK\sin m\varphi^\star}{(r^\star)^{l+1}}&=-\Sigma[\cos\nu].
\end{align*}
Multiplying the second line by $\I$ and adding it to the first,
\[
  \frac{\PK\,e^{\I m\varphi^\star}}{(r^\star)^{l+1}}=\Sigma\big[\sin\nu-\I\cos\nu\big]=\Sigma\big[-\I\,e^{\I\nu}\big].
\]

\paragraph{Case $m=0$.} Here $e^{\I m\varphi^\star}=1$, and only the evaluation at $\varphi=0$ is available. The left-hand side is real. The complex form then requires the imaginary part of the right-hand side to vanish:
\begin{itemize}
  \item $\Sigma[\sin\nu]=0$ for $l$ even;
  \item $\Sigma[\cos\nu]=0$ for $l$ odd.
\end{itemize}
Both hold because the terms $(p,q)$ and $(l-p,-q)$ cancel in pairs. Indeed $F_{l0p}=(-1)^lF_{l0,l-p}$, $G_{lpq}=G_{l,l-p,-q}$ and $\nu_{l0pq}=-\nu_{l0,l-p,-q}$.

\paragraph{Both parities together.} By \eqref{eq:flm}, $e^{\I f_{lm}}=1$ for $l-m$ even and $-\I$ for $l-m$ odd. The two cases are therefore the single formula
\begin{equation}
  \boxed{\frac{\PK_{lm}(\cos\theta^\star)\,e^{\I m\varphi^\star}}{(r^\star)^{l+1}}=\frac{1}{(a^\star)^{l+1}}\sum_{p=0}^{l}\sum_{q\in\mathbb Z}F_{lmp}(i^\star)G_{lpq}(e^\star)\,e^{\I\zeta^\star_{lmpq}},\qquad \zeta^\star_{lmpq}=\nu^\star_{lmpq}+f_{lm}.} \label{eq:kaula-complex}
\end{equation}
where:
\begin{itemize}
  \item $(r^\star,\theta^\star,\varphi^\star)$ are the distance, colatitude and body-fixed longitude of the star;
  \item $(a^\star,e^\star,i^\star)$ are its semi-major axis, eccentricity and inclination on the planet's equator;
  \item $F_{lmp}$ and $G_{lpq}$ are Kaula's inclination and eccentricity functions (Appendix~\ref{app:kaula});
  \item $\nu^\star_{lmpq}$ is Kaula's angle~\eqref{eq:nu} and $f_{lm}$ the parity phase~\eqref{eq:flm}.
\end{itemize}

\section{Spherical-harmonic form}
With $Y_l^m=N_{lm}P_l^me^{\I m\varphi}=(-1)^mN_{lm}\PK_{lm}e^{\I m\varphi}$,
\begin{equation}
  \boxed{\frac{Y_l^m(\theta^\star,\varphi^\star)}{(r^\star)^{l+1}}=\frac{(-1)^m}{(a^\star)^{l+1}}\sqrt{\frac{2l+1}{4\pi}\frac{(l-m)!}{(l+m)!}}\sum_{p=0}^{l}\sum_{q\in\mathbb Z}F_{lmp}(i^\star)G_{lpq}(e^\star)\,e^{\I\zeta^\star_{lmpq}}\qquad(0\le m\le l).} \label{eq:kaula-Ylm}
\end{equation}
The conjugate has $e^{-\I\zeta}$. The factor $(-1)^m$ is essential: without it every odd-$m$ term changes sign.

\paragraph{Consistency check for any $l$.} We reconstruct the real form \eqref{eq:kaula-real} from \eqref{eq:kaula-Ylm} and the conjugate harmonic of the field point, $Y_l^{m*}(\theta,\varphi)=(-1)^mN_{lm}\PK_{lm}(\cos\theta)e^{-\I m\varphi}$. The two factors $(-1)^m$ cancel, and
\[
  \Re\big[e^{-\I m\varphi}e^{\I(\nu+f_{lm})}\big]=\Re\big[e^{\I f_{lm}}e^{\I(\nu-m\varphi)}\big]=\begin{cases}\cos(\nu-m\varphi),&l-m\ \text{even},\\ \Re[-\I e^{\I(\nu-m\varphi)}]=\sin(\nu-m\varphi),&l-m\ \text{odd},\end{cases}
\]
which is \eqref{eq:kaula-real}.

\section{The tidal-raising potential in orbital elements}
Inserting \eqref{eq:kaula-Ylm} into \eqref{eq:U-harmonics} and using $\frac{4\pi}{2l+1}N_{lm}=\sqrt{\frac{4\pi}{2l+1}\frac{(l-m)!}{(l+m)!}}$ gives $U^{\rm tid}=\sum_{l\ge2}\sum_{m,p,q}U_{lmpq}$, with
\begin{align}
  U_{lmpq}(\vr)&=-\G M^\star\frac{r^l}{(a^\star)^{l+1}}(-1)^m\sqrt{\frac{4\pi}{2l+1}\frac{(l-m)!}{(l+m)!}}(2-\delta_{0,m})\notag\\
  &\qquad\times F_{lmp}(i^\star)G_{lpq}(e^\star)\,\Re\Big[Y_l^m(\theta,\varphi)^*e^{\I\zeta^\star_{lmpq}}\Big] \label{eq:Ulmpq-complex}\\
  &=-\G M^\star\frac{r^l}{(a^\star)^{l+1}}\frac{(l-m)!}{(l+m)!}(2-\delta_{0,m})\PK_{lm}(\cos\theta)\notag\\
  &\qquad\times F_{lmp}(i^\star)G_{lpq}(e^\star)\Big\{{\cos\atop\sin}\Big\}^{l-m\ \text{even}}_{l-m\ \text{odd}}\big(\nu^\star_{lmpq}-m\varphi\big). \label{eq:Ulmpq-real}
\end{align}
where $\vr=(r,\theta,\varphi)$ is the field point in the planet (body-fixed longitude $\varphi$), $M^\star$ the mass of the star, and the elements $(a^\star,e^\star,i^\star)$ and the angle $\nu^\star_{lmpq}$ are those of the star's orbit, as in~\eqref{eq:kaula-complex}.

\paragraph{Forcing frequency of each term.} This property is needed in Chapter~\ref{ch:tidal} to apply the Love number. Fix a point of the planet: $r$, $\theta$ and $\varphi$ are constant, since $\varphi$ rotates with the planet. In \eqref{eq:Ulmpq-real}, the only time-dependent part is then the phase $\nu^\star_{lmpq}(t)-m\varphi$, whose rate is
\[
  \frac{\dd}{\dd t}\big(\nu^\star_{lmpq}-m\varphi\big)=\dot\nu^\star_{lmpq}=\sigma^\star_{lmpq}\qquad\text{by \eqref{eq:sigma}}.
\]
Each term $U_{lmpq}$ is therefore a periodic forcing of frequency $\sigma^\star_{lmpq}$ for the material of the planet. This is the frequency at which the planet responds to that term, and the argument of the Love number in Chapter~\ref{ch:tidal}.

The series in $q$ converges for any $e<1$, more slowly as $e$ grows; its truncation is discussed in Chapter~\ref{ch:nbody}.

\chapter{The tidal potential}\label{ch:tidal}

\section{Response of one term}
At a body-fixed point of the surface, $(R,\theta,\varphi)$, each term \eqref{eq:Ulmpq-complex} can be written
\[
  U_{lmpq}(R,\theta,\varphi;t)=\Re\big[\mathcal U_{lmpq}(R,\theta,\varphi)\,e^{\I\nu^\star_{lmpq}(t)}\big].
\]
Here $\mathcal U_{lmpq}$ is a complex amplitude that depends on time only through the slow variation of $a^\star,e^\star,i^\star$ (hypothesis H3). Its phase advances as $\nu^\star_{lmpq}(t)$, at the rate $\sigma^\star_{lmpq}$, and because $\varphi$ is body-fixed this is the frequency seen by the material.

By hypothesis H1, the response to $\Re[\mathcal Ue^{\I\sigma t}]$ is $\Re[\kb_l(\sigma)\,\mathcal Ue^{\I\sigma t}]$, where $\kb_l$ is the Fourier transform of the (real) response kernel. Hence \eqref{eq:love-def}: $\kb_l(-\sigma)=\kb_l(\sigma)^*$. In terms of the physical frequency $\chi=|\sigma|$,
\begin{equation}
  \kb_l(\sigma^\star_{lmpq})=\begin{cases}\kb_l(\chi^\star_{lmpq}),&\sigma^\star_{lmpq}>0,\\ \kb_l(\chi^\star_{lmpq})^*,&\sigma^\star_{lmpq}<0,\end{cases}\qquad\Longleftrightarrow\qquad \epsilon_l(\sigma)=\sgn(\sigma)\,\epsilon_l(\chi). \label{eq:love-chi}
\end{equation}
Here $\sigma^\star_{lmpq}$ is the signed frequency~\eqref{eq:sigma}, $\chi^\star_{lmpq}=|\sigma^\star_{lmpq}|$ the physical frequency, and $\kb_l=|\kb_l|e^{-\I\epsilon_l}$ (Appendix~\ref{app:conventions}).

\paragraph{Sign of the lag.} With $\epsilon_l(\chi)>0$, the response lags behind the forcing for both signs of $\sigma$. For instance, a constant time lag $\Delta t>0$ gives $\epsilon_l=\sigma\Delta t$, whose sign follows $\sigma$, and the time-averaged work done on the bulge is then dissipative. Using $\epsilon_l(\chi)$ for $\sigma<0$ without conjugation would turn the lag into an advance.

\paragraph{Exterior potential.} The degree-$l$ potential of the deformed planet:
\begin{itemize}
  \item satisfies Laplace's equation outside the planet;
  \item vanishes at infinity;
  \item is continuous at $r=R$.
\end{itemize}
It is therefore proportional to $r^{-(l+1)}Y_l^m$, with the factor $(R/r)^{l+1}$.

\section{Result}
The tidal potential of the term $(l,m,p,q)$ is
\begin{equation}
  T_{lmpq}(\vr)=\Big(\frac{R}{r}\Big)^{l+1}\Re\Big[\kb_l(\sigma^\star_{lmpq})\,\mathcal U_{lmpq}(R,\theta,\varphi)\,e^{\I\nu^\star_{lmpq}}\Big],\qquad r\ge R. \label{eq:T-def}
\end{equation}
The lag enters only once, through the argument of $\kb_l$. With $r^l\to R^l$ and $R^l(R/r)^{l+1}=R^{2l+1}/r^{l+1}$,
\begin{equation}
  \boxed{\begin{aligned}T_{lmpq}(\vr)=-|\kb_l|\,\G M^\star\frac{R^{2l+1}}{(a^\star)^{l+1}r^{l+1}}&(-1)^m\sqrt{\frac{4\pi}{2l+1}\frac{(l-m)!}{(l+m)!}}(2-\delta_{0,m})\\
  &\times F_{lmp}(i^\star)G_{lpq}(e^\star)\,\Re\Big[Y_l^m(\theta,\varphi)^*e^{\I(\zeta^\star_{lmpq}-\epsilon_l)}\Big],\end{aligned}} \label{eq:T-complex}
\end{equation}
where:
\begin{itemize}
  \item $\vr=(r,\theta,\varphi)$ is a point outside the planet, $r\ge R$, with $R$ the radius of the planet and $\varphi$ the body-fixed longitude;
  \item $|\kb_l|\equiv|\kb_l(\chi^\star_{lmpq})|$ and $\epsilon_l\equiv\epsilon_l(\sigma^\star_{lmpq})$ are the modulus and the phase lag of the Love number, evaluated at the frequency~\eqref{eq:sigma} of the term;
  \item $\zeta^\star_{lmpq}=\nu^\star_{lmpq}+f_{lm}$ (Kaula's angle plus parity phase, Chapter~\ref{ch:kaula});
  \item $a^\star,e^\star,i^\star$ are the elements of the star's orbit on the planet's equator;
  \item the total tidal potential is $T=\sum_{l\ge2}\sum_{m,p,q}T_{lmpq}$.
\end{itemize}

\paragraph{Real form.}
\begin{multline}
  T_{lmpq}(\vr)=-|\kb_l|\G M^\star\frac{R^{2l+1}}{(a^\star)^{l+1}r^{l+1}}\frac{(l-m)!}{(l+m)!}(2-\delta_{0,m})\PK_{lm}(\cos\theta)\\
  \times F_{lmp}G_{lpq}\Big\{{\cos\atop\sin}\Big\}^{l-m\ \text{even}}_{l-m\ \text{odd}}\big(\nu^\star_{lmpq}-m\varphi-\epsilon_l\big). \label{eq:T-real}
\end{multline}

\paragraph{Love-number split.} From $|\kb|e^{-\I\epsilon}=\Re\kb+\I\Im\kb$ and $\Re[(x+\I y)z]=x\Re z-y\Im z$,
\begin{equation}
  |\kb_l|\Re\big[Y_l^{m*}e^{\I(\zeta-\epsilon_l)}\big]=\Re\kb_l\,\Re\big[Y_l^{m*}e^{\I\zeta}\big]-\Im\kb_l\,\Im\big[Y_l^{m*}e^{\I\zeta}\big],\qquad \kb_l=\kb_l(\sigma^\star_{lmpq}). \label{eq:love-split}
\end{equation}

\chapter{Gradient of the tidal potential and the force}\label{ch:gradient}

\section{Components of a vector}\label{sec:components}
Several bases are used for the gradient and the force. For any vector $\mathbf V$, we write its components as follows.
\begin{itemize}
  \item \textbf{Local polar basis} at the point $(r,\theta,\varphi)$: $\mathbf V=V_r\er+V_\theta\et+V_\varphi\ep$, where $\er$, $\et$, $\ep$ point outward, southward and eastward (Appendix~\ref{app:conventions}).
  \item \textbf{Body-fixed Cartesian basis}: $\mathbf V=V_x\mathbf e_x+V_y\mathbf e_y+V_z\mathbf e_z$, with $\mathbf e_z$ along the spin axis and $\mathbf e_x$ towards $\varphi=0$. The two horizontal components are combined into one complex number:
  \begin{equation}
    V_+\equiv V_x+\I V_y. \label{eq:Vplus}
  \end{equation}
\end{itemize}
With $\er=\sin\theta\cos\varphi\,\mathbf e_x+\sin\theta\sin\varphi\,\mathbf e_y+\cos\theta\,\mathbf e_z$, $\et=\cos\theta\cos\varphi\,\mathbf e_x+\cos\theta\sin\varphi\,\mathbf e_y-\sin\theta\,\mathbf e_z$ and $\ep=-\sin\varphi\,\mathbf e_x+\cos\varphi\,\mathbf e_y$, the two sets of components are related by
\begin{equation}
  V_r=\sin\theta\,\Re\big[V_+e^{-\I\varphi}\big]+\cos\theta\,V_z,\qquad
  V_\theta=\cos\theta\,\Re\big[V_+e^{-\I\varphi}\big]-\sin\theta\,V_z,\qquad
  V_\varphi=\Im\big[V_+e^{-\I\varphi}\big]. \label{eq:polar-from-cart}
\end{equation}
Indeed $V_+e^{-\I\varphi}=(V_x\cos\varphi+V_y\sin\varphi)+\I(V_y\cos\varphi-V_x\sin\varphi)$: the real part is the horizontal component along the meridian plane, the imaginary part the component along $\ep$.

When the gradient of $T$ is written component by component, $(\nabla T)_r$, $(\nabla T)_\theta$, $(\nabla T)_\varphi$, $(\nabla T)_z$ and $(\nabla T)_+$ denote these components of $\mathbf V=\nabla T$.

\section{What has to be differentiated}
Start from the complex form of the tidal potential, eq.~\eqref{eq:T-complex}. Separate the factor that depends on the field point $\vr$ from the factor that does not. Each term then reads
\begin{equation}
  T_{lmpq}(\vr)=\Re\big[W_{lmpq}\,g_{lm}(\vr)\big],\qquad g_{lm}(\vr)\equiv\frac{Y_l^m(\theta,\varphi)^*}{r^{l+1}},\qquad W_{lmpq}\equiv\mathcal A_{lmpq}\,e^{\I(\zeta^\star_{lmpq}-\epsilon_l)}, \label{eq:T-W-g}
\end{equation}
with the real coefficient
\begin{equation}
  \mathcal A_{lmpq}=-|\kb_l|\G M^\star\frac{R^{2l+1}}{(a^\star)^{l+1}}(-1)^m\sqrt{\frac{4\pi}{2l+1}\frac{(l-m)!}{(l+m)!}}(2-\delta_{0,m})F_{lmp}(i^\star)G_{lpq}(e^\star). \label{eq:Almpq}
\end{equation}
Only $g_{lm}$ depends on the field point; $W_{lmpq}$ depends on the star and on the response. Since $\Re$ is linear and $\nabla$ is a real operator,
\begin{equation}
  \nabla T_{lmpq}=\Re\big[W_{lmpq}\,\nabla g_{lm}\big]. \label{eq:gradT-g}
\end{equation}
The whole problem is therefore to compute $\nabla g_{lm}$. This is done in two ways below: in the local polar basis (Section~\ref{sec:polar}), and in the body-fixed Cartesian basis (Section~\ref{sec:cart}). The function $g_{lm}$ of \eqref{eq:T-W-g} is an exterior solid harmonic: it is harmonic for $r>0$, homogeneous of degree $-(l+1)$ in $\vr$, and proportional to $e^{-\I m\varphi}$.

\paragraph{A compact real form.} It is convenient to introduce a second real coefficient and the phase at the field point:
\begin{equation}
  \mathcal K_{lmpq}\equiv\G M^\star\frac{R^{2l+1}}{(a^\star)^{l+1}}\frac{(l-m)!}{(l+m)!}(2-\delta_{0,m})F_{lmp}(i^\star)G_{lpq}(e^\star),\qquad
  \Lambda_{lmpq}\equiv\zeta^\star_{lmpq}-m\varphi. \label{eq:Klmpq}
\end{equation}
They are related to $\mathcal A_{lmpq}$ by $\mathcal A_{lmpq}(-1)^mN_{lm}=-|\kb_l|\mathcal K_{lmpq}$. Since $\cos(x+f_{lm})$ equals $\cos x$ for $l-m$ even and $\sin x$ for $l-m$ odd, the real form \eqref{eq:T-real} becomes, for both parities,
\begin{equation}
  T_{lmpq}(\vr)=-|\kb_l|\frac{\mathcal K_{lmpq}}{r^{l+1}}\PK_{lm}(\cos\theta)\cos\big(\Lambda_{lmpq}-\epsilon_l\big). \label{eq:T-compact}
\end{equation}

\section{Gradient in the local polar basis}\label{sec:polar}
\paragraph{Real form.} Apply $\nabla=\er\partial_r+\et\,r^{-1}\partial_\theta+\ep\,(r\sin\theta)^{-1}\partial_\varphi$ to \eqref{eq:T-compact}. Three facts are needed:
\begin{itemize}
  \item $\partial_r r^{-(l+1)}=-(l+1)r^{-(l+2)}$;
  \item only $\PK_{lm}(\cos\theta)$ depends on $\theta$;
  \item only $\Lambda_{lmpq}$ depends on $\varphi$, with $\partial_\varphi\Lambda_{lmpq}=-m$, so $\partial_\varphi\cos(\Lambda-\epsilon)=m\sin(\Lambda-\epsilon)$.
\end{itemize}
This gives
\begin{equation}
  \boxed{\begin{aligned}
  (\nabla T_{lmpq})_r&=(l+1)|\kb_l|\frac{\mathcal K_{lmpq}}{r^{l+2}}\PK_{lm}(\cos\theta)\cos(\Lambda_{lmpq}-\epsilon_l),\\
  (\nabla T_{lmpq})_\theta&=-|\kb_l|\frac{\mathcal K_{lmpq}}{r^{l+2}}\frac{\dd\PK_{lm}}{\dd\theta}\cos(\Lambda_{lmpq}-\epsilon_l),\\
  (\nabla T_{lmpq})_\varphi&=-|\kb_l|\frac{\mathcal K_{lmpq}}{r^{l+2}}\frac{m\PK_{lm}}{\sin\theta}\sin(\Lambda_{lmpq}-\epsilon_l).
  \end{aligned}} \label{eq:grad-polar-real}
\end{equation}
where:
\begin{itemize}
  \item $\vr=(r,\theta,\varphi)$ is the field point, $\varphi$ body-fixed;
  \item $\mathcal K_{lmpq}$ is the real coefficient~\eqref{eq:Klmpq}, which contains the masses, the radius $R$, $a$, $F_{lmp}(i)$ and $G_{lpq}(e)$;
  \item $\Lambda_{lmpq}=\zeta_{lmpq}-m\varphi$ is the phase of the term at the field point;
  \item $|\kb_l|$ and $\epsilon_l$ are the modulus and lag of the Love number at the frequency $\sigma_{lmpq}$ of~\eqref{eq:sigma}.
\end{itemize}
The two angular functions are proved in Appendix~\ref{app:harmonics} (eq.~\eqref{eq:dPK}--\eqref{eq:mPK-sin}):
\begin{align}
  \frac{\dd\PK_{lm}}{\dd\theta}&=\tfrac12\big[(l+m)(l-m+1)\PK_{l,m-1}-\PK_{l,m+1}\big]\ \ (m\ge1),\qquad \frac{\dd\PK_{l0}}{\dd\theta}=-\PK_{l1}, \label{eq:dPK-main}\\
  \frac{m\PK_{lm}}{\sin\theta}&=\tfrac12\big[\PK_{l+1,m+1}+(l-m+1)(l-m+2)\PK_{l+1,m-1}\big]\ \ (m\ge1),\qquad 0\ \ (m=0). \label{eq:mPK-main}
\end{align}
The azimuthal factor $m\PK_{lm}/\sin\theta$ contains no $1/\sin\theta$: it is a combination of degree-$(l+1)$ functions and is regular at the poles.

\paragraph{Spherical-harmonic form.} The same gradient can be written with spherical harmonics. \citet[\S7.3, eq.~25--27 and 51]{Varshalovich1988} give the gradient of $f(r)Y_{LM}$:
\[
  \nabla\{f(r)Y_{LM}\}=f'(r)\,\mathbf Y^{(-1)}_{LM}+\sqrt{L(L+1)}\,\frac{f}{r}\,\mathbf Y^{(1)}_{LM},
\]
with the polar components $[\mathbf Y^{(-1)}_{LM}]_r=Y_{LM}$, $[\mathbf Y^{(1)}_{LM}]_\theta=\partial_\theta Y_{LM}/\sqrt{L(L+1)}$ and $[\mathbf Y^{(1)}_{LM}]_\varphi=\I M\,Y_{LM}/(\sqrt{L(L+1)}\sin\theta)$. The $\theta$-derivative is \citep[eq.~5.8(5)]{Varshalovich1988}
\[
  \partial_\theta Y_{LM}=\tfrac12\sqrt{L(L+1)-M(M+1)}\,Y_{L,M+1}e^{-\I\varphi}-\tfrac12\sqrt{L(L+1)-M(M-1)}\,Y_{L,M-1}e^{\I\varphi}.
\]
Apply this to $g_{lm}$ of \eqref{eq:T-W-g}, i.e.\ $f=r^{-(l+1)}$ and the \emph{conjugate} harmonic. Conjugation changes $\partial_\varphi Y_l^m=\I mY_l^m$ into $\partial_\varphi Y_l^{m*}=-\I mY_l^{m*}$, and swaps $e^{\mp\I\varphi}$. Then insert in \eqref{eq:gradT-g}, with $A^\pm_{lm}=\frac12\sqrt{l(l+1)-m(m\pm1)}$ and $\Phi_{lmpq}=\zeta^\star_{lmpq}-\epsilon_l$:
\begin{equation}
  \begin{aligned}
  (\nabla T_{lmpq})_r&=-(l+1)\frac{\mathcal A_{lmpq}}{r^{l+2}}\Re\big[Y_l^{m*}e^{\I\Phi_{lmpq}}\big],\\
  (\nabla T_{lmpq})_\theta&=\frac{\mathcal A_{lmpq}}{r^{l+2}}\Re\Big[\big(A^+_{lm}Y_l^{m+1\,*}e^{\I\varphi}-A^-_{lm}Y_l^{m-1\,*}e^{-\I\varphi}\big)e^{\I\Phi_{lmpq}}\Big]\quad(m\ge1),\\
  (\nabla T_{l0pq})_\theta&=\frac{\mathcal A_{l0pq}}{r^{l+2}}\sqrt{l(l+1)}\,\Re\big[Y_l^{1*}e^{\I\varphi}e^{\I\Phi_{l0pq}}\big],\\
  (\nabla T_{lmpq})_\varphi&=\frac{\mathcal A_{lmpq}}{r^{l+2}}\frac{m}{\sin\theta}\Im\big[Y_l^{m*}e^{\I\Phi_{lmpq}}\big]\\
  &=-\frac{\mathcal A_{lmpq}}{2r^{l+2}}\Im\Big[\big(c^+_{lm}Y_{l+1}^{m+1\,*}e^{\I\varphi}+c^-_{lm}Y_{l+1}^{m-1\,*}e^{-\I\varphi}\big)e^{\I\Phi_{lmpq}}\Big].
  \end{aligned} \label{eq:grad-polar-cplx}
\end{equation}
\begin{itemize}
  \item The $m=0$ meridional component uses $\partial_\theta Y_l^{0*}=\sqrt{l(l+1)}Y_l^{1*}e^{\I\varphi}$ \citep[eq.~5.8(5)]{Varshalovich1988}, so no negative order appears.
  \item The $\varphi$-component: $\Re[-\I m\,z]=m\,\Im z$, hence the $\Im$.
  \item The last line uses \eqref{eq:Y-over-sin} (Appendix~\ref{app:harmonics}), with the constants $c^\pm_{lm}$ of \eqref{eq:c-constants} below.
\end{itemize}

\section{Gradient in the body-fixed Cartesian basis}\label{sec:cart}
\paragraph{Ladder operators.} In the body-fixed Cartesian basis, define $D_\pm\equiv\partial_x\pm\I\partial_y$. Then, by \eqref{eq:Vplus}, $(\nabla T)_+=D_+T$ for any real $T$. The operators $\partial_z$ and $D_\pm$:
\begin{itemize}
  \item commute with the Laplacian, so they map harmonic functions to harmonic functions;
  \item lower the homogeneity degree by one;
  \item multiply the azimuthal dependence by $1$ and $e^{\pm\I\varphi}$ respectively.
\end{itemize}
The function $g_{lm}$ of \eqref{eq:T-W-g} is harmonic, of degree $-(l+1)$ and proportional to $e^{-\I m\varphi}$. Its images are therefore harmonic, of degree $-(l+2)$ and proportional to $e^{-\I m\varphi}$ or $e^{-\I(m\mp1)\varphi}$, i.e.\ each is a single degree-$(l+1)$ solid harmonic $g_{l+1,m'}$:
\begin{equation}
  \partial_zg_{lm}=-c^0_{lm}\,g_{l+1,m},\qquad D_-g_{lm}=c^+_{lm}\,g_{l+1,m+1},\qquad D_+g_{lm}=-c^-_{lm}\,g_{l+1,m-1}\quad(m\ge1), \label{eq:ladder}
\end{equation}
with the constants
\begin{equation}
  c^0_{lm}=\sqrt{\tfrac{2l+1}{2l+3}\big[(l+1)^2-m^2\big]},\qquad c^\pm_{lm}=\sqrt{\tfrac{2l+1}{2l+3}(l\pm m+1)(l\pm m+2)}. \label{eq:c-constants}
\end{equation}
For $m=0$, $g_{l0}$ is real, so $D_+g_{l0}=(D_-g_{l0})^*=c^+_{l0}\,Y_{l+1}^1/r^{l+2}$ and no negative order appears.

\paragraph{Origin of the constants.} They follow from \citet[eq.~7.3(51), first line]{Varshalovich1988} with $f=r^{-(l+1)}$. The first term vanishes because $(\dd/\dd r+(l+1)/r)\,r^{-(l+1)}=0$, which leaves
\[
  \nabla(r^{-(l+1)}Y_{lm})=\sqrt{(l+1)(2l+1)}\,r^{-(l+2)}\mathbf Y^{l+1}_{lm}.
\]
Expanding $\mathbf Y^{l+1}_{lm}$ with Clebsch--Gordan coefficients on the covariant basis $\mathbf e_{\pm1}=\mp(\mathbf e_x\pm\I\mathbf e_y)/\sqrt2$, $\mathbf e_0=\mathbf e_z$ \citep[eq.~1.1(20)]{Varshalovich1988}, and taking the complex conjugate, gives \eqref{eq:ladder}--\eqref{eq:c-constants}.

\paragraph{Real Cartesian components.} Insert \eqref{eq:ladder} in \eqref{eq:gradT-g}, express $g_{l+1,m'}$ with $\PK_{l+1,m'}$, and use $\Re[Wz]$ for $\partial_z$. For $D_+$, use $D_+\Re[Wg]=\tfrac12\big[W\,D_+g+W^*(D_-g)^*\big]$. With $X\equiv\Lambda_{lmpq}-\epsilon_l$:
\begin{equation}
  \boxed{\begin{aligned}
  (\nabla T_{lmpq})_z&=(l-m+1)|\kb_l|\frac{\mathcal K_{lmpq}}{r^{l+2}}\PK_{l+1,m}(\cos\theta)\cos X,\\
  (\nabla T_{lmpq})_+&=|\kb_l|\frac{\mathcal K_{lmpq}}{2r^{l+2}}e^{\I\varphi}\big[\PK_{l+1,m+1}\,e^{-\I X}-(l-m+1)(l-m+2)\PK_{l+1,m-1}\,e^{\I X}\big]\quad(m\ge1),\\
  (\nabla T_{l0pq})_+&=|\kb_l|\frac{\mathcal K_{l0pq}}{r^{l+2}}e^{\I\varphi}\PK_{l+1,1}\cos X.
  \end{aligned}} \label{eq:grad-cart}
\end{equation}
where $V_z$ and $V_+=V_x+\I V_y$ are the body-fixed Cartesian components of Section~\ref{sec:components}, $\vr=(r,\theta,\varphi)$ the field point, $\mathcal K_{lmpq}$ the coefficient~\eqref{eq:Klmpq}, $X=\Lambda_{lmpq}-\epsilon_l$, and $\PK_{l+1,m'}=\PK_{l+1,m'}(\cos\theta)$.

All components involve only degree-$(l+1)$ solid harmonics, which is the input of the second Kaula transformation (Chapter~\ref{ch:second}). Projected with \eqref{eq:polar-from-cart}, \eqref{eq:grad-cart} gives back \eqref{eq:grad-polar-real}.

\section{The force}
By the force convention \eqref{eq:force-convention}, $\Fsp=+M^\star\nabla T$ evaluated at the star, $\vr=\vr^\star=(r^\star,\theta^\star,\varphi^\star)$. Evaluate \eqref{eq:grad-polar-real} there, i.e.\ with $\Lambda^\star_{lmpq}\equiv\zeta^\star_{lmpq}-m\varphi^\star$. Then split $|\kb_l|\cos(\Lambda-\epsilon_l)$ and $|\kb_l|\sin(\Lambda-\epsilon_l)$ with \eqref{eq:love-split}:
\begin{equation}
  \begin{aligned}
  F^{\star\to\mathrm p}_r&=M^\star\sum_{lmpq}(l+1)\frac{\mathcal K_{lmpq}}{(r^\star)^{l+2}}\PK_{lm}(\cos\theta^\star)\big[\Re\kb_l\cos\Lambda^\star_{lmpq}-\Im\kb_l\sin\Lambda^\star_{lmpq}\big],\\
  F^{\star\to\mathrm p}_\theta&=-M^\star\sum_{lmpq}\frac{\mathcal K_{lmpq}}{(r^\star)^{l+2}}\frac{\dd\PK_{lm}}{\dd\theta}(\cos\theta^\star)\big[\Re\kb_l\cos\Lambda^\star_{lmpq}-\Im\kb_l\sin\Lambda^\star_{lmpq}\big],\\
  F^{\star\to\mathrm p}_\varphi&=-M^\star\sum_{lmpq}\frac{\mathcal K_{lmpq}}{(r^\star)^{l+2}}\frac{m\PK_{lm}}{\sin\theta}(\cos\theta^\star)\big[\Re\kb_l\sin\Lambda^\star_{lmpq}+\Im\kb_l\cos\Lambda^\star_{lmpq}\big],
  \end{aligned} \label{eq:force-polar}
\end{equation}
where:
\begin{itemize}
  \item $(r^\star,\theta^\star,\varphi^\star)$ is the position of the star seen from the planet, $\varphi^\star$ body-fixed;
  \item $\Lambda^\star_{lmpq}=\zeta^\star_{lmpq}-m\varphi^\star$ is the phase of the term at the star;
  \item $\kb_l=\kb_l(\sigma^\star_{lmpq})$ is the complex Love number at the frequency of the term, split into its real (elastic) and imaginary (dissipative) parts;
  \item $\Fps=-\Fsp$ is the force on the star.
\end{itemize}

\section{Physical check: circular equatorial orbit}\label{sec:physical-check}
We check the signs of \eqref{eq:force-polar} on the simplest case.
\paragraph{Setting.} The orbit is circular ($e=0$) and lies in the equator ($i=0$); the tidal lag is a constant time lag $\Delta t>0$; the degree is $l=2$.

\paragraph{Surviving terms.}
\begin{itemize}
  \item For $e=0$: $G_{2pq}(0)=1$ for $q=0$ and $0$ otherwise, so only $q=0$ remains.
  \item For $i=0$: $F_{2mp}(0)\neq0$ only for $(m,p)=(0,1)$, with $F_{201}(0)=-\tfrac12$, and for $(m,p)=(2,0)$, with $F_{220}(0)=3$ (Table~\ref{tab:F2mp}).
\end{itemize}
Two terms remain: $(m,p,q)=(0,1,0)$ and $(2,0,0)$.

\paragraph{Their frequencies,} from \eqref{eq:sigma} with $\dot\omega=\dot\Omega=0$ and $\dot\M=n$:
\[
  \sigma_{2010}=(2-2)\,n=0,\qquad \sigma_{2200}=2n-2\Op=-2(\Op-n).
\]
With a constant time lag, $\epsilon_2(\sigma)=\sigma\Delta t$, so
\[
  \kb_2(\sigma)=k_2\,e^{-\I\sigma\Delta t},\qquad \kb_2(\sigma_{2010})=k_2,
\]
\[
  \Im\kb_2(\sigma_{2200})=-k_2\sin(\sigma_{2200}\Delta t)=k_2\sin\big(2(\Op-n)\Delta t\big).
\]

\paragraph{Geometry.} The star is in the equator, $\theta^\star=\pi/2$. Hence $\dd\PK_{2m}/\dd\theta=0$ for the two surviving terms, and $F_\theta=0$. Moreover $\Lambda^\star_{2200}=\nu_{2200}-2\varphi^\star=0$: for a circular equatorial orbit the star's body-fixed longitude is $\varphi^\star=\omega+\Omega+\M-\thetap$, and $\nu_{2200}=2\omega+2\M+2(\Omega-\thetap)=2\varphi^\star$.

\paragraph{Azimuthal component.} With $\Lambda^\star=0$, the bracket of $F_\varphi$ in \eqref{eq:force-polar} reduces to $\Im\kb_2(\sigma_{2200})$. The other factors of the $(2,0,0)$ term are $\mathcal K_{2200}=\G M^\star\frac{R^5}{a^3}\cdot\tfrac1{12}\cdot3$ and $m\PK_{22}/\sin\theta|_{\pi/2}=2\cdot3$. Hence
\[
  F^{\star\to\mathrm p}_\varphi=-\frac{3}{2}\frac{\G M^{\star2}R^5}{a^7}\,k_2\sin\big(2(\Op-n)\Delta t\big),\qquad
  F^{\mathrm p\to\star}_\varphi=+\frac{3}{2}\frac{\G M^{\star2}R^5}{a^7}\,k_2\sin\big(2(\Op-n)\Delta t\big).
\]
\paragraph{Reading.}
\begin{itemize}
  \item If the planet spins faster than the orbit, $\Op>n$, the bulge leads the star. Then $F^{\mathrm p\to\star}_\varphi>0$: the bulge pulls the star forward, along its orbital motion ($\ep$ is eastward, the direction of the orbital motion for a prograde orbit).
  \item If $\Op<n$, the sign reverses.
  \item The radial component is $F^{\mathrm p\to\star}_r=-3k_2\G M^{\star2}R^5a^{-7}\big[\tfrac14+\tfrac34\cos2(\Op-n)\Delta t\big]<0$: an attraction.
\end{itemize}
These are the classical tidal signs \citep{Hut1981}. The same result is derived in general form in Chapter~\ref{ch:coplanar}.

\chapter{Second Kaula transformation: the force in orbital elements}\label{ch:second}

\section{Principle and identification of the elements}
The gradient \eqref{eq:grad-cart} is a sum of degree-$(l+1)$ solid harmonics of the field point,
\[
  \frac{\PK_{l+1,m'}(\cos\theta)\,e^{\pm\I m'\varphi}}{r^{l+2}},\qquad 0\le m'\le l+1.
\]
We describe the field point by its \emph{own} osculating elements $(a,e,i,\omega,\Omega,\M)$, and apply the degree-$(l+1)$ Kaula transformation \eqref{eq:kaula-complex} with indices $(k,j)$:
\begin{equation}
  \frac{\PK_{l+1,m'}(\cos\theta)e^{\I m'\varphi}}{r^{l+2}}=\frac{1}{a^{l+2}}\sum_{k=0}^{l+1}\sum_{j\in\mathbb Z}F_{l+1,m'k}(i)G_{l+1,kj}(e)e^{\I(\nu_{l+1,m'kj}+f_{l+1,m'})}, \label{eq:second-kaula}
\end{equation}
where $\nu_{l+1,m'kj}=(l+1-2k)\omega+(l+1-2k+j)\M+m'(\Omega-\thetap)$. The rotation angle $\thetap$ is the same in both transformations, since it is the same planet and the same body-fixed frame.

\paragraph{Identification $\vr=\vr^\star$.} The force is evaluated at the star, so the field point has the orbital elements of the star:
\begin{equation}
  a=a^\star,\quad e=e^\star,\quad i=i^\star,\quad \omega=\omega^\star,\quad \Omega=\Omega^\star,\quad \M=\M^\star. \label{eq:identification}
\end{equation}
From here on, the $\star$ on the elements is dropped. Three consequences:
\begin{enumerate}
  \item No geometric coordinate remains. The factor $1/r^{l+2}$ is represented exactly by $a^{-(l+2)}$ times the series in $G_{l+1,kj}(e)$; $r$ is never replaced by $a$.
  \item The lag $\epsilon_l$ and the argument of $\kb_l$ belong to the first transformation, with indices $(l,m,p,q)$. They do not depend on $(k,j)$.
  \item No secular simplification is made: all $(p,q,k,j)$ are kept.
\end{enumerate}
This order of operations — gradient with respect to the field point first, identification of the elements afterwards — is the one required to obtain correct perturbation equations \citep{EfroimskyMakarov2013}.

\section{Phases after identification}
Each product in the gradient combines a phase of the first transformation, $\nu_{lmpq}$ (degree $l$), with a phase of the second one, $\nu_{l+1,m',kj}$ (degree $l+1$). Their difference contains $m(\Omega-\thetap)-m'(\Omega-\thetap)$. The factor $e^{\pm\I\varphi}$ of \eqref{eq:grad-cart} gives $m'=m\pm1$, and the projection on the nodal frame (next section) supplies the missing $e^{\mp\I(\Omega-\thetap)}$. After both steps, $\Omega-\thetap$ cancels, and only two combinations of angles remain:
\begin{align}
  D_{pqkj}&\equiv\nu_{lmpq}-\nu_{l+1,m,kj}=(2k-2p-1)\,\omega+(2k-2p+q-j-1)\,\M, \label{eq:D}\\
  S_{pqkj}&\equiv\nu_{l0pq}+\nu_{l+1,1,kj}-(\Omega-\thetap)=(2l+1-2p-2k)\,\omega+(2l+1-2p-2k+q+j)\,\M. \label{eq:S}
\end{align}
$D$ is independent of $m$, $\Omega$ and $\thetap$. $S$, needed only for $m=0$, is independent of $\Omega$ and $\thetap$. Two further relations are used below:
\begin{equation}
  \nu_{l+1,m\pm1,kj}=\nu_{l+1,m,kj}\pm(\Omega-\thetap),\qquad\text{hence}\qquad \nu_{l+1,m\pm1,kj}-\nu_{lmpq}=-D_{pqkj}\pm(\Omega-\thetap). \label{eq:nu-shift}
\end{equation}

\paragraph{Parity phases.} By \eqref{eq:flm}, $f_{lm}$ depends only on the parity of $l-m$: $f_{lm}=0$ if $l-m$ is even, $-\pi/2$ if it is odd.
\begin{itemize}
  \item $(l+1)-(m\pm1)$ has the parity of $l-m$, so $f_{l+1,m\pm1}=f_{lm}$.
  \item $(l+1)-m$ has the opposite parity, so $f_{lm}-f_{l+1,m}=+\pi/2$ for $l-m$ even and $-\pi/2$ for $l-m$ odd:
  \begin{equation}
    f_{lm}-f_{l+1,m}=(-1)^{l-m}\,\frac{\pi}{2}. \label{eq:f-diff}
  \end{equation}
  The same holds for $f_{lm}-f_{l,m\pm1}$, since $l-(m\pm1)$ also has the opposite parity.
\end{itemize}
For any angle $x$, \eqref{eq:f-diff} gives
\begin{equation}
  \Re\,e^{\I(x+f_{lm}-f_{l+1,m})}=\cos\big(x+(-1)^{l-m}\tfrac\pi2\big)=(-1)^{l-m+1}\sin x. \label{eq:re-shift}
\end{equation}

\section{The force in the nodal frame}
\paragraph{The nodal frame.} The body-fixed axes $\mathbf e_x$, $\mathbf e_y$ still depend on $\Omega-\thetap$. We remove this dependence by projecting on the nodal frame:
\begin{itemize}
  \item $\mathbf e_Z$ is the spin axis;
  \item $\mathbf e_X$ points towards the ascending node of the orbit on the equator, i.e.\ towards the body-fixed longitude $\varphi=\Omega-\thetap$;
  \item rotating the axes by the angle $\Omega-\thetap$ about $\mathbf e_z$ gives
  \begin{equation}
    V_X+\I V_Y=e^{-\I(\Omega-\thetap)}\big(V_x+\I V_y\big)=e^{-\I(\Omega-\thetap)}\,V_+,\qquad V_Z=V_z. \label{eq:nodal-rot}
  \end{equation}
\end{itemize}

\paragraph{Method.} In each component of \eqref{eq:grad-cart}:
\begin{enumerate}
  \item write the real trigonometric function as the real part of an exponential;
  \item gather $\PK_{l+1,m'}(\cos\theta)$, $e^{\mp\I m'\varphi}$ and $1/r^{l+2}$ into one solid harmonic, and replace it with \eqref{eq:second-kaula} (or its complex conjugate);
  \item identify the elements \eqref{eq:identification} and collect the phases with \eqref{eq:D}--\eqref{eq:re-shift}.
\end{enumerate}
Recall that $X=\zeta_{lmpq}-m\varphi-\epsilon_l$ with $\zeta_{lmpq}=\nu_{lmpq}+f_{lm}$. Since $\PK$, $F$ and $G$ are real, $\Re$ can be moved through the sums.

\paragraph{Step 1: the $Z$-component.} With $\cos X=\Re\,e^{\I X}$,
\[
  \frac{\PK_{l+1,m}\cos X}{r^{l+2}}=\Re\Big[e^{\I(\zeta_{lmpq}-\epsilon_l)}\,\frac{\PK_{l+1,m}\,e^{-\I m\varphi}}{r^{l+2}}\Big].
\]
The second factor is the complex conjugate of \eqref{eq:second-kaula} with $m'=m$:
\[
  \frac{\PK_{l+1,m}\,e^{-\I m\varphi}}{r^{l+2}}=\frac{1}{a^{l+2}}\sum_{k,j}F_{l+1,mk}G_{l+1,kj}\,e^{-\I(\nu_{l+1,m,kj}+f_{l+1,m})}.
\]
The total phase is $\nu_{lmpq}-\nu_{l+1,m,kj}+f_{lm}-f_{l+1,m}-\epsilon_l=D_{pqkj}-\epsilon_l+(f_{lm}-f_{l+1,m})$. By \eqref{eq:re-shift},
\[
  \frac{\PK_{l+1,m}\cos X}{r^{l+2}}=\frac{(-1)^{l-m+1}}{a^{l+2}}\sum_{k,j}F_{l+1,mk}G_{l+1,kj}\sin(D_{pqkj}-\epsilon_l).
\]
Multiplying by $(l-m+1)|\kb_l|\mathcal K_{lmpq}$ gives the first line of \eqref{eq:grad-node}.

\paragraph{Step 2: the horizontal components, $m\ge1$.} The bracket of \eqref{eq:grad-cart} contains two terms.
\begin{itemize}
  \item \emph{Order $m+1$.} Since $e^{\I\varphi}e^{-\I X}=e^{-\I(\zeta_{lmpq}-\epsilon_l)}e^{\I(m+1)\varphi}$,
  \begin{align*}
    \frac{e^{\I\varphi}\PK_{l+1,m+1}e^{-\I X}}{r^{l+2}}&=e^{-\I(\zeta_{lmpq}-\epsilon_l)}\frac{\PK_{l+1,m+1}e^{\I(m+1)\varphi}}{r^{l+2}}\\
    &=\frac{1}{a^{l+2}}\sum_{k,j}F_{l+1,m+1,k}G_{l+1,kj}\,e^{\I(\nu_{l+1,m+1,kj}-\nu_{lmpq}+\epsilon_l)},
  \end{align*}
  where \eqref{eq:second-kaula} was used with $m'=m+1$ and $f_{l+1,m+1}=f_{lm}$. By \eqref{eq:nu-shift}, the phase is $-(D_{pqkj}-\epsilon_l)+(\Omega-\thetap)$.
  \item \emph{Order $m-1$.} Since $e^{\I\varphi}e^{\I X}=e^{\I(\zeta_{lmpq}-\epsilon_l)}e^{-\I(m-1)\varphi}$, the conjugate of \eqref{eq:second-kaula} with $m'=m-1$ and $f_{l+1,m-1}=f_{lm}$ gives
  \[
    \frac{e^{\I\varphi}\PK_{l+1,m-1}e^{\I X}}{r^{l+2}}=\frac{1}{a^{l+2}}\sum_{k,j}F_{l+1,m-1,k}G_{l+1,kj}\,e^{\I(\nu_{lmpq}-\nu_{l+1,m-1,kj}-\epsilon_l)},
  \]
  with phase $(D_{pqkj}-\epsilon_l)+(\Omega-\thetap)$ by \eqref{eq:nu-shift}.
\end{itemize}
Both terms carry the same factor $e^{\I(\Omega-\thetap)}$. The rotation \eqref{eq:nodal-rot} removes it, which gives the second line of \eqref{eq:grad-node}. For $m=l$, the order $m+1=l+1$ is still allowed at degree $l+1$.

\paragraph{Step 3: the horizontal components, $m=0$.} Here $(\nabla T_{l0pq})_+\propto e^{\I\varphi}\PK_{l+1,1}\cos X$ with $X=\zeta_{l0pq}-\epsilon_l$. Write $\cos X=\tfrac12(e^{-\I X}+e^{\I X})$ and use \eqref{eq:second-kaula} with $m'=1$:
\[
  \frac{e^{\I\varphi}\PK_{l+1,1}\cos X}{r^{l+2}}=\frac{1}{2a^{l+2}}\sum_{k,j}F_{l+1,1k}G_{l+1,kj}\,e^{\I(\nu_{l+1,1,kj}+f_{l+1,1})}\big[e^{-\I(\zeta_{l0pq}-\epsilon_l)}+e^{\I(\zeta_{l0pq}-\epsilon_l)}\big].
\]
\begin{itemize}
  \item First term: $f_{l+1,1}=f_{l0}$, and the phase is $\nu_{l+1,1,kj}-\nu_{l0pq}+\epsilon_l=-(D_{pqkj}-\epsilon_l)+(\Omega-\thetap)$.
  \item Second term: the phase is $\nu_{l0pq}+\nu_{l+1,1,kj}-\epsilon_l+2f_{l0}=S_{pqkj}-\epsilon_l+(\Omega-\thetap)+2f_{l0}$ by \eqref{eq:S}, and $e^{2\I f_{l0}}=(-1)^l$.
\end{itemize}
After the rotation \eqref{eq:nodal-rot}, this gives the third line of \eqref{eq:grad-node}.

\paragraph{Result.} With $\mathcal K_{lmpq}$ from \eqref{eq:Klmpq}, evaluated with the identified elements,
\begin{equation}
  \boxed{\begin{aligned}
  (\nabla T_{lmpq})_Z={}&(-1)^{l-m+1}(l-m+1)\frac{|\kb_l|\mathcal K_{lmpq}}{a^{l+2}}\\
  &\qquad\times\sum_{k=0}^{l+1}\sum_jF_{l+1,mk}G_{l+1,kj}\sin(D_{pqkj}-\epsilon_l),\\
  (\nabla T_{lmpq})_X+\I(\nabla T_{lmpq})_Y={}&\frac{|\kb_l|\mathcal K_{lmpq}}{2a^{l+2}}\sum_{k,j}G_{l+1,kj}\Big[F_{l+1,m+1,k}e^{-\I(D_{pqkj}-\epsilon_l)}\\
  &\qquad-(l-m+1)(l-m+2)F_{l+1,m-1,k}e^{\I(D_{pqkj}-\epsilon_l)}\Big],\quad m\ge1,\\
  (\nabla T_{l0pq})_X+\I(\nabla T_{l0pq})_Y={}&\frac{|\kb_l|\mathcal K_{l0pq}}{2a^{l+2}}\sum_{k,j}F_{l+1,1k}G_{l+1,kj}\\
  &\qquad\times\Big[e^{-\I(D_{pqkj}-\epsilon_l)}+(-1)^le^{\I(S_{pqkj}-\epsilon_l)}\Big].
  \end{aligned}} \label{eq:grad-node}
\end{equation}
where:
\begin{itemize}
  \item $V_Z$ is along the spin axis and $V_X+\I V_Y$ in the equator, $\mathbf e_X$ pointing to the ascending node~\eqref{eq:nodal-rot};
  \item $a,e,i,\omega,\M$ are the elements of the star's orbit on the planet's equator (identified with those of the field point,~\eqref{eq:identification});
  \item $\mathcal K_{lmpq}$ is the coefficient~\eqref{eq:Klmpq}; $F_{l+1,m'k}(i)$ and $G_{l+1,kj}(e)$ are the degree-$(l+1)$ Kaula functions of the second transformation;
  \item $D_{pqkj}$ and $S_{pqkj}$ are the combined angles~\eqref{eq:D}--\eqref{eq:S};
  \item $|\kb_l|$, $\epsilon_l$ are the modulus and lag of the Love number at $\sigma_{lmpq}$, the frequency of the first transformation.
\end{itemize}
These components depend only on $(a,e,i,\omega,\M)$. The node longitude $\Omega$ and the rotation angle $\thetap$ enter only through $\kb_l(\sigma_{lmpq})$.

\paragraph{Orbital components.} With $u=\omega+f$, the radial, transverse and normal components are
\begin{equation}
  \begin{aligned}
  V_{\rm R}&=V_X\cos u+V_Y\sin u\cos i+V_Z\sin u\sin i,\\
  V_{\rm T}&=-V_X\sin u+V_Y\cos u\cos i+V_Z\cos u\sin i,\\
  V_{\rm N}&=-V_Y\sin i+V_Z\cos i.
  \end{aligned} \label{eq:RTN}
\end{equation}

\section{The force in the local spherical basis}
The local basis follows the star, so two geometric quantities necessarily remain in it. Both are functions of the elements through the true anomaly $f(\M,e)$:
\begin{equation}
  r^\star=\frac{a(1-e^2)}{1+e\cos f},\qquad \cos\theta^\star=\sin i\sin u,\qquad \varphi^\star=(\Omega-\thetap)+\lambda_N, \label{eq:star-geometry}
\end{equation}
where $\lambda_N$ is the longitude of the star on the equator, measured from the node:
\[
  \cos\lambda_N=\frac{\cos u}{\sin\theta^\star},\qquad \sin\lambda_N=\frac{\cos i\sin u}{\sin\theta^\star}
\]
(Kaula's orbit--equator--meridian triangle, \citealt[Fig.~4]{Kaula2000}).

\paragraph{Which degree for which component.} We start from the polar form \eqref{eq:grad-polar-real} evaluated at the star, $X^\star=\zeta_{lmpq}-m\varphi^\star-\epsilon_l$, with $1/r^{l+2}=(1/r^\star)\cdot1/(r^\star)^{l+1}$.
\begin{itemize}
  \item The radial and meridional components contain the degree-$l$ functions $\PK_{lm}$ and $\PK_{l,m\pm1}$ of \eqref{eq:dPK-main}. They are transformed at degree $l$. This absorbs $1/(r^\star)^{l+1}$ and leaves a factor $1/r^\star$.
  \item The azimuthal component is the degree-$(l+1)$ combination \eqref{eq:mPK-main}. It is transformed at degree $l+1$, which absorbs the whole $1/(r^\star)^{l+2}$.
\end{itemize}
In the degree-$l$ transformation, the indices of the second expansion are again called $(k,j)$, and the phase differences are
\begin{equation}
  \begin{gathered}
  \nu_{lmpq}-\nu_{l,m,kj}=\alpha_{pqkj},\qquad \nu_{lmpq}-\nu_{l,m\pm1,kj}=\alpha_{pqkj}\mp(\Omega-\thetap),\\
  \alpha_{pqkj}\equiv(2k-2p)\omega+(2k-2p+q-j)\M.
  \end{gathered} \label{eq:alpha}
\end{equation}
Each $e^{\mp\I\varphi^\star}$ left over is split with \eqref{eq:star-geometry}: $\varphi^\star=(\Omega-\thetap)+\lambda_N$, and the part $\Omega-\thetap$ cancels the one in \eqref{eq:alpha}.

\paragraph{Step 4: radial component.} With $\cos X^\star=\Re\,e^{\I X^\star}$ and the conjugate of \eqref{eq:kaula-complex} at order $m$,
\[
  \frac{\PK_{lm}(\cos\theta^\star)\cos X^\star}{(r^\star)^{l+1}}=\Re\Big[e^{\I(\zeta_{lmpq}-\epsilon_l)}\frac{\PK_{lm}e^{-\I m\varphi^\star}}{(r^\star)^{l+1}}\Big]=\frac{1}{a^{l+1}}\sum_{k,j}F_{lmk}G_{lkj}\cos(\alpha_{pqkj}-\epsilon_l).
\]
Both $f_{lm}$ and $m(\Omega-\thetap)$ cancel. Multiplying by $(l+1)|\kb_l|\mathcal K_{lmpq}/r^\star$ gives \eqref{eq:local-r}.

\paragraph{Step 5: meridional component.} By \eqref{eq:dPK-main}, $\dd\PK_{lm}/\dd\theta$ contains $\PK_{l,m-1}$ and $\PK_{l,m+1}$. Insert $e^{-\I(m\mp1)\varphi^\star}e^{\mp\I\varphi^\star}=e^{-\I m\varphi^\star}$:
\[
  \PK_{l,m\mp1}\cos X^\star=\Re\Big[e^{\I(\zeta_{lmpq}-\epsilon_l)}\,e^{\mp\I\varphi^\star}\,\PK_{l,m\mp1}e^{-\I(m\mp1)\varphi^\star}\Big].
\]
Apply the conjugate of \eqref{eq:kaula-complex} at order $m\mp1$. The phase is
\[
  \nu_{lmpq}-\nu_{l,m\mp1,kj}\mp\varphi^\star-\epsilon_l+(f_{lm}-f_{l,m\mp1})=\alpha_{pqkj}\mp\lambda_N-\epsilon_l+(-1)^{l-m}\tfrac\pi2,
\]
by \eqref{eq:alpha}, \eqref{eq:star-geometry} and \eqref{eq:f-diff}. By \eqref{eq:re-shift},
\[
  \frac{\PK_{l,m\mp1}\cos X^\star}{(r^\star)^{l+1}}=\frac{(-1)^{l-m+1}}{a^{l+1}}\sum_{k,j}F_{l,m\mp1,k}G_{lkj}\sin(\alpha_{pqkj}\mp\lambda_N-\epsilon_l).
\]
The prefactor of $(\nabla T)_\theta$ in \eqref{eq:grad-polar-real} is $-|\kb_l|\mathcal K_{lmpq}/r^{l+2}$, and \eqref{eq:dPK-main} gives the weights $+\tfrac12(l+m)(l-m+1)$ for $m-1$ and $-\tfrac12$ for $m+1$. The two minus signs combine into $(-1)^{l-m}$, which gives \eqref{eq:local-theta}.

For $m=0$, $\dd\PK_{l0}/\dd\theta=-\PK_{l1}$. Only the order $1$ appears, with the lower sign ($+\lambda_N$), and $f_{l0}-f_{l1}=(-1)^l\pi/2$. Hence the $m=0$ formula below \eqref{eq:local-phi}.

\paragraph{Step 6: azimuthal component.} With $\sin X^\star=\Im\,e^{\I X^\star}$ and \eqref{eq:mPK-main},
\[
  \PK_{l+1,m\pm1}\sin X^\star=\Im\Big[e^{\I(\zeta_{lmpq}-\epsilon_l)}\,e^{\pm\I\varphi^\star}\,\PK_{l+1,m\pm1}e^{-\I(m\pm1)\varphi^\star}\Big].
\]
Apply the conjugate of \eqref{eq:second-kaula} at order $m\pm1$, with $f_{l+1,m\pm1}=f_{lm}$. By \eqref{eq:nu-shift} and \eqref{eq:star-geometry}, the phase is
\[
  \nu_{lmpq}-\nu_{l+1,m\pm1,kj}\pm\varphi^\star-\epsilon_l=D_{pqkj}\mp(\Omega-\thetap)\pm(\Omega-\thetap)\pm\lambda_N-\epsilon_l=D_{pqkj}\pm\lambda_N-\epsilon_l.
\]
Hence
\[
  \frac{\PK_{l+1,m\pm1}\sin X^\star}{(r^\star)^{l+2}}=\frac{1}{a^{l+2}}\sum_{k,j}F_{l+1,m\pm1,k}G_{l+1,kj}\sin(D_{pqkj}\pm\lambda_N-\epsilon_l).
\]
With the prefactor $-|\kb_l|\mathcal K_{lmpq}$ and the weights $\tfrac12$ and $\tfrac12(l-m+1)(l-m+2)$ of \eqref{eq:mPK-main}, this gives \eqref{eq:local-phi}. No factor $1/r^\star$ remains.

\paragraph{Result.}
\begin{align}
  (\nabla T_{lmpq})_r&=(l+1)\frac{|\kb_l|\mathcal K_{lmpq}}{r^\star a^{l+1}}\sum_{k=0}^{l}\sum_jF_{lmk}G_{lkj}\cos(\alpha_{pqkj}-\epsilon_l), \label{eq:local-r}\\
  (\nabla T_{lmpq})_\theta&=(-1)^{l-m}\frac{|\kb_l|\mathcal K_{lmpq}}{2r^\star a^{l+1}}\sum_{k,j}G_{lkj}\Big[(l+m)(l-m+1)F_{l,m-1,k}\sin(\alpha_{pqkj}-\lambda_N-\epsilon_l)\notag\\
  &\hspace{12em}-F_{l,m+1,k}\sin(\alpha_{pqkj}+\lambda_N-\epsilon_l)\Big], \label{eq:local-theta}\\
  (\nabla T_{lmpq})_\varphi&=-\frac{|\kb_l|\mathcal K_{lmpq}}{2a^{l+2}}\sum_{k=0}^{l+1}\sum_jG_{l+1,kj}\Big[F_{l+1,m+1,k}\sin(D_{pqkj}+\lambda_N-\epsilon_l)\notag\\
  &\hspace{9em}+(l-m+1)(l-m+2)F_{l+1,m-1,k}\sin(D_{pqkj}-\lambda_N-\epsilon_l)\Big]. \label{eq:local-phi}
\end{align}
Here $r^\star$ is the distance of the star, $\lambda_N$ its longitude counted from the node~\eqref{eq:star-geometry}, $\alpha_{pqkj}$ and $D_{pqkj}$ the combined angles~\eqref{eq:alpha} and~\eqref{eq:D}, and $(\er,\et,\ep)$ the local basis at the star. The last two hold for $m\ge1$, with the convention $F_{l,l+1,k}\equiv0$ (for $m=l$, the order $l+1$ does not exist at degree $l$). For $m=0$:
\[
  (\nabla T_{l0pq})_\theta=-(-1)^l\frac{|\kb_l|\mathcal K_{l0pq}}{r^\star a^{l+1}}\sum_{k,j}F_{l1k}G_{lkj}\sin(\alpha_{pqkj}+\lambda_N-\epsilon_l),\qquad (\nabla T_{l0pq})_\varphi=0.
\]
The force is $\Fsp=+M^\star\sum_{lmpq}\nabla T_{lmpq}$. The orbital components follow from
\[
  F_{\rm R}=F_r,\qquad F_{\rm T}=-bF_\theta+cF_\varphi,\qquad F_{\rm N}=-cF_\theta-bF_\varphi,
\]
with $b=\sin i\cos u/\sin\theta^\star$ and $c=\cos i/\sin\theta^\star$.

\section{Degree two}
Here $\mathcal K_{2mpq}=\G M^\star\frac{R^5}{a^3}c_mF_{2mp}G_{2pq}$, with $c_0=1$, $c_1=\tfrac13$, $c_2=\tfrac1{12}$, and $\epsilon=\epsilon_2(\sigma_{2mpq})$:
\begin{align}
  (\nabla T_{2mpq})_r&=\frac{3|\kb_2|\mathcal K_{2mpq}}{r^\star a^3}\sum_{k=0}^{2}\sum_jF_{2mk}G_{2kj}\cos(\alpha_{pqkj}-\epsilon),\\
  (\nabla T_{20pq})_\theta&=-\frac{|\kb_2|\mathcal K_{20pq}}{r^\star a^3}\sum_{k,j}F_{21k}G_{2kj}\sin(\alpha_{pqkj}+\lambda_N-\epsilon),\\
  (\nabla T_{21pq})_\theta&=\frac{|\kb_2|\mathcal K_{21pq}}{r^\star a^3}\sum_{k,j}G_{2kj}\big[-3F_{20k}\sin(\alpha_{pqkj}-\lambda_N-\epsilon)+\tfrac12F_{22k}\sin(\alpha_{pqkj}+\lambda_N-\epsilon)\big],\\
  (\nabla T_{22pq})_\theta&=\frac{2|\kb_2|\mathcal K_{22pq}}{r^\star a^3}\sum_{k,j}F_{21k}G_{2kj}\sin(\alpha_{pqkj}-\lambda_N-\epsilon),\\
  (\nabla T_{21pq})_\varphi&=-\frac{|\kb_2|\mathcal K_{21pq}}{a^4}\sum_{k=0}^{3}\sum_jG_{3kj}\big[\tfrac12F_{32k}\sin(D_{pqkj}+\lambda_N-\epsilon)+3F_{30k}\sin(D_{pqkj}-\lambda_N-\epsilon)\big],\\
  (\nabla T_{22pq})_\varphi&=-\frac{|\kb_2|\mathcal K_{22pq}}{a^4}\sum_{k=0}^{3}\sum_jG_{3kj}\big[\tfrac12F_{33k}\sin(D_{pqkj}+\lambda_N-\epsilon)+F_{31k}\sin(D_{pqkj}-\lambda_N-\epsilon)\big],
\end{align}
and $(\nabla T_{20pq})_\varphi=0$.

\chapter{The coplanar case}\label{ch:coplanar}

The only approximation of this chapter is $i=0$: the star's orbit lies in the planet's equator, i.e.\ the obliquity is zero. No secular averaging is made.

\section{Inclination functions, phases and coefficients}
At $i=0$ only the terms with $l-m-2t=0$ survive in Kaula's eq.~(3.62) (Appendix~\ref{app:kaula}):
\begin{equation}
  F_{lmp}(0)=\begin{cases}\PK_{lm}(0)=\dfrac{(-1)^{(l-m)/2}(l+m)!}{2^l\big(\frac{l-m}{2}\big)!\big(\frac{l+m}{2}\big)!},& l-m\ \text{even},\ p=p_0\equiv\frac{l-m}{2},\\[2mm]0,&\text{otherwise}.\end{cases} \label{eq:F0}
\end{equation}
For example $F_{201}(0)=-\tfrac12$, $F_{220}(0)=3$, $F_{311}(0)=-\tfrac32$, $F_{330}(0)=15$. Consequences:
\begin{itemize}
  \item only the orders $m\equiv l\pmod2$ raise tides, with $p=p_0$;
  \item the star stays in the equator ($\theta^\star=\pi/2$), so the meridional component vanishes identically;
  \item $\ep$ is the transverse direction.
\end{itemize}
The node is undefined at $i=0$. With the longitude of pericentre $\varpi=\Omega+\omega$,
\begin{equation}
  \nu_{lmp_0q}=m(\varpi-\thetap)+(m+q)\M,\qquad \sigma_{lmp_0q}=m(\dot\varpi-\Op)+(m+q)n. \label{eq:nu-coplanar}
\end{equation}
where $\varpi=\Omega+\omega$ is the longitude of pericentre, $\M$ the mean anomaly, $n$ the mean motion, $\thetap$ and $\Op$ the rotation angle and spin rate of the planet. At $l=2$ this is $\sigma_{201q}=qn$ and $\sigma_{220q}=(2+q)n-2\Op+2\dot\varpi$, the usual coplanar form; Appendix~\ref{app:frequency} shows how it relates to the general form when $i\neq0$.
The coefficients combine into
\begin{equation}
  C_{lm}\equiv\frac{(l-m)!}{(l+m)!}(2-\delta_{0,m})\big[\PK_{lm}(0)\big]^2,\qquad C_{20}=\tfrac14,\ C_{22}=\tfrac34,\ C_{31}=\tfrac38,\ C_{33}=\tfrac58, \label{eq:Clm}
\end{equation}
which satisfy $\sum_mC_{lm}=1$ (the addition theorem at $\Theta=0$ for points in the equator).

\section{General degree}
\paragraph{Double-sum form.} With $p=k=p_0$, we have $\alpha_{p_0qp_0j}=(q-j)\M$. Moreover, since $\sin\theta^\star=1$, the azimuthal component can be written at degree $l$ through $m\PK_{lm}/\sin\theta|_{\pi/2}=m\PK_{lm}(0)$. Then
\begin{equation}
  \begin{aligned}
  (\nabla T_{lmp_0q})_r&=(l+1)\frac{\G M^\star R^{2l+1}}{r^\star a^{2l+2}}C_{lm}G_{lp_0q}(e)\sum_jG_{lp_0j}(e)\,|\kb_l|\cos\big((q-j)\M-\epsilon_l\big),\\
  (\nabla T_{lmp_0q})_\theta&=0,\\
  (\nabla T_{lmp_0q})_\varphi&=-m\frac{\G M^\star R^{2l+1}}{r^\star a^{2l+2}}C_{lm}G_{lp_0q}(e)\sum_jG_{lp_0j}(e)\,|\kb_l|\sin\big((q-j)\M-\epsilon_l\big).
  \end{aligned} \label{eq:coplanar-A}
\end{equation}
where $r^\star$ is the distance of the star, $a,e,\M$ its semi-major axis, eccentricity and mean anomaly, $C_{lm}$ the coefficient~\eqref{eq:Clm}, $p_0=(l-m)/2$, and $|\kb_l|$, $\epsilon_l$ the modulus and lag of the Love number at the frequency $\sigma_{lmp_0q}$ of~\eqref{eq:nu-coplanar}.

\paragraph{Closed form.} Kaula's definition of $G$ \citep[p.~35]{Kaula2000}, read backwards, gives
\[
  \sum_jG_{lp_0j}e^{-\I j\M}=\Big(\frac{a}{r^\star}\Big)^{l+1}e^{\I m(\M-f)}.
\]
The $j$-sum is therefore the instantaneous position of the star:
\begin{equation}
  \boxed{\begin{aligned}
  (\nabla T_{lmp_0q})_r&=(l+1)\frac{\G M^\star R^{2l+1}}{a^{l+1}(r^\star)^{l+2}}C_{lm}G_{lp_0q}(e)\,|\kb_l|\cos\big((m+q)\M-mf-\epsilon_l\big),\\
  (\nabla T_{lmp_0q})_\varphi&=-m\frac{\G M^\star R^{2l+1}}{a^{l+1}(r^\star)^{l+2}}C_{lm}G_{lp_0q}(e)\,|\kb_l|\sin\big((m+q)\M-mf-\epsilon_l\big).
  \end{aligned}} \label{eq:coplanar-B}
\end{equation}
with the same notation, $f$ being the true anomaly of the star. Here $(m+q)\M-mf=\nu-m\varphi^\star$. Since $\varphi^\star=\varpi-\thetap+f$, the angles $\varpi$ and $\thetap$ cancel; they remain only in $\sigma$. The force is
\[
  \Fsp=M^\star\sum_{l}\sum_{m\equiv l\,(2)}\sum_q\big[(\nabla T)_r\er+(\nabla T)_\varphi\ep\big].
\]

\section{Degree two}
Let $\kb_{20q}=\kb_2(\sigma_{201q})$ with $\sigma_{201q}=qn$, and $\kb_{22q}=\kb_2(\sigma_{220q})$ with $\sigma_{220q}=2(\dot\varpi-\Op)+(2+q)n$; their lags are $\epsilon_{20q}$ and $\epsilon_{22q}$. Then
\begin{equation}
  \boxed{\begin{aligned}
  (\nabla T_2)_r&=\frac{\G M^\star R^5}{a^3(r^\star)^4}\sum_{q}\Big[\tfrac34G_{21q}(e)|\kb_{20q}|\cos(q\M-\epsilon_{20q})\\
  &\hspace{8em}+\tfrac94G_{20q}(e)|\kb_{22q}|\cos\big((2+q)\M-2f-\epsilon_{22q}\big)\Big],\\
  (\nabla T_2)_\theta&=0,\\
  (\nabla T_2)_\varphi&=-\frac32\frac{\G M^\star R^5}{a^3(r^\star)^4}\sum_qG_{20q}(e)|\kb_{22q}|\sin\big((2+q)\M-2f-\epsilon_{22q}\big),
  \end{aligned}} \label{eq:coplanar-l2}
\end{equation}
and $\Fsp=M^\star\nabla T_2$. Here $r^\star$ and $f$ are the distance and true anomaly of the star, $a,e,\M$ its elements, $R$ the radius of the planet; $\kb_{20q}$ acts at the frequency $qn$ (the $m=0$ tide) and $\kb_{22q}$ at $(2+q)n-2\Op+2\dot\varpi$ (the $m=2$ tide).

\paragraph{Classical limit.} Take a circular orbit and a constant time lag $\Delta t$: $q=0$, $\epsilon_{200}=0$, $\epsilon_{220}=2(n-\Op)\Delta t$. Then
\[
  \Fps=-\frac{3k_2\G M^{\star2}R^5}{a^7}\Big[\frac14+\frac34\cos2(\Op-n)\Delta t\Big]\er+\frac{3k_2\G M^{\star2}R^5}{2a^7}\sin2(\Op-n)\Delta t\,\ep,
\]
the standard result \citep{Hut1981}.

\chapter{Use in an $N$-body code}\label{ch:nbody}

This chapter gives the equations to be implemented in an $N$-body integrator for the tide raised by the star in the planet, at degree two. It does not depend on a particular code. The tide raised in the star uses the same equations (Chapter~\ref{ch:star}), and the implementation in Posidonius is the subject of Chapter~\ref{ch:posidonius}. Every formula of this chapter was checked numerically against an independent computation of the exact linear response of the planet, which does not use Kaula's expansion (Section~\ref{sec:validation}).

\section{Inputs, outputs and frame of the force}\label{sec:io}
\paragraph{Inputs,} all in inertial axes, at the time $t$ of the force evaluation:
\begin{itemize}
  \item $\vr=\vr_{\mathrm p}-\vr_\star$ and $\mathbf v=\mathbf v_{\mathrm p}-\mathbf v_\star$, position and velocity of the planet relative to the star;
  \item $\mathbf s$, spin vector of the planet;
  \item $M^\star$, $M_{\mathrm p}$, $R$, and the gravitational constant $\G$;
  \item the Love-number spectra $\texttt{rek2}(\chi)$, $\texttt{imk2}(\chi)$ of the planet (Section~\ref{sec:love});
  \item the history of the accepted steps, from which the mean rates $\bar n$, $\dot{\bar\lambda}$, $\dot{\bar\Omega}$ are computed (Section~\ref{sec:rates}), and the mean motions of the other orbiting bodies, which fix the averaging window.
\end{itemize}

\paragraph{Outputs.}
\begin{itemize}
  \item The force of the star on the planet, $\Fsp=M^\star\nabla T|_{\vr^\star}$, in inertial axes (convention~\eqref{eq:force-convention}).
  \item The torque on the spin, $\boldsymbol\Gamma=\dd(I\mathbf s)/\dd t$ (Section~\ref{sec:torque}).
\end{itemize}

\paragraph{How the force is applied.} The tidal interaction is an internal force of the star--planet pair:
\begin{itemize}
  \item the planet receives $+\Fsp/M_{\mathrm p}$ and the star $-\Fsp/M^\star$, in the barycentric frame;
  \item the relative acceleration is then $\Fsp(1/M_{\mathrm p}+1/M^\star)$ automatically;
  \item in a multi-planet system, the reaction on the star produces the indirect terms in the motion of the other planets.
\end{itemize}
No extra factor and no indirect term must be added by hand.

\section{Geometry from the state vectors}\label{sec:geometry}
All geometric quantities are computed from $(\vr,\mathbf v,\mathbf s)$. No Euler angle of the spin axis and no inertial reference direction is needed.

\paragraph{Orbit of the star seen from the planet.} The star is at $\vr^\star=-\vr$, with velocity $\mathbf v^\star=-\mathbf v$. Its orbit has the angular momentum per unit reduced mass $\mathbf h=\vr^\star\times\mathbf v^\star=\vr\times\mathbf v$, the same as that of the planet. With $\mu_{\rm orb}=\G(M^\star+M_{\mathrm p})$ and $r^\star=\|\vr^\star\|$,
\begin{equation}
  a=\Big(\frac{2}{r^\star}-\frac{\|\mathbf v\|^2}{\mu_{\rm orb}}\Big)^{-1},\qquad
  \mathbf e_{\rm v}=\frac{\mathbf v^\star\times\mathbf h}{\mu_{\rm orb}}-\frac{\vr^\star}{r^\star},\qquad
  e=\|\mathbf e_{\rm v}\|,\qquad \bar n_{\rm osc}=\sqrt{\mu_{\rm orb}/a^3}. \label{eq:elements-a-e}
\end{equation}

\paragraph{Nodal frame.} Let $\hat{\mathbf s}=\mathbf s/\|\mathbf s\|$ be the spin axis, $\Op=\|\mathbf s\|$ the spin rate, and $\hat{\mathbf h}=\mathbf h/\|\mathbf h\|$ the orbit normal. The node of the orbit on the equator is the intersection of the two planes, i.e.\ the direction $\hat{\mathbf s}\times\hat{\mathbf h}$, whose norm is $\sin i$. The frame is built in three steps.
\begin{enumerate}
  \item \textbf{Direction of the node.}
  \[
    \mathbf N_0=\begin{cases}\hat{\mathbf s}\times\hat{\mathbf h}, & \text{if}\ \|\hat{\mathbf s}\times\hat{\mathbf h}\|\ge\varepsilon_N,\\[1mm] \hat{\mathbf s}\times\mathbf e_k, & \text{otherwise}.\end{cases}
  \]
  $\mathbf e_k$ is the inertial axis with the smallest $|\hat s_k|$, and $\varepsilon_N=10^{-10}$.
  \item \textbf{Projection onto the equator and normalisation.}
  \[
    \mathbf N_\perp=\mathbf N_0-(\mathbf N_0\cdot\hat{\mathbf s})\,\hat{\mathbf s},\qquad
    \hat{\mathbf N}=\frac{\mathbf N_\perp}{\|\mathbf N_\perp\|}.
  \]
  \item \textbf{Third axis.}
  \[
    \hat{\mathbf Y}=\hat{\mathbf s}\times\hat{\mathbf N}.
  \]
\end{enumerate}
In compact form,
\begin{equation}
  \hat{\mathbf N}=\frac{\mathbf N_0-(\mathbf N_0\cdot\hat{\mathbf s})\,\hat{\mathbf s}}{\big\|\mathbf N_0-(\mathbf N_0\cdot\hat{\mathbf s})\,\hat{\mathbf s}\big\|},\qquad
  \hat{\mathbf Y}=\hat{\mathbf s}\times\hat{\mathbf N}. \label{eq:node-frame}
\end{equation}
\begin{itemize}
  \item $(\hat{\mathbf N},\hat{\mathbf Y},\hat{\mathbf s})$ is the nodal frame of Chapter~\ref{ch:second}: $\hat{\mathbf N}$ points to the ascending node of the star's orbit on the equator.
  \item \textbf{The denominator.} In the general case, $\mathbf N_0=\hat{\mathbf s}\times\hat{\mathbf h}$ is already orthogonal to $\hat{\mathbf s}$, so $\mathbf N_\perp=\mathbf N_0$ and $\|\mathbf N_\perp\|=\|\hat{\mathbf s}\times\hat{\mathbf h}\|=\sin i\ge\varepsilon_N$. In the other case, $\|\mathbf N_\perp\|=\|\hat{\mathbf s}\times\mathbf e_k\|=\sqrt{1-\hat s_k^2}\ge\sqrt{2/3}$, because the smallest component of a unit vector satisfies $\hat s_k^2\le1/3$. The denominator is therefore never zero.
  \item \textbf{The projection.} Mathematically $\mathbf N_\perp=\mathbf N_0$. The projection only removes the rounding error, so that $\hat{\mathbf N}\cdot\hat{\mathbf s}=0$ to machine precision.
  \item \textbf{The threshold.} When $\hat{\mathbf s}\parallel\hat{\mathbf h}$ ($i=0$ or $\pi$), the node is undefined and any equatorial direction can be used: the force does not depend on it (Section~\ref{sec:validation}, test V7). The threshold $\varepsilon_N$ only avoids dividing by a vanishing $\sin i$; it does not switch to another formula for the force.
\end{itemize}

\paragraph{Inclination and angles in the orbital plane.}
\begin{equation}
  i=\operatorname{atan2}\big(\|\hat{\mathbf s}\times\hat{\mathbf h}\|,\ \hat{\mathbf h}\cdot\hat{\mathbf s}\big)\in[0,\pi],\qquad
  \hat{\mathbf W}=\frac{\hat{\mathbf h}\times\hat{\mathbf N}}{\|\hat{\mathbf h}\times\hat{\mathbf N}\|},\qquad
  \hat{\mathbf N}'=\hat{\mathbf W}\times\hat{\mathbf h}. \label{eq:inclination}
\end{equation}
$\hat{\mathbf N}'$ and $\hat{\mathbf W}$ span the orbital plane, with $\hat{\mathbf N}'=\hat{\mathbf N}$ when the node is defined. A retrograde rotation corresponds to $i>\pi/2$; no sign convention is needed for $\Op$.

\paragraph{Elements referred to the equator.} With $\hat{\vr}^\star=\vr^\star/r^\star$,
\begin{gather}
  u=\operatorname{atan2}\big(\hat{\vr}^\star\cdot\hat{\mathbf W},\ \hat{\vr}^\star\cdot\hat{\mathbf N}'\big),\qquad
  \omega=\operatorname{atan2}\big(\mathbf e_{\rm v}\cdot\hat{\mathbf W},\ \mathbf e_{\rm v}\cdot\hat{\mathbf N}'\big),\qquad f=u-\omega, \label{eq:elements-angles}\\
  E=\operatorname{atan2}\big(\sqrt{1-e^2}\,\sin f,\ e+\cos f\big),\qquad \M=E-e\sin E,\qquad u_\M\equiv\omega+\M. \notag
\end{gather}
\begin{itemize}
  \item $\omega$ is the argument of pericentre of the \emph{star's} orbit, referred to the planet's equator. It equals that of the planet plus $\pi$.
  \item $u_\M$ is the argument of the mean latitude. It is defined even when $e=0$: the force depends on $\omega$ and $\M$ only through $u_\M$ and through $\M$ multiplied by $G_{2pq}(e)=\mathcal O(e^{|q|})$. If $e=0$, set $\omega=0$.
\end{itemize}

\paragraph{Position of the star in the nodal frame.}
\begin{equation}
  \rho\equiv\hat{\vr}^\star\cdot\hat{\mathbf N}+\I\,\hat{\vr}^\star\cdot\hat{\mathbf Y}=\sin\theta^\star e^{\I\lambda_N},\qquad
  c^\star\equiv\cos\theta^\star=\hat{\vr}^\star\cdot\hat{\mathbf s}, \label{eq:rho}
\end{equation}
where $\theta^\star$ is the colatitude of the star and $\lambda_N$ its longitude counted from the node, as in~\eqref{eq:star-geometry}.

\section{From the gradient to the implementable form}\label{sec:implementable}
We start from the Cartesian gradient~\eqref{eq:grad-cart} at $l=2$, evaluated at the star. Three steps make it implementable: the change to the nodal frame, the elimination of $\Omega-\thetap$, and the polynomial form of the angular functions.

\paragraph{Notation for $l=2$.} With the elements of Section~\ref{sec:geometry},
\begin{equation}
  \mathcal K_{2mpq}=\G M^\star\frac{R^5}{a^3}\,c_m\,F_{2mp}(i)\,G_{2pq}(e),\qquad c_0=1,\ c_1=\tfrac13,\ c_2=\tfrac1{12}, \label{eq:K2}
\end{equation}
from~\eqref{eq:Klmpq}, and $\kb_2\equiv\kb_2(\bar\sigma_{2mpq})$ (Section~\ref{sec:love}). Only the terms with $F_{2mp}(i)\neq0$ and $|q|\le q_{\max}(e)$ (Section~\ref{sec:truncation}) are kept.

\paragraph{Step 1: phases at the star.} By~\eqref{eq:star-geometry}, the body-fixed longitude of the star is $\varphi^\star=(\Omega-\thetap)+\lambda_N$. In the phase of~\eqref{eq:grad-cart}, $\Lambda_{2mpq}=\zeta_{2mpq}-m\varphi^\star$, the term $m(\Omega-\thetap)$ of $\nu_{2mpq}$ cancels:
\begin{equation}
  \Lambda_{2mpq}=\Psi_{mpq}-m\lambda_N,\qquad \Psi_{mpq}\equiv(2-2p)\,\omega+(2-2p+q)\,\M+f_{2m}=(2-2p)\,u_\M+q\,\M+f_{2m}. \label{eq:Psi}
\end{equation}
The second form uses $(2-2p)\omega+(2-2p+q)\M=(2-2p)(\omega+\M)+q\M$ and is regular as $e\to0$.

\paragraph{Step 2: phase factor.} With $|\kb_2|e^{-\I\epsilon_2}=\kb_2$, the factor $|\kb_2|e^{\pm\I X}$ of~\eqref{eq:grad-cart}, $X=\Lambda-\epsilon_2$, becomes
\begin{equation}
  |\kb_2|\,e^{\I X}=\mathcal E_{mpq}\,e^{-\I m\lambda_N},\qquad \mathcal E_{mpq}\equiv\kb_2(\bar\sigma_{2mpq})\,e^{\I\Psi_{mpq}}, \label{eq:Ecal}
\end{equation}
and $|\kb_2|e^{-\I X}=\mathcal E^*_{mpq}e^{\I m\lambda_N}$.

\paragraph{Step 3: angular functions.} The degree-3 functions factor as
\[
  \PK_{3m'}(\cos\theta)=\sin^{m'}\theta\;\mathcal P_{3m'}(\cos\theta),\qquad \mathcal P_{3m'}=\frac{\dd^{m'}P_3}{\dd x^{m'}}
\]
(Appendix~\ref{app:harmonics}), with
\begin{equation}
  \mathcal P_{30}(x)=\tfrac12(5x^3-3x),\qquad \mathcal P_{31}(x)=\tfrac32(5x^2-1),\qquad \mathcal P_{32}(x)=15x,\qquad \mathcal P_{33}(x)=15. \label{eq:P3m}
\end{equation}
Hence, at the star, by~\eqref{eq:rho},
\[
  \PK_{3m'}(c^\star)\,e^{\pm\I m'\lambda_N}=\rho^{m'}\mathcal P_{3m'}(c^\star)\quad(+),\qquad \bar\rho^{\,m'}\mathcal P_{3m'}(c^\star)\quad(-).
\]
\begin{itemize}
  \item $Z$-component: $|\kb_2|\,\PK_{3m}\cos X=\Re\big[\PK_{3m}\,e^{-\I m\lambda_N}\mathcal E_{mpq}\big]=\mathcal P_{3m}\,\Re\big[\bar\rho^{\,m}\mathcal E_{mpq}\big]$.
  \item Order $m+1$: $|\kb_2|\,e^{\I\lambda_N}\PK_{3,m+1}e^{-\I X}=\PK_{3,m+1}\,e^{\I(m+1)\lambda_N}\mathcal E^*_{mpq}=\rho^{m+1}\mathcal P_{3,m+1}\,\mathcal E^*_{mpq}$.
  \item Order $m-1$: $|\kb_2|\,e^{\I\lambda_N}\PK_{3,m-1}e^{\I X}=\PK_{3,m-1}\,e^{-\I(m-1)\lambda_N}\mathcal E_{mpq}=\bar\rho^{\,m-1}\mathcal P_{3,m-1}\,\mathcal E_{mpq}$.
\end{itemize}
In the last two lines, $e^{\I\varphi^\star}=e^{\I(\Omega-\thetap)}e^{\I\lambda_N}$ was used; the factor $e^{\I(\Omega-\thetap)}$ is removed by the rotation~\eqref{eq:nodal-rot} to the nodal frame. Steps 1--3 were checked numerically term by term.

\paragraph{Result.} The gradient of $T$ at the star, in the nodal frame, is
\begin{equation}
  \boxed{\begin{aligned}
  V_Z&=\frac{1}{(r^\star)^4}\sum_{m,p,q}(3-m)\,\mathcal K_{2mpq}\,\mathcal P_{3m}(c^\star)\,\Re\big[\bar\rho^{\,m}\mathcal E_{mpq}\big],\\
  V_X+\I V_Y&=\frac{1}{(r^\star)^4}\sum_{p,q}\mathcal K_{20pq}\,\mathcal P_{31}(c^\star)\,\rho\,\Re\,\mathcal E_{0pq}\\
  &\quad+\frac{1}{2(r^\star)^4}\sum_{m=1}^{2}\sum_{p,q}\mathcal K_{2mpq}\Big[\mathcal P_{3,m+1}(c^\star)\,\rho^{m+1}\mathcal E^*_{mpq}\\
  &\hspace{9em}-(3-m)(4-m)\,\mathcal P_{3,m-1}(c^\star)\,\bar\rho^{\,m-1}\mathcal E_{mpq}\Big].
  \end{aligned}} \label{eq:impl-single}
\end{equation}
where:
\begin{itemize}
  \item $V_Z$ is the component of $\nabla T$ along the spin axis $\hat{\mathbf s}$, and $V_X+\I V_Y$ its components along $\hat{\mathbf N}$ and $\hat{\mathbf Y}$ (nodal frame~\eqref{eq:node-frame});
  \item $r^\star$ is the distance star--planet, $\rho=\sin\theta^\star e^{\I\lambda_N}$ and $c^\star=\cos\theta^\star$ give the direction of the star in the nodal frame~\eqref{eq:rho};
  \item $\mathcal K_{2mpq}$ is the real coefficient~\eqref{eq:K2}: masses, radius, $a$, $F_{2mp}(i)$, $G_{2pq}(e)$;
  \item $\mathcal P_{3m'}$ are the polynomials~\eqref{eq:P3m};
  \item $\mathcal E_{mpq}=\kb_2(\bar\sigma_{2mpq})\,e^{\I\Psi_{mpq}}$ combines the Love number of the term and its phase $\Psi_{mpq}=(2-2p)u_\M+q\M+f_{2m}$ \eqref{eq:Psi}--\eqref{eq:Ecal};
  \item the sums run over $m,p\in\{0,1,2\}$ and $|q|\le q_{\max}(e)$ (Section~\ref{sec:truncation}).
\end{itemize}
Properties:
\begin{itemize}
  \item It holds for any obliquity $0\le i\le\pi$. No separate formula is needed at $i=0$. A faster form for a coplanar orbit is given in Section~\ref{sec:coplanar-nbody}, as an optimisation only.
  \item It contains no division by $\sin i$, $\sin\theta^\star$ or $e$. The angular functions are polynomials in $(\rho,\bar\rho,c^\star)$.
  \item The exact geometry is used: $r^\star$ and the direction of the star come from the state vectors, not from the series. The truncation acts only on the Kaula expansion of the tidal-raising potential.
\end{itemize}
The double transformation~\eqref{eq:grad-node}, which uses the elements only, is given in Appendix~\ref{app:double} for reference.

\section{Tidal frequencies and Love numbers}\label{sec:love}
\subsection{Frequencies with mean rates}\label{sec:rates}
The Love number is defined for a harmonic forcing of fixed frequency (H1). Each Kaula term is quasi-harmonic, with a slowly varying rate (H3). The frequency used in $\kb_2$ is therefore the rate of $\nu_{2mpq}$ averaged over the \emph{fast} angles: the orbital period and the synodic periods with the other planets. The slow angles (secular precessions, resonant angles, super-periods) are not averaged; they are part of the forcing frequency. The phases $\Psi_{mpq}$ always use the osculating $u_\M$ and $\M$ of the current instant.

\paragraph{Regular form.} The rate~\eqref{eq:sigma} of Kaula's angle can be written in three equivalent forms, derived in Appendix~\ref{app:frequency}:
\begin{itemize}
  \item Kaula's form: $\sigma_{2mpq}=(2-2p+q)\,n-m\,\Op+(2-2p)\,\dot\omega+m\,\dot\Omega$;
  \item with the longitude of pericentre $\varpi=\Omega+\omega$: $\sigma_{2mpq}=(2-2p+q)\,n-m\,\Op+(2-2p)\,\dot\varpi+(m-2+2p)\,\dot\Omega$. In the coplanar case only $k=m-2+2p=0$ occurs and this is the usual form $(2-2p+q)n-m\Op+(2-2p)(\dot\Omega+\dot\omega)$;
  \item with the mean longitude $\lambda=\Omega+\omega+\M$, the form used here.
\end{itemize}
With $\lambda\equiv\Omega+\omega+\M$, the mean longitude of the star counted in the equator, and $k\equiv m-2+2p$, the rate at $l=2$ reads, identically,
\begin{equation}
  \boxed{\bar\sigma_{2mpq}=(2-2p)\,\dot{\bar\lambda}+q\,\bar n+k\,\dot{\bar\Omega}-m\,\Op,\qquad k=m-2+2p.} \label{eq:sigma-regular}
\end{equation}
where (a bar denotes the average over the fast angles, defined below):
\begin{itemize}
  \item $\dot{\bar\lambda}$ is the mean rate of the mean longitude $\lambda$ of the star; it is close to $n$ and contains the precession of the pericentre and of the node;
  \item $\bar n$ is the mean motion $\sqrt{\G(M^\star+M_{\mathrm p})/a^3}$ averaged over the window; it replaces $\dot\M$;
  \item $\dot{\bar\Omega}$ is the mean precession rate of the node of the star's orbit on the planet's equator, zero for an isolated Kepler orbit;
  \item $\Op=\|\mathbf s\|$ is the spin rate of the planet;
  \item $m,p,q$ are the indices of the term.
\end{itemize}
Properties:
\begin{itemize}
  \item $\dot{\bar\lambda}$ stays defined as $e\to0$ and as $i\to0$. The osculating $\dot\omega$ alone behaves like $1/e$, and $\dot\Omega$ like $1/\sin i$.
  \item $\dot{\bar\Omega}$ appears only with $k\neq0$. These terms have $F_{2mp}(i)=\mathcal O(i^{|k|})$ and vanish as $i\to0$, where $\Omega$ is undefined. For $m=2,p=0$ and $m=0,p=1$ ($k=0$, the only terms of the coplanar case), the frequency depends on $\dot{\bar\lambda}$, $\bar n$ and $\Op$ only.
\end{itemize}

\paragraph{Averaging window.} The window $\Delta T$ must remove the fast oscillations and keep the slow ones:
\begin{itemize}
  \item longer than the orbital period $P_i=2\pi/n_i$ of the planet, and than its synodic periods with the other bodies, $P_{{\rm syn},ij}=2\pi/|n_i-n_j|$, which are the periods of the conjunctions;
  \item much shorter than the slow time scales (secular precessions, libration of a resonance), which are part of the forcing.
\end{itemize}
It is computed by the code, at each accepted step, from the osculating mean motions $n_j$ of all the other orbiting bodies:
\begin{equation}
  \boxed{\Delta T_i=\min\Big(\max\big(\max_{j\neq i}P_{{\rm syn},ij},\ P_i\big),\ 10\,P_i\Big),\qquad P_{{\rm syn},ij}=\frac{2\pi}{|n_i-n_j|},\quad P_i=\frac{2\pi}{n_i}.} \label{eq:window}
\end{equation}
\begin{itemize}
  \item The longest synodic period covers all the conjunction frequencies of the planet.
  \item The lower bound $P_i$ applies to a single planet, the upper bound $10P_i$ to two planets with very close periods, whose synodic period would otherwise become longer than the slow time scales.
  \item The user gives no parameter.
\end{itemize}

\paragraph{Estimation in a run.} Keep, for each tidally active planet, a history of the accepted steps $(t_j,\lambda_j,\Omega_j,n_{{\rm osc},j})$ covering at least $\Delta T$:
\begin{enumerate}
  \item $\Omega_j$ is accumulated from the rotation of $\hat{\mathbf N}$ about $\hat{\mathbf s}$: $\Delta\Omega=\operatorname{atan2}\big(\hat{\mathbf s}\cdot(\hat{\mathbf N}_{j-1}\times\hat{\mathbf N}_j),\ \hat{\mathbf N}_{j-1}\cdot\hat{\mathbf N}_j\big)$;
  \item $\lambda_j=\lambda_{j-1}+\operatorname{wrap}\big(\Delta\Omega+u_{\M,j}-u_{\M,j-1}\big)$, with $\operatorname{wrap}$ to $(-\pi,\pi]$;
  \item $\dot{\bar\lambda}=[\lambda(t)-\lambda(t-\Delta T)]/\Delta T$ and $\dot{\bar\Omega}=[\Omega(t)-\Omega(t-\Delta T)]/\Delta T$, the values at $t-\Delta T$ being interpolated linearly between the two stored samples around it; $\bar n$ is the mean of $n_{\rm osc}$ over the window;
  \item before one window has elapsed: $\dot{\bar\lambda}=\bar n=n_{\rm osc}$, $\dot{\bar\Omega}=0$.
\end{enumerate}

\paragraph{Accuracy.}
\begin{itemize}
  \item Two planets near the 3:2 resonance: the rule~\eqref{eq:window} gives one synodic period. One integration step ($\Delta T=P/20$) gives the osculating rate, with an error of $4.5\times10^{-4}n$ on $\bar\sigma_{2200}$ and up to $45\,\%$ on the force near synchronism. One synodic period gives $9\times10^{-5}n$ and $4\,\%$.
  \item TRAPPIST-1 b, c, d, e, with their resonant chain: the rule gives $2.0$ to $3.0\,P_i$. The error on $\bar\sigma_{2200}$ is $2\times10^{-5}n$ to $6\times10^{-5}n$, against $5\times10^{-4}n$ to $8\times10^{-4}n$ for one step and $2\times10^{-4}n$ for one orbital period. Over the four planets, the worst error of the rule is within $10\,\%$ of the best of the fixed windows tested ($P_i$ to $10P_i$). No single window is the best for every planet.
  \item The history is updated on \emph{accepted} steps only. Intermediate force evaluations (sub-steps, predictor--corrector iterations) only read it, so that the rates do not depend on the trial states of the integrator.
\end{itemize}

\subsection{Love numbers from tabulated spectra}
The spectra $\texttt{rek2}(\chi)$ and $\texttt{imk2}(\chi)\ge0$ are given at physical frequencies $\chi=|\sigma|\ge0$, with $\texttt{imk2}>0$ for a lag. By~\eqref{eq:love-spectra},
\begin{equation}
  \kb_2(\bar\sigma)=\texttt{rek2}(|\bar\sigma|)-\I\,\sgn(\bar\sigma)\,\texttt{imk2}(|\bar\sigma|),\qquad \kb_2(0)=\texttt{rek2}(0). \label{eq:love-impl}
\end{equation}

\paragraph{Preparation of the table.} Keep the strictly increasing prefix $\chi_1<\dots<\chi_N$ of the tabulated frequencies: it stops at a padding that repeats the last row or that falls back to zero. Inside $[\chi_1,\chi_N]$, interpolate linearly.

\paragraph{Extrapolation above $\chi_N$.} Real rheologies have power-law tails at high frequency: $\texttt{imk2}\propto\chi^{-\alpha}$ and $\texttt{rek2}\to k_\infty+C\chi^{b}$, with $b=-\alpha$ (Andrade) or $b=-2\alpha$ (Maxwell, $\alpha=1$). With the last two (three) distinct points:
\begin{equation}
  \begin{gathered}
  \beta=\frac{\ln(\texttt{imk2}_N/\texttt{imk2}_{N-1})}{\ln(\chi_N/\chi_{N-1})}\ \text{clamped to}\ [-1,0],\qquad
  \texttt{imk2}(\chi)=\texttt{imk2}_N\Big(\frac{\chi}{\chi_N}\Big)^{\beta},\\
  \texttt{rek2}(\chi)=\texttt{rek2}_N+(\texttt{rek2}_N-\texttt{rek2}_{N-1})\,\frac{\chi^{b}-\chi_N^{b}}{\chi_N^{b}-\chi_{N-1}^{b}},
  \end{gathered} \label{eq:extrap}
\end{equation}
where $\chi_N>\chi_{N-1}>\chi_{N-2}$ are the last three distinct frequencies of the table, $\texttt{rek2}_j$ and $\texttt{imk2}_j$ the tabulated values there, $\beta$ the high-frequency slope of $\texttt{imk2}$ in log--log, and $b$ solves $\dfrac{\chi_N^b-\chi_{N-1}^b}{\chi_{N-1}^b-\chi_{N-2}^b}=\dfrac{\texttt{rek2}_N-\texttt{rek2}_{N-1}}{\texttt{rek2}_{N-1}-\texttt{rek2}_{N-2}}$ (a fit of $\texttt{rek2}=k_\infty+C\chi^b$ through the last three points), clamped between $2\beta$ and $\beta$, i.e.\ between the Maxwell and the Andrade behaviours.
\begin{itemize}
  \item If $\texttt{imk2}_N$ or $\texttt{imk2}_{N-1}$ is not positive, or $\beta\ge0$, both values are held constant. If $\texttt{rek2}$ is not monotonic over the last three points, $b=\beta$ is used.
  \item For $\beta<0$, $\chi^b\to0$: $\texttt{rek2}$ tends to a finite limit and $\texttt{imk2}$ to zero. The extrapolated values stay finite and $\texttt{imk2}\ge0$ for any $\chi$.
\end{itemize}

\paragraph{Extrapolation below $\chi_1$} (only if $\chi_1>0$). The same rule is applied with the first points, $\beta$ clamped to $[0,1]$ ($\texttt{imk2}\propto\chi^{\beta}$ below the dissipation peak), and $\texttt{imk2}(0)=0$.

\paragraph{Accuracy.} Tested on spectra truncated by one decade and compared with the exact values:
\begin{center}\small
\begin{tabular}{lcc}\toprule
Spectrum (cut one decade before its end) & error on \texttt{rek2} & error on \texttt{imk2}\\\midrule
Maxwell, homogeneous body & $3\times10^{-14}$ & $4\times10^{-9}$\\
Andrade, $\alpha=0.2$ / $0.3$ & $3\times10^{-5}$ / $1\times10^{-5}$ & $1.3\times10^{-3}$\\
Burgers & $6\times10^{-5}$ & $4\times10^{-2}$\\
TRAPPIST-1 b--h spectra (Revol et al.\ 2024) & $\le5\times10^{-3}$ & $\le8\times10^{-3}$\\
Low-frequency end, Maxwell / Andrade (cut below the peak) & $\le2\times10^{-7}$ & $\le5\times10^{-3}$\\
Holding the end values (no extrapolation), same cases & up to $7\times10^{-2}$ & $0.6$ to $9$\\\bottomrule
\end{tabular}
\end{center}
No rule can predict a feature that lies outside the data, such as a second relaxation peak beyond $\chi_N$: the spectra should cover the forcing frequencies of the run, with the extrapolation as a safeguard.

\section{Truncation and eccentricity functions}\label{sec:truncation}
The force must be accurate for any Love-number spectrum, since the spectrum can vary strongly with frequency. In the worst case, each term may have an arbitrary $\kb_{2,mpq}$ with $|\kb_{2,mpq}|\le\kappa_{\max}$. The force of a term is linear in $(\Re\kb,\Im\kb)$, $\vF_{mpq}=\mathsf L_{mpq}(\Re\kb,\Im\kb)^{\mathsf T}$, so the truncation error is bounded by
\[
  \|\delta\vF\|\le\kappa_{\max}\sum_{|q|>q_{\max}}\sum_{m,p}\|\mathsf L_{mpq}\|_2.
\]
The table below keeps $|q|\le q_{\max}(e)$. It guarantees, for any spectrum with $|\kb_2|\le\kappa_{\max}$,
\begin{equation}
  \|\delta\vF\|\le10^{-5}\,\kappa_{\max}F_0(r^\star),\qquad \|\delta\boldsymbol\Gamma\|\le10^{-5}\,\kappa_{\max}\,r^\star F_0(r^\star),\qquad F_0(r^\star)=\frac{3\G M^{\star2}R^5}{(r^\star)^7}, \label{eq:guarantee}
\end{equation}
at every point of the orbit, where $F_0(r^\star)$ is the instantaneous elastic force per unit Love number:
\begin{center}
\renewcommand{\arraystretch}{1.2}\setlength{\tabcolsep}{4pt}\small
\begin{tabular}{l|cccccccccccc}\toprule
$e\le$ & 0.001 & 0.01 & 0.02 & 0.05 & 0.10 & 0.15 & 0.20 & 0.25 & 0.30 & 0.35 & 0.40 & 0.50\\\midrule
$q_{\max}$ & 2 & 3 & 3 & 5 & 7 & 10 & 12 & 15 & 18 & 23 & 28 & 42\\\bottomrule
\end{tabular}
\end{center}
\begin{itemize}
  \item Each value is the worst over the points of a 35-point grid in $e$ that fall in the bin. For each $e$ there are 714 states: 7 inclinations from 0 to $\pi$, 3 arguments of pericentre, and 34 positions on the orbit (32 uniform in $\M$, plus apocentre and pericentre). The worst state is always near the apocentre, where the force is weakest. The bound was checked again on independent states.
  \item The torque bound, divided by $r^\star F_0$, never exceeds the force bound divided by $F_0$, since $\|\vr^\star\times\delta\vF\|\le r^\star\|\delta\vF\|$.
  \item The guarantee is relative to $\kappa_{\max}F_0$, a scale set mostly by the elastic part, while the secular evolution is driven by the dissipative part $\Im\kb_2$. The tolerance $10^{-5}$ resolves the dissipative part to $10^{-3}$ in the worst case when $|\Im\kb_2|/\kappa_{\max}\gtrsim10^{-2}$; in general, choose a tolerance of order $10^{-3}|\Im\kb_2|/\kappa_{\max}$. The tables for tolerances $10^{-2}$ to $10^{-6}$ are in Appendix~\ref{app:truncation}.
  \item The table covers $e\le0.5$. Beyond, $q_{\max}$ must be extended with the same method.
\end{itemize}

\paragraph{Eccentricity functions.} $G_{2pq}(e)$ is needed up to $|q|=42$. The power series truncated at $e^{7}$ is not sufficient: its error reaches $8\times10^{-4}$ at $e=0.2$ and $2\times10^{-2}$ at $e=0.3$. $G_{2pq}$ is computed exactly with Tisserand's formula (Appendix~\ref{app:kaula}, eq.~\eqref{eq:tisserand}), with at most 22 terms for $e\le0.5$ and an absolute error below $10^{-13}$. The symmetries $G_{22q}=G_{20,-q}$ and $G_{21q}=G_{21,-q}$ halve the work. The inclination functions $F_{2mp}(i)$ are the closed forms of Table~\ref{tab:F2mp}.

\paragraph{Fast evaluation of $G_{2pq}$.} Tisserand's formula costs up to $10^{-4}$\,s per force evaluation at $e=0.4$, and a cache keyed on $e$ does not help: the eccentricity rebuilt from the state vectors changes in its last bits at every call. For $e\le0.5$ the functions are therefore read from Chebyshev tables of their smooth part,
\begin{equation}
  G_{2pq}(e)=\beta_e^{|q|}\,H_{pq}(u),\qquad u=\beta_e^2,\qquad \beta_e=\frac{e}{1+\sqrt{1-e^2}},\qquad
  H_{pq}(u)\simeq\sum_{k=0}^{8}c^{(j)}_{pqk}\,T_k\big(t_j(u)\big), \label{eq:G-table}
\end{equation}
where:
\begin{itemize}
  \item $u\in[0,u_{\max}]$, $u_{\max}=\beta_e^2(0.5)=0.0718$, is cut into 8 equal intervals $j$; $t_j(u)\in[-1,1]$ is the local variable of interval $j$ and $T_k$ are the Chebyshev polynomials (Clenshaw's recurrence);
  \item $H_{pq}=G_{2pq}/\beta_e^{|q|}$ is a smooth function of $u$, since $G_{2pq}=\mathcal O(e^{|q|})$; the factor $\beta_e^{|q|}$ is applied afterwards;
  \item the coefficients $c^{(j)}_{pqk}$ are computed from Tisserand's formula (45 terms) at the 9 Chebyshev nodes of each interval, once per run (about 0.1\,s).
\end{itemize}
The absolute error is $3\times10^{-15}$ for $e\le0.5$ and $|q|\le q_{\max}(e)$. Above $e=0.5$, Tisserand's formula is used directly.

\section{Force and torque in inertial axes}\label{sec:torque}
\paragraph{Force.} From~\eqref{eq:impl-single} and the convention~\eqref{eq:force-convention},
\begin{equation}
  \boxed{\Fsp=M^\star\Big[\Re(V_X+\I V_Y)\,\hat{\mathbf N}+\Im(V_X+\I V_Y)\,\hat{\mathbf Y}+V_Z\,\hat{\mathbf s}\Big].} \label{eq:impl-force}
\end{equation}
where $\Fsp$ is the force of the star on the planet in inertial axes, $M^\star$ the mass of the star, $V_X,V_Y,V_Z$ the components~\eqref{eq:impl-single} of $\nabla T$ at the star, and $(\hat{\mathbf N},\hat{\mathbf Y},\hat{\mathbf s})$ the nodal frame expressed in inertial axes~\eqref{eq:node-frame}.

\paragraph{Torque.} Neglect the tide in the star. The only internal interaction of the pair is $\Fsp=-\Fps$, applied at the centres of mass. With the reduced mass $\mu=M^\star M_{\mathrm p}/(M^\star+M_{\mathrm p})$, the orbital angular momentum $\mathbf L_{\rm orb}=\mu\,\vr\times\dot\vr$ evolves as $\dot{\mathbf L}_{\rm orb}=\vr\times\Fsp$, because gravity is central. The conservation of $\mathbf L_{\rm orb}+\mathbf L_{\rm spin}$ then imposes
\begin{equation}
  \boxed{\frac{\dd(I\mathbf s)}{\dd t}=\boldsymbol\Gamma=-\,\vr\times\Fsp=M^\star\,\vr^\star\times\nabla T\big|_{\vr^\star}.} \label{eq:torque}
\end{equation}
where $I$ is the moment of inertia of the planet, $\mathbf s$ its spin vector, $\vr=\vr_{\mathrm p}-\vr_\star$ the position of the planet relative to the star, $\vr^\star=-\vr$ the position of the star seen from the planet, and $\Fsp$ the force~\eqref{eq:impl-force}.
\begin{itemize}
  \item The torque must be computed with \emph{the same instantaneous} force as the one applied to the orbit. A torque built from an orbit-averaged force exchanges, at each step, a different angular momentum from the one taken from the orbit, and $\mathbf L_{\rm tot}$ drifts.
  \item The component along the spin axis, $\boldsymbol\Gamma\cdot\hat{\mathbf s}$, changes the spin rate; the perpendicular components change the direction of the spin axis.
  \item The constant-time-lag model \citep{Hut1981,Bolmont2015} satisfies~\eqref{eq:torque} exactly.
\end{itemize}

\section{Coplanar orbit: a faster form}\label{sec:coplanar-nbody}
When the orbit lies in the equator of the planet ($i=0$), only $(m,p)=(0,1)$ and $(2,0)$ survive in~\eqref{eq:impl-single}, since $F_{201}(0)=-\tfrac12$, $F_{220}(0)=3$ and the other $F_{2mp}(0)$ vanish. The star is in the equator: $c^\star=0$, $|\rho|=1$, $\mathcal P_{31}(0)=-\tfrac32$, $\mathcal P_{33}(0)=15$ and $\mathcal P_{30}(0)=\mathcal P_{32}(0)=0$. Hence $V_Z=0$, and $V_X+\I V_Y=\rho\,(w_r+\I w_\varphi)$ with real $w_r$, $w_\varphi$. The gradient has a radial and an azimuthal part only:
\begin{equation}
\boxed{\begin{gathered}
  \nabla T\big|_{\vr^\star}=\frac{\G M^\star R^5}{a^3(r^\star)^4}\big(w_r\,\hat\vr^\star+w_\varphi\,\hat{\boldsymbol\varphi}\big),\qquad \hat{\boldsymbol\varphi}=\hat{\mathbf s}\times\hat\vr^\star,\\
  w_r=\sum_{|q|\le q_{\max}}\Big[\tfrac34\,G_{21q}(e)\,\Re\big(\kb_{01q}\,e^{\I q\M}\big)+\tfrac94\,G_{20q}(e)\,\Re Z_q\Big],\\
  w_\varphi=-\tfrac32\sum_{|q|\le q_{\max}}G_{20q}(e)\,\Im Z_q,\qquad Z_q=\kb_{20q}\,e^{\I[(2+q)\M-2f]},
\end{gathered}} \label{eq:coplanar-nbody}
\end{equation}
where:
\begin{itemize}
  \item $\kb_{01q}=\kb_2(q\bar n)$ and $\kb_{20q}=\kb_2(2\dot{\bar\lambda}+q\bar n-2\Op)$ are the Love numbers at the frequencies~\eqref{eq:sigma-regular} with $i=0$;
  \item $f$ is the true anomaly of the star: $\cos f=\hat{\mathbf e}_{\rm v}\cdot\hat\vr^\star$, $\sin f=\hat{\mathbf h}\cdot(\hat{\mathbf e}_{\rm v}\times\hat\vr^\star)$, with $\hat{\mathbf e}_{\rm v}=\mathbf e_{\rm v}/e$; if $e=0$, $f=\M=0$;
  \item only $a$, $e$, $r^\star$, $\hat\vr^\star$, $\M$ and $f$ are needed: no node, no $\omega$, no $F_{2mp}$ and no $\mathcal P_{3m}$.
\end{itemize}
The force and the torque follow from~\eqref{eq:force-convention} and~\eqref{eq:torque}: $\Fsp=M^\star\nabla T|_{\vr^\star}$. This is the form~\eqref{eq:coplanar-l2} of Chapter~\ref{ch:coplanar} written with $|\kb|\cos(\cdot-\epsilon)=\Re(\kb\,e^{\I\,\cdot})$; it equals~\eqref{eq:impl-single} at $i=0$ to rounding error.

\paragraph{When to use it.} It costs $2(2q_{\max}+1)$ terms instead of $9(2q_{\max}+1)$. Its error with respect to~\eqref{eq:impl-single} is linear in the obliquity. For any Love numbers with $|\kb_2|\le\kappa_{\max}$ (the assumption of the truncation guarantee~\eqref{eq:guarantee}),
\begin{equation}
  \|\vF_{\rm coplanar}-\vF_{\rm single}\|\le C(e)\;i\;\kappa_{\max}F_0(r^\star), \label{eq:coplanar-switch}
\end{equation}
where $C(e)$ is the worst case over the orbit, reached at apocentre: $C=2.0$ ($e=0$), $2.7$ ($e=0.05$), $3.6$ ($0.1$), $7.1$ ($0.2$), $13$ ($0.3$), $24$ ($0.4$), $55$ ($0.5$). Realistic rheologies stay well below (Maxwell, constant time lag: $C\le2.5$ at $e=0.49$). The coplanar form is used for $i\le i_s=10^{-7}$\,rad: the jump at the switch is then below $5.5\times10^{-6}\,\kappa_{\max}F_0$ for all $e\le0.5$, i.e.\ below the truncation tolerance $10^{-5}$. It applies to a prograde orbit only ($\hat{\mathbf h}\cdot\hat{\mathbf s}>0$); near $i=\pi$ the surviving terms are $(0,1)$ and $(2,2)$, and the general form~\eqref{eq:impl-single} is used.

\section{Dissipated power}\label{sec:power}
$N$-body codes often output the power dissipated by the tide in the planet. For the constant-time-lag model it is $-(\mathbf F_{\rm diss}\cdot\mathbf v+\boldsymbol\Gamma\cdot\mathbf s)$, with the dissipative part of the force only \citep{Bolmont2015}. The same definition is used here: the dissipative part is the part of the force carried by $\Im\kb_2$,
\begin{equation}
  \boxed{P_{\rm diss}=-\big(\mathbf F_{\Im}\cdot\mathbf v+\boldsymbol\Gamma_{\Im}\cdot\mathbf s\big),} \label{eq:power}
\end{equation}
where $\mathbf v=\mathbf v_{\mathrm p}-\mathbf v_\star$, and $\mathbf F_{\Im}$, $\boldsymbol\Gamma_{\Im}$ are the force~\eqref{eq:impl-force} and the torque~\eqref{eq:torque} computed with every $\kb_2(\bar\sigma_{2mpq})$ replaced by $\I\,\Im\kb_2(\bar\sigma_{2mpq})$. Since the force is linear in $(\Re\kb_2,\Im\kb_2)$, $\mathbf F_{\Im}$ and $\boldsymbol\Gamma_{\Im}$ are obtained by a second accumulation of the same terms.
\begin{itemize}
  \item \textbf{Orbit average.} Over an orbit with constant rates,
  \begin{equation}
    \langle P_{\rm diss}\rangle=\frac{\G M^{\star2}R^5}{a^6}\sum_{m,p,q}c_m\,F_{2mp}^2(i)\,G_{2pq}^2(e)\;\bar\sigma_{2mpq}\,\big(-\Im\kb_2(\bar\sigma_{2mpq})\big)\ \ge0, \label{eq:power-mean}
  \end{equation}
  non-negative since $\sigma\,\Im\kb_2(\sigma)\le0$ for a lag.
  \item \textbf{Instantaneous value.} $P_{\rm diss}$ is not positive at every instant for a general rheology (minimum $-0.8\,\%$ of the mean in an Andrade case at $e=0.3$, against $-860\,\%$ for the power of the total force). The split into $\Re$ and $\Im$ parts is a convention, the one that extends the constant-time-lag output: in that limit~\eqref{eq:power} equals the constant-time-lag power up to $\mathcal O\big((n\Delta t)^2\big)$. For an elastic planet it is exactly zero.
\end{itemize}

\section{Algorithm}\label{sec:algorithm}
For each planet with the tide active, at each force evaluation:
\begin{enumerate}
  \item \textbf{Geometry} (Section~\ref{sec:geometry}): $\hat{\mathbf s}$, $\Op$, $\hat{\mathbf h}$; nodal frame $\hat{\mathbf N},\hat{\mathbf Y}$; $i$, $\hat{\mathbf W}$, $\hat{\mathbf N}'$; $a$, $e$, $\bar n_{\rm osc}$, $u$, $\omega$, $f$, $\M$, $u_\M$; $r^\star$, $\rho$, $c^\star$.
  \item \textbf{Mean rates} (Section~\ref{sec:rates}): $\bar n$, $\dot{\bar\lambda}$, $\dot{\bar\Omega}$ from the history (read only).
  \item \textbf{Kaula functions}: $q_{\max}(e)$ from the table of Section~\ref{sec:truncation}; $F_{2mp}(i)$ for $m,p=0,1,2$; $G_{2pq}(e)$ for $p=0,1$ and $|q|\le q_{\max}$, from~\eqref{eq:G-table} for $e\le0.5$.
  \item \textbf{Terms}: for each term $(m,p,q)$ with $F_{2mp}(i)\neq0$, compute $\bar\sigma_{2mpq}$~\eqref{eq:sigma-regular}, $\kb_2(\bar\sigma)$~\eqref{eq:love-impl}, $\mathcal K_{2mpq}$~\eqref{eq:K2}, $\mathcal E_{mpq}$~\eqref{eq:Ecal}; accumulate the single form~\eqref{eq:impl-single}, or the coplanar form~\eqref{eq:coplanar-nbody} if $i\le i_s$.
  \item \textbf{Output}: $\Fsp$~\eqref{eq:impl-force} and $\boldsymbol\Gamma$~\eqref{eq:torque}, and their dissipative parts for~\eqref{eq:power}.
\end{enumerate}
At each \emph{accepted} step, update $\lambda$ and $\Omega$, recompute the window~\eqref{eq:window}, and append $(t,\lambda,\Omega,n_{\rm osc})$ to the history.

\paragraph{Exact savings.} Two identities halve the cost of the Love numbers and of the phases, without changing the result beyond rounding:
\begin{itemize}
  \item for $m=0$, $\bar\sigma_{201,-q}=-\bar\sigma_{201q}$ and $\bar\sigma_{202q}=-\bar\sigma_{200,-q}$ exactly, and $\kb_2(-\sigma)=\kb_2(\sigma)^*$ (hypothesis H1): one evaluation of the spectrum per pair;
  \item $e^{\I q\M}$ is obtained by successive products of $e^{\I\M}$ (rounding error $\approx|q|\times10^{-16}$).
\end{itemize}

\paragraph{Cost.} At most $9\times(2q_{\max}+1)$ terms: 45 for $e\le0.001$, 63 for $e\le0.02$ and 765 for $e=0.5$. In the coplanar form, $2\times(2q_{\max}+1)$. Each term needs one interpolation in the spectra and one complex product.

\section{Validation tests}\label{sec:validation}
Any implementation should pass the following tests. Each one checks a property of the exact physics, independent of the Kaula derivation. The reference implementation (Python) and Posidonius pass all of them (Section~\ref{sec:pos-tests}).
\begin{description}
  \item[V1 Circular equatorial orbit, constant time lag.] The closed forms of Section~\ref{sec:physical-check}:
  \[
    F^{\mathrm p\to\star}_r=-\frac{3k_2\G M^{\star2}R^5}{a^7}\Big[\tfrac14+\tfrac34\cos2(\Op-n)\Delta t\Big],\qquad
    F^{\mathrm p\to\star}_\varphi=\frac{3k_2\G M^{\star2}R^5}{2a^7}\sin2(\Op-n)\Delta t,
  \]
  and $F_\theta=0$.
  Agreement $5\times10^{-16}$.
  \item[V2 Constant time lag, eccentric and oblique.] With $\kb_2=k_2e^{-\I\sigma\Delta t}$ and Kepler rates, the force equals the exact force of the delayed bulge: the bulge points to the star position $\Delta t$ earlier, rotated with the planet by $\Op\Delta t$. Agreement $10^{-10}$ with $|q|\le40$.
  \item[V3 Small time lag.] The difference with the Hut force \citep{Hut1981} decreases as $\Delta t$.
  \item[V4 Elastic planet.] With a constant real $k_2$, for any rates and any obliquity, $\boldsymbol\Gamma=0$ and
  \[
    \Fsp=-3k_2\G M^{\star2}R^5(r^\star)^{-7}\,\hat{\vr}.
  \]
  \item[V5 Rotation of the inertial axes.] Rotating $(\vr,\mathbf v,\mathbf s)$ rotates $\Fsp$ and $\boldsymbol\Gamma$ in the same way. This test detects any hidden reference to the inertial axes or to an inertial plane.
  \item[V6 Change of length unit.] Multiply all lengths by $\lambda$ and $\G$ by $\lambda^3$, which keeps the time unit. Then $\Fsp\to\lambda\Fsp$ and $\boldsymbol\Gamma\to\lambda^2\boldsymbol\Gamma$. This test detects dimensional errors.
  \item[V7 Obliquity $\to0$.] The force is continuous at $i=0$, with $|\vF(i)-\vF(0)|=\mathcal O(i)$, and independent of the node direction chosen at $i=0$.
  \item[V8 Signs.] For a coplanar orbit and a constant time lag, the orbit-averaged $\boldsymbol\Gamma\cdot\hat{\mathbf s}$ is negative for $\Op>n$ and positive for $\Op<n$. For an oblique spin, the averaged torque equals the averaged Hut torque to $\mathcal O(\Delta t)$.
  \item[V9 Angular momentum.] In a two-body integration with the instantaneous torque~\eqref{eq:torque}, $\mathbf L_{\rm orb}+I\mathbf s$ is conserved up to the integration error of the scheme.
\end{description}

\section{Assumptions and limits}\label{sec:limits}
\begin{itemize}
  \item \textbf{Physics:} hypotheses H1--H3, response at the mean frequency. The result is exact for a constant time lag. For a general rheology, the error is of order $\partial_\sigma\kb_2\cdot\delta\sigma$, where $\delta\sigma$ is the neglected fast modulation of the rates.
  \item \textbf{Elements:} two-body osculating elements with $\mu_{\rm orb}=\G(M^\star+M_{\mathrm p})$, referred to the instantaneous equator.
  \item \textbf{Degree:} degree two only. The neglected degree-three terms are of relative order $(R/a)^2$.
  \item \textbf{Scope:} the tide in the planet. The tide in the star uses the same equations with the roles of the bodies exchanged (Chapter~\ref{ch:star}).
  \item \textbf{Eccentricity:} $e\le0.5$ with the truncation table above. Beyond, the code falls back to the largest $q_{\max}$ of the table ($42$) and Tisserand's formula, without guarantee.
  \item \textbf{Spin--orbit synchronism.} Hypothesis H3 fails for a term whose mean frequency is zero. At the 1:1 spin--orbit resonance, $\bar\sigma_{2200}=2(\dot{\bar\lambda}-\Op)$ uses the instantaneous spin rate $\Op$, while the spin librates with the orbital period: $\bar\sigma_{2200}$ changes sign twice per orbit. Its Love number is read at a frequency that oscillates around zero: $\Im\kb_2$ changes sign with $\bar\sigma_{2200}$, and jumps if the spectrum does not tend to zero at low frequency. The torque of this term then behaves like a dry friction (stick--slip): the spin equation becomes stiff and discontinuous, and the secular rates obtained in a run depend on the integration step. With an obliquity, the terms $(1,0,-1)$, $(1,1,1)$ and $(2,1,2)$ also have frequencies close to zero. A treatment of the response in the time domain (the bulge as a state variable with its own relaxation equations) would be needed in this regime; it is not part of this document.
\end{itemize}

\chapter{The tide raised in the star}\label{ch:star}

The previous chapters treat the tide raised by the star in the planet. A planet also raises a tide in its star, which acts back on the orbit and on the spin of the star. This chapter shows that the same equations apply with the roles of the two bodies exchanged, and states what changes: the spectrum of the star depends on the sign of the frequency.

\section{Exchange of the roles}\label{sec:star-roles}
Nothing in Chapters~\ref{ch:potential}--\ref{ch:nbody} uses the fact that the deformed body is a planet. Call D the deformed body and P the point-mass perturber that raises the tide. The derivation gives the force of P on D, and the torque on the spin of D, from the state of P relative to D, the spin of D, the radius and the Love number of D, and the mass of P. For the tide raised in the star:
\begin{center}\small
\renewcommand{\arraystretch}{1.25}
\begin{tabular}{p{5.6cm}p{4.2cm}p{4.2cm}}\toprule
Quantity & Tide in the planet (D $=$ planet) & Tide in the star (D $=$ star)\\\midrule
relative position $\vr=\vr_{\rm D}-\vr_{\rm P}$ & $\vr_{\mathrm p}-\vr_\star$ & $\vr_\star-\vr_{\mathrm p}$\\
relative velocity $\mathbf v$ & $\mathbf v_{\mathrm p}-\mathbf v_\star$ & $\mathbf v_\star-\mathbf v_{\mathrm p}$\\
spin $\mathbf s$ & spin of the planet & spin of the star $\boldsymbol\Omega_\star$\\
radius $R$ & $R_{\mathrm p}$ & $R_\star$\\
mass of the perturber, $M^\star$ in the formulas & $M_\star$ & $M_{\mathrm p}$\\
Love number & $\kb_2$ of the planet & $\kb_\star$ of the star (Section~\ref{sec:star-spectrum})\\
force~\eqref{eq:impl-force} & $\Fsp$, on the planet & $\mathbf F_{{\rm p}\to\star}$, on the star\\
torque~\eqref{eq:torque} & on the spin of the planet & on the spin of the star\\\bottomrule
\end{tabular}
\end{center}
\begin{itemize}
  \item The geometry of Section~\ref{sec:geometry} is computed with these vectors: the inclination $i$ is that of the orbit on the equator of the star, and the elements are those of the planet seen from the star.
  \item The force~\eqref{eq:impl-force} is then the force of the planet on the star; the planet receives the opposite force. The torque~\eqref{eq:torque}, $\boldsymbol\Gamma_\star=-\vr\times\mathbf F_{{\rm p}\to\star}$, acts on the spin of the star, and $\dot{\mathbf L}_{\rm orb}+\boldsymbol\Gamma_\star=0$ exactly.
  \item With several planets, each planet raises its own tide in the star. The force and the torque are summed over the planets; the cross terms (the bulge raised by one planet acting on another planet) are neglected.
  \item The coplanar form (Section~\ref{sec:coplanar-nbody}) applies when the orbit lies in the equator of the star; the dissipated power~\eqref{eq:power} is the power dissipated in the star.
\end{itemize}
This exchange was checked against the exact linear response of a rotating deformed body to a point mass, computed without Kaula's expansion: the force agrees to $4.5\times10^{-15}$ and the angular momentum is conserved exactly.

\section{Love number of the star}\label{sec:star-spectrum}
\paragraph{Why the sign of the frequency matters.} For a planet, the response depends only on $|\sigma|$, and $\kb_2(-\sigma)=\kb_2(\sigma)^*$ (hypothesis H1). The tidal response of a rotating star, with its inertial waves, depends on whether the tidal pattern rotates in the same direction as the star or in the opposite direction. The stellar spectra are therefore tabulated for both signs of the frequency, and they are not symmetric: $\kb_\star(-\sigma)\neq\kb_\star(\sigma)^*$.

\paragraph{Convention.} A term $(m,p,q)$ of the expansion is a pattern $e^{\I(\sigma t-m\varphi)}$ in the frame of the deformed body, with frequency $\bar\sigma_{2mpq}$~\eqref{eq:sigma-regular}. $\bar\sigma>0$ is a pattern prograde with respect to the rotation of the body. The term takes the Love number at its own signed frequency. This was checked against the exact response of a body whose response depends on the order $m$ and on the sign of the frequency: the reading $\kb_\star(\bar\sigma_{2mpq})$ is exact to $1.9\times10^{-15}$, while the reading with the opposite sign is wrong by $7\,\%$.

\paragraph{Reading of the tabulated spectrum.} The stellar spectrum gives, at signed frequencies $x$ (rad/s), the values $\texttt{rek2}(x)$ and $\texttt{imk2}(x)$, with $\texttt{imk2}$ of the sign of $x$ (a lag for both signs). It is read as
\begin{equation}
  \boxed{\kb_\star(\bar\sigma)=\texttt{rek2}(x)-\I\,\texttt{imk2}(x)\Big(\frac{\Omega_\star}{\Omega_{\rm ref}}\Big)^2,\qquad
  x=\frac{\bar\sigma}{\mathrm{DAY}}\,\frac{\Omega_{\rm ref}}{\Omega_\star},} \label{eq:kbar-star}
\end{equation}
where:
\begin{itemize}
  \item $\bar\sigma$ is the signed frequency of the term, in rad/day, and DAY the number of seconds in a day;
  \item $\Omega_\star=\|\boldsymbol\Omega_\star\|$ is the spin rate of the star, and $\Omega_{\rm ref}$ the spin rate for which the spectrum was computed. The spectrum is rescaled to the current spin: frequencies in proportion to $\Omega_\star$, $\texttt{imk2}$ as $\Omega_\star^2$. Without $\Omega_{\rm ref}$, there is no rescaling ($\Omega_{\rm ref}/\Omega_\star=1$);
  \item $\Im\kb_\star<0$ for $\bar\sigma>0$, as for the planet~\eqref{eq:love-impl}. For a spectrum symmetric in $x$, \eqref{eq:kbar-star} reduces to~\eqref{eq:love-impl}.
\end{itemize}
The frequency of the files is taken to be the tidal frequency in the frame of the star, with the sign convention above. If the supplier of the spectra used the opposite sign convention, the dissipative part of the tide would change (the elastic part would not, since $\texttt{rek2}$ is nearly constant in these files); this convention should be confirmed with the supplier.

\paragraph{Table.}
\begin{itemize}
  \item The table is the strictly increasing prefix $x_1<\dots<x_N$ of the tabulated frequencies, with $x_1<0<x_N$ and $\texttt{imk2}\,x\ge0$ at every row; the code stops on an invalid table.
  \item Linear interpolation inside $[x_1,x_N]$, also across $x=0$: a non-zero $\texttt{imk2}(0)$ given by the table is kept.
  \item Outside the table, the extrapolation~\eqref{eq:extrap} is applied in $|x|$ at each end, and the sign of $\texttt{imk2}$ restored. A warning is printed once per run, since the accuracy is then not guaranteed.
\end{itemize}

\paragraph{Orders $m=0,1,2$.} The same spectrum is used for every order $m$, as hypothesis H2 does for the planet. For $m=0$ the identity $\kb(-\sigma)=\kb(\sigma)^*$ no longer holds, so each term is evaluated separately (the saving of Section~\ref{sec:algorithm} is not used). The single form~\eqref{eq:impl-single} uses $\Re\,\mathcal E_{0pq}$ for $m=0$: it keeps the Hermitian part $\tfrac12\big[\kb_\star(\sigma)+\kb_\star(-\sigma)^*\big]$ of the input, which is also what the exact response does for an axisymmetric pattern. The choice of the spectrum for $m=0$ and $m=1$ is a physical assumption of the model.

\section{Mean rates}
The frequencies of the tide in the star use the same estimator as for the planet (Section~\ref{sec:rates}), applied to the orbit seen in the equator of the star. One history is kept for each pair (star, planet), updated at the accepted steps, with the averaging window~\eqref{eq:window} of the planet.

\section{Checks}
\begin{itemize}
  \item 16 random states with a tabulated stellar spectrum: the implementation equals the reference single form to $2.7\times10^{-15}$, and the exact response at $|q|\le42$ to $4.2\times10^{-15}$.
  \item The torque on the star is the sum over the planets of the torques~\eqref{eq:torque}, exactly.
  \item Two-body integrations: $\mathbf L_{\rm orb}+\mathbf L_\star$ conserved to $8\times10^{-17}$ of the exchanged angular momentum; the secular rates of the energy and of the spin of the star equal their orbit averages to $5\times10^{-5}$.
\end{itemize}

\chapter{Implementation in Posidonius}\label{ch:posidonius}

This chapter describes how the equations of Chapters~\ref{ch:nbody} and~\ref{ch:star} are implemented in the $N$-body code Posidonius (Rust). It gives the files and functions, the units, the data stored with the bodies, the messages printed by the code, the tests, and the known limitations. The model is selected per body with the tidal model \script{Kaula} of the Python interface; the input format of the spectra is unchanged.

\section{Files and functions}\label{sec:pos-files}
The equations are written once, for a deformed body D and a perturber P (Section~\ref{sec:star-roles}), and used for both tides.
\begin{center}\small
\renewcommand{\arraystretch}{1.2}
\begin{tabular}{p{6.7cm}p{7.8cm}}\toprule
File & Content\\\midrule
\script{src/effects/tides/kaula/kaula_tide.rs} & the equations: geometry, mean rates, Love numbers, Kaula functions, truncation, single and coplanar forms, force, torque and dissipative parts. Pure functions of the state, independent of \script{Particle}\\
\script{src/effects/tides/kaula.rs} & the connection to Posidonius: planetary tide (D $=$ body, P $=$ tidal host) and stellar tide (D $=$ tidal host, P $=$ body), stored states, update of the rate histories\\
\script{src/effects/tides/kaula/love_number.rs} & storage of the tabulated spectrum (\script{LoveNumber}) and read access to it\\
\script{src/effects/tides.rs} & calls of the planetary and stellar tides, dissipated power (\script{calculate_denergy_dt})\\
\script{src/particles/universe.rs} & \script{update_kaula_rate_histories}, called by the integrators\\
\script{src/integrator/ias15.rs}, \script{whfast.rs}, \script{leapfrog.rs} & one call per accepted step to \script{update_kaula_rate_histories}\\\bottomrule
\end{tabular}
\end{center}
The former Kaula code (\script{components_2d.rs}, \script{components_3d.rs}, \script{polynomials.rs} and the Love-number cache) is removed.

\paragraph{Equations and functions.} The functions of \script{kaula_tide.rs} implement:
\begin{center}\small
\renewcommand{\arraystretch}{1.2}
\begin{tabular}{p{7cm}p{7.5cm}}\toprule
Function & Equations\\\midrule
\script{nodal_geometry} & \eqref{eq:elements-a-e}--\eqref{eq:rho}: $a$, $e$, $n_{\rm osc}$, nodal frame, $i$, $\omega$, $\M$, $f$, $u_\M$, $r^\star$, $\rho$, $c^\star$\\
\script{update_rate_history}\newline\script{mean_rates} & estimator of Section~\ref{sec:rates}: accumulation of $\lambda$ and $\Omega$, ring buffer, $\bar n$, $\dot{\bar\lambda}$, $\dot{\bar\Omega}$\\
\script{mean_motion_about_host}\newline\script{averaging_window} & window~\eqref{eq:window}\\
\script{frequency} & $\bar\sigma_{2mpq}$~\eqref{eq:sigma-regular}\\
\script{prepare_spectrum}\newline\script{evaluate_spectrum}\newline\script{love_number_at} & unsigned spectrum of a planet: table, interpolation, extrapolation~\eqref{eq:extrap}, $\kb_2$~\eqref{eq:love-impl}\\
\script{prepare_signed_spectrum}\newline\script{evaluate_signed_spectrum}\newline\script{love_number_signed} & signed spectrum of a star, $\kb_\star$~\eqref{eq:kbar-star}\\
\script{q_max} & truncation table of Section~\ref{sec:truncation}\\
\script{inclination_functions_f2} & $F_{2mp}(i)$, Table~\ref{tab:F2mp}\\
\script{eccentricity_functions_g2}\newline\script{eccentricity_functions_g2_tisserand} & $G_{2pq}(e)$: tables~\eqref{eq:G-table} for $e\le0.5$, Tisserand's formula~\eqref{eq:tisserand} otherwise\\
\script{tidal_gradient_nodal}\newline\script{tide_single} & single form~\eqref{eq:K2}--\eqref{eq:impl-single}, force~\eqref{eq:impl-force}, torque~\eqref{eq:torque}\\
\script{tidal_gradient_coplanar}\newline\script{tide_coplanar} & coplanar form~\eqref{eq:coplanar-nbody}\\
\script{tide}\newline\script{tide_tabulated} & choice of the form ($i\le i_s=10^{-7}$: coplanar), Love numbers from the tabulated spectrum; returns the force, the torque and their dissipative parts~\eqref{eq:power}\\\bottomrule
\end{tabular}
\end{center}
The functions of the forms take the Love number as a closure, so that the tests can use analytic rheologies (Maxwell, constant time lag, elastic) as well as tabulated spectra.

\section{Call chain}\label{sec:pos-flow}
At each force evaluation, every integrator calls \script{Universe::calculate_additional_effects}, which calls \script{tides::calculate_tidal_acceleration} and \script{tides::calculate_dangular_momentum_dt_due_to_tides}.
\begin{itemize}
  \item \textbf{Planetary tide.} For each body with \script{TidesEffect::OrbitingBody(TidalModel::Kaula(_))}, \script{kaula::calculate_planetary_tidal_force} returns $\Fsp$. \script{tides.rs} applies $+\Fsp/M_{\mathrm p}$ to the body and $-\Fsp/M_\star$ to the tidal host (Section~\ref{sec:io}).
  \item \textbf{Stellar tide.} If the tidal host has \script{TidesEffect::CentralBody(TidalModel::Kaula(_))}, \script{kaula::calculate_stellar_tidal_force} returns, for each orbiting body, the force of the body on the host, $\mathbf F_{{\rm p}\to\star}$. The host receives $+\mathbf F_{{\rm p}\to\star}/M_\star$ and the body the opposite. Every \script{OrbitingBody} feels the stellar tide, whatever its own tidal model (\script{DisabledModel} included); bodies with \script{TidesEffect::Disabled} have no tidal coordinates and are skipped.
  \item \textbf{Torques.} \script{kaula::calculate_torque_due_to_tides} returns $\vr\times\vF=-\boldsymbol\Gamma$, which \script{tides.rs} negates and stores in \script{dangular_momentum_dt}; the integrators integrate the angular momentum of the spin with it.
\end{itemize}

\paragraph{Same state for the force and the torque.} Each computed tide is stored with the state $(\vr,\mathbf v,\mathbf s)$ that produced it (\script{TideState}). The torque reuses it when the state is the same, and recomputes it otherwise: LeapFrog, for instance, asks for the torque without the force. The force and the torque therefore always come from the same instantaneous state, which is the condition for the conservation of the total angular momentum (Section~\ref{sec:torque}).

\section{Units and variables}\label{sec:pos-units}
Posidonius uses the units of the Mercury code (\script{constants.rs}):
\begin{center}\small
\begin{tabular}{lll}\toprule
Quantity & Unit & Constant\\\midrule
length & AU & \script{AU} $=1.495\,978\,707\times10^{11}$ m\\
time & day & \script{DAY} $=86\,400$ s\\
mass & $M_\odot$ & \script{M_SUN} $=1.9818\times10^{30}$ kg\\
$\G$ & AU$^3\,M_\odot^{-1}$\,day$^{-2}$ & \script{G} $=\G_{\rm SI}M_\odot\,\mathrm{DAY}^2/\mathrm{AU}^3=2.9493\times10^{-4}$\\
spin, $n$, $\dot{\bar\lambda}$, $\dot{\bar\Omega}$, $\bar\sigma$ & rad/day & \\
force & $M_\odot$\,AU\,day$^{-2}$ & \\
torque & $M_\odot$\,AU$^2$\,day$^{-2}$ & \\
dissipated power & $M_\odot$\,AU$^2$\,day$^{-3}$ & \\\bottomrule
\end{tabular}
\end{center}
The spectra are tabulated in rad/s. The only conversion is at the call of the spectrum: $\chi\,[\mathrm{rad/s}]=|\bar\sigma|\,[\mathrm{rad/day}]/\mathrm{DAY}$, and $x=\bar\sigma/\mathrm{DAY}$ (rescaled by the spin, \eqref{eq:kbar-star}) for a star.

\paragraph{Correspondence of the variables} (planetary tide; for the stellar tide, exchange the roles as in Section~\ref{sec:star-roles}).
\begin{center}\small
\renewcommand{\arraystretch}{1.2}
\begin{tabular}{>{$}l<{$}p{7.4cm}p{5cm}}\toprule
\text{Symbol} & Posidonius & Note\\\midrule
\vr=\vr_{\mathrm p}-\vr_\star & \script{particle.tides.coordinates.position} & relative to the tidal host\\
\mathbf v=\mathbf v_{\mathrm p}-\mathbf v_\star & \script{particle.tides.coordinates.velocity} & idem\\
\mathbf s & \script{particle.spin} & rad/day, inertial axes\\
M_{\mathrm p},\ R & \script{particle.mass}, \script{particle.radius} & $M_\odot$, AU\\
M_\star & \script{tidal_host_particle.mass} & $M_\odot$\\
\texttt{rek2},\ \texttt{imk2},\ \chi & \script{KaulaParameters::love_numbers} & three arrays of 1024 values\\
\Omega_{\rm ref} & \script{stellar_spectrum_spin_rate} & rad/s, stellar spectra only\\\bottomrule
\end{tabular}
\end{center}
All quantities, the orbital elements included, are computed from \script{tides.coordinates}, i.e.\ relative to the tidal host. They differ from the \script{heliocentric_position} relative to the most massive body only when the tidal host is not the most massive body (a satellite system, for example).

\section{Data stored with the bodies}\label{sec:pos-state}
\paragraph{Planetary tide,} in \script{KaulaParameters} of the body:
\begin{itemize}
  \item \script{love_numbers}: the spectrum, as given by the user;
  \item \script{spectrum} (\script{TideSpectrum}): the prepared spectrum (number $N$ of distinct frequencies, extrapolation exponents at both ends), computed at the first force evaluation;
  \item \script{rate_history} (\script{RateHistory}): the history of the mean rates;
  \item \script{tide_state} (\script{TideState}): the last force and torque, with their dissipative parts and the state that produced them;
  \item \script{tidal_force}: no longer used and always zero, kept so that the files written by earlier versions can still be read.
\end{itemize}
\paragraph{Stellar tide,} one per pair (tidal host, body), stored in the orbiting body: \script{Tides::host_kaula_tide} (\script{HostTide}), with the rate history of the pair and the last tide of the pair. The prepared signed spectrum of the host is stored in the \script{KaulaParameters} of the host.

\paragraph{Rate history.} \script{Particle} is \script{Copy}, so the history is a fixed-size array: $\lambda$ and $\Omega$ are accumulated at every accepted step (to keep the full turns), and a sample $(t,\lambda,\Omega,n_{\rm osc})$ is stored every $\Delta T/32$ in a ring buffer of 40 entries, which covers $1.25\,\Delta T$. The values at $t-\Delta T$ are interpolated linearly. With these settings the rates agree with those of the full history to $2\times10^{-5}n$ on $\bar\sigma_{2200}$, five times below the error of the estimator itself.

\paragraph{Update at the accepted steps.} \script{Universe::update_kaula_rate_histories} is called once per step, at the beginning of \script{iterate} of each integrator (IAS15, WHFast, LeapFrog), before the force evaluations of the step, which only read the history. It computes:
\begin{itemize}
  \item the osculating mean motions of all the bodies orbiting the tidal host, with or without a tide, from their inertial coordinates relative to the host;
  \item for each body with a planetary Kaula tide, and for each pair with a stellar Kaula tide, the window~\eqref{eq:window} and the geometry of the accepted state, with the spin $\mathbf L/I$ of that state; then it appends the step to the history.
\end{itemize}

\paragraph{Snapshots.} All the new fields are serialised, so that a run resumed from a recovery snapshot continues with the same rates and gives the same trajectory as an uninterrupted run. The fields have default values, so that the case files written by the Python interface, which give only \script{love_numbers}, are read unchanged. The \script{.bin} snapshots written by earlier versions of the Kaula model cannot be read (the binary format has no field names); the \script{.json} files can.

\section{Love-number spectra}\label{sec:pos-spectra}
The spectra are given in \script{KaulaParameters::love_numbers} as three arrays of 1024 values: \script{spectrum_excitation_frequency} ($\chi$ or $x$ in rad/s), \script{spectrum_real_part} (\texttt{rek2}) and \script{spectrum_imaginary_part} (\texttt{imk2}). They are read from three-column files and padded to 1024 values, by repeating the last row or with zeros; the padding is detected and ignored.
\begin{itemize}
  \item \textbf{Planet} (unsigned spectrum): $\chi\ge0$, increasing; $\texttt{imk2}\ge0$ for a lag. The values are used as given, with no change of sign: $\kb_2=\texttt{rek2}-\I\,\sgn(\bar\sigma)\,\texttt{imk2}$~\eqref{eq:love-impl}. A first row of physical parameters, present in some files, must be removed before loading.
  \item \textbf{Star} (signed spectrum): frequencies of both signs, increasing, with $\texttt{imk2}$ of the sign of the frequency, and optionally \script{stellar_spectrum_spin_rate}; reading~\eqref{eq:kbar-star}.
\end{itemize}

\section{Messages}\label{sec:pos-messages}
\begin{center}\small
\renewcommand{\arraystretch}{1.2}
\begin{tabular}{p{2cm}p{12.5cm}}\toprule
Type & Condition\\\midrule
warning & $e>0.5$: $q_{\max}$ is kept at 42 and the truncation error is no longer guaranteed (printed once per run)\\
warning & a forcing frequency is more than a decade outside an unsigned spectrum (above $10\chi_N$, or below $\chi_1/10$ when $\chi_1>0$): the extrapolated value is not guaranteed (once per run)\\
warning & a forcing frequency is outside a signed spectrum (once per run)\\
stop & $e\ge1$: the Kaula expansion does not exist\\
stop & invalid spectrum: fewer than 3 distinct frequencies, non-finite values, unsorted frequencies; for a planet, a negative frequency or $\texttt{imk2}<0$; for a star, no frequency of one sign, or $\texttt{imk2}$ of the sign opposite to the frequency\\
stop & a spectrum rescaled by the spin (\script{stellar_spectrum_spin_rate}) used for a body that does not rotate\\\bottomrule
\end{tabular}
\end{center}

\section{Dissipated power}\label{sec:pos-power}
The output \script{denergy_dt} of a body with a planetary Kaula tide is~\eqref{eq:power}, computed with the dissipative parts stored with the last tide. For a host with a stellar Kaula tide, it is the sum of~\eqref{eq:power} over the orbiting bodies, i.e.\ the power dissipated in the star. Before the first force evaluation, and for a body without tide, it is zero.

\section{Cost}\label{sec:pos-cost}
Time of one force evaluation for one planet, single core:
\begin{center}\small
\begin{tabular}{lccc}\toprule
$e$ & $G_{2pq}$ (tables) & single form & coplanar form\\\midrule
0.006 & 0.11\,$\mu$s & 1.9\,$\mu$s & 0.55\,$\mu$s\\
0.06 & 0.20\,$\mu$s & 5.3\,$\mu$s & 2.1\,$\mu$s\\
0.2 & 0.34\,$\mu$s & 10.3\,$\mu$s & 3.1\,$\mu$s\\
0.4 & 0.78\,$\mu$s & 28\,$\mu$s & 7.4\,$\mu$s\\\bottomrule
\end{tabular}
\end{center}
Tisserand's formula alone costs 1 to 117\,$\mu$s over the same range. The tables of $G_{2pq}$ are built once per run (about 0.1\,s).

\section{Tests}\label{sec:pos-tests}
\begin{center}\small
\renewcommand{\arraystretch}{1.2}
\begin{tabular}{p{5.4cm}p{9.1cm}}\toprule
File & Content\\\midrule
\script{tests/test_kaula_planet.rs} & functions of \script{kaula_tide.rs} against reference values computed independently in Python: nine reference cases and random states (force and torque), $F_{2mp}$, $G_{2pq}$ (tables against Tisserand), spectra (interpolation, extrapolation, padding), truncation table, mean rates, window\\
\script{tests/test_kaula_planet_coplanar.rs} & coplanar form against Python and against the single form at $i=0$; continuity at the switch\\
\script{tests/test_kaula_planet_integration.rs} & integrations (IAS15, WHFast, LeapFrog): conservation of $\mathbf L_{\rm orb}+I\mathbf s$ (V9), secular rates against the orbit averages, constant-time-lag limit against \script{ConstantTimeLag} (V3), restart from a snapshot, \script{denergy_dt}\\
\script{tests/test_kaula_stellar.rs} & signed spectra, reading~\eqref{eq:kbar-star}, stellar tide against Python and against the exact response, one tide per body, conservation and secular rates\\
\script{tests/data/kaula_planet/}, \script{tests/data/kaula_stellar/} & reference values (inputs, intermediate values, expected force and torque)\\\bottomrule
\end{tabular}
\end{center}

\paragraph{Reference cases.} \script{tests/data/kaula_planet/kaula_l2_reference_cases.json} contains nine cases in Posidonius units:
\begin{center}\small
\setlength{\tabcolsep}{4pt}
\begin{tabular}{p{5.6cm}lcccp{3.2cm}}\toprule
Case & Spectrum & $e$ & $i$ (rad) & $q_{\max}$ & Checks\\\midrule
\script{circular_equatorial_ctl} & const.\ time lag & 0 & 0 & 2 & V1\\
\script{eccentric_coplanar_tilted_plane_ctl} & const.\ time lag & 0.05 & 0 & 5 & plane $\neq xy$\\
\script{eccentric_oblique_maxwell} & Maxwell & 0.1 & 0.5 & 10 & mean rates $\neq$ Kepler\\
\script{retrograde_spin_maxwell} & Maxwell & 0.2 & 2.6 & 12 & $i>\pi/2$\\
\script{upside_down_ctl} & const.\ time lag & 0.05 & 0 & 5 & orbit and spin along $-z$\\
\script{elastic_oblique} & elastic & 0.15 & 0.8 & 12 & V4\\
\script{high_eccentricity_maxwell} & Maxwell & 0.45 & 0.3 & 42 & large $q_{\max}$\\
\script{tabulated_spectrum_inside} & test file & 0.05 & 0.2 & 7 & interpolation\\
\script{tabulated_spectrum_extrapolated} & test file & 0.1 & 0.1 & 7 & 108 of 135 terms extrapolated\\\bottomrule
\end{tabular}
\end{center}
For each case the file gives the inputs ($\vr$, $\mathbf v$, $\mathbf s$, masses, radius, $\G$, mean rates, spectrum), intermediate values ($a$, $e$, $i$, $\omega$, $\M$, $u_\M$, $n$, $\Op$), and the expected $\Fsp$ and $\boldsymbol\Gamma$ with the truncation table and with $|q|\le70$ (converged to $10^{-15}$).

\paragraph{Tolerances.}
\begin{itemize}
  \item With the same truncation, the code reproduces $\Fsp$ to $10^{-12}$ relative to $\|\Fsp\|$, and $\boldsymbol\Gamma$ to $10^{-12}$ relative to $r^\star\|\Fsp\|$; the measured differences are below $10^{-14}$.
  \item Against the converged values, the difference is the truncation error, bounded by~\eqref{eq:guarantee}: below $1.1\times10^{-6}$ for these cases.
  \item The torque is not compared relative to $\|\boldsymbol\Gamma\|$: for an elastic planet $\boldsymbol\Gamma=0$ exactly, and for a weakly dissipative one $\|\boldsymbol\Gamma\|\ll r^\star\|\Fsp\|$.
\end{itemize}

\paragraph{Stored trajectories.} The tests \script{test_tides_kaula} compare short runs with positions stored in \script{tests/data/test_tides_kaula-enabled_*}. These positions were regenerated with the present model, since the force has changed.

\section{Known limitations}\label{sec:pos-limits}
\begin{itemize}
  \item \textbf{Spin--orbit synchronism} (Section~\ref{sec:limits}). Near the 1:1 resonance the spin equation of the model becomes stiff. WHFast then prints ``convergence issue with the integration of the additional forces'', and the secular rates depend on the step for both WHFast and IAS15. Results in this regime should be checked by changing the step.
  \item \textbf{Mean rates and the integrator.} The history is sampled at the accepted steps, so the mean rates depend slightly on the step sequence of the integrator during the first windows of a run (relative difference of the spin $10^{-5}$ after 10 orbits between IAS15 and WHFast in a test case, $6\times10^{-8}$ after 200).
  \item \textbf{Error control of IAS15.} The step control of IAS15 in Posidonius uses one error estimate for the accelerations and the torques together. For a planet, the torques are many orders of magnitude smaller than the accelerations, so the integration error of the spin is not controlled by the step size.
  \item \textbf{Eccentricity} above 0.5: no truncation guarantee (warning).
  \item \textbf{Coplanar form} for prograde orbits only; a retrograde coplanar orbit uses the single form, which is correct but slower.
\end{itemize}


\appendix
\chapter{Setting, conventions and notation}\label{app:conventions}

\section{Physical setting}
We consider a star of mass $M^\star$ and a planet of mass $M_{\mathrm p}$ and radius $R$. The star is treated as a point mass. The planet is deformed by the tidal field of the star and develops a tidal bulge. We compute:
\begin{itemize}
  \item the gravitational force exerted by this bulge on the star, $\Fps$;
  \item its reaction on the planet, $\Fsp=-\Fps$;
  \item the torque that the same interaction exerts on the spin of the planet.
\end{itemize}
The tide raised on the star is not treated.

The response of the planet is assumed to satisfy three hypotheses:
\begin{description}
  \item[H1 — linear and time-invariant.] The deformation depends linearly on the forcing, through a response function that does not change over the forcing periods.
  \item[H2 — spherically symmetric body.] The planet is radially stratified at rest; its rotational flattening is neglected in the response. A forcing of degree $l$ and order $m$ then produces a response of the same $(l,m)$, with a transfer function that depends on $l$ but not on $m$.
  \item[H3 — quasi-harmonic forcing.] Each term of the expansion of the forcing has a slowly varying amplitude and a phase that advances at a slowly varying rate (Chapter~\ref{ch:tidal}).
\end{description}
The transfer function is the complex Love number $\kb_l(\sigma)$ \citep{EfroimskyMakarov2013,EfroimskyMakarov2014}.

\section{Frames and coordinates}
\begin{itemize}
  \item \textbf{Planet-centred spherical coordinates} $(r,\theta,\varphi)$:
  \begin{itemize}
    \item $\theta$ is the colatitude, measured from the spin axis of the planet;
    \item $\varphi$ is the \emph{body-fixed} longitude, i.e. it rotates with the planet. This is Kaula's convention: $\varphi=\alpha-\thetap$, where $\alpha$ is the longitude measured from a reference direction of the equatorial plane and $\thetap$ is the rotation angle of the planet from the same direction \citep[eq.~3.53--3.54]{Kaula2000}.
  \end{itemize}
  \item \textbf{Points.} The field point is $\vr=(r,\theta,\varphi)$. The star is at $\vr^\star=(r^\star,\theta^\star,\varphi^\star)$.
  \item \textbf{Local basis.} $(\er,\et,\ep)$ points outward, southward and eastward respectively.
  \item \textbf{Body-fixed Cartesian basis} $(\mathbf e_x,\mathbf e_y,\mathbf e_z)$: $\mathbf e_z$ along the spin axis, $\mathbf e_x$ towards $\varphi=0$.
  \item \textbf{Nodal frame} $(\mathbf e_X,\mathbf e_Y,\mathbf e_Z)$: $\mathbf e_Z$ along the spin axis, $\mathbf e_X$ towards the ascending node of the star's orbit on the planet's equator (Chapter~\ref{ch:second}).
\end{itemize}

\section{Orbital elements and rotation}
The relative orbit of the star around the planet is described by osculating Keplerian elements \emph{referred to the planet's equator}:
\begin{itemize}
  \item $a$ semi-major axis, $e$ eccentricity;
  \item $i$ inclination, i.e. the angle between the orbital and the spin angular momenta (the obliquity);
  \item $\omega$ argument of pericentre, $\Omega$ longitude of the ascending node, $\M$ mean anomaly;
  \item $n$ mean motion, $f$ true anomaly, $u=\omega+f$ argument of latitude.
\end{itemize}
When the two transformations must be distinguished (Chapter~\ref{ch:second}), the elements of the star carry a $\star$. The rotation angle of the planet is $\thetap$ and its spin rate is $\Op=\dot\thetap$. Only the combination $\Omega-\thetap$ ever appears, so the choice of the reference direction drops out.

\section{Spherical harmonics and Kaula functions}
The associated Legendre functions $P_l^m$ include the Condon--Shortley phase. Spherical harmonics are $Y_l^m=N_{lm}P_l^me^{\I m\varphi}$, and Kaula's associated function is $\PK_{lm}=(-1)^mP_l^m$. The definitions, the two phase conventions and all identities are collected in Appendix~\ref{app:harmonics}. Kaula's functions $F_{lmp}(i)$ and $G_{lpq}(e)$ are given in Appendix~\ref{app:kaula}.

The indices are $l\ge2$, $0\le m\le l$, $0\le p\le l$ and $q\in\mathbb Z$. A second set $(k,j)$ is used for the transformation at the field point. Negative orders $m$ never appear in the results.

\section{Love number}
The response to a term of signed frequency $\sigma$ is
\begin{equation}
  \kb_l(\sigma)=|\kb_l(\sigma)|e^{-\I\epsilon_l(\sigma)}=\Re\kb_l+\I\Im\kb_l,\qquad
  \kb_l(-\sigma)=\kb_l(\sigma)^*. \label{eq:love-def}
\end{equation}
\begin{itemize}
  \item The response kernel is real. Hence $\Re\kb_l$ is even and $\Im\kb_l$ odd in $\sigma$.
  \item With the physical frequency $\chi=|\sigma|$, we have $\kb_l(\sigma)=\kb_l(\chi)$ for $\sigma>0$ and $\kb_l(\chi)^*$ for $\sigma<0$.
  \item When the spectra are given as $\texttt{rek2}(\chi)\ge0$ and $\texttt{imk2}(\chi)\ge0$, the latter positive for a lag,
  \begin{equation}
    \kb_2(\sigma)=\texttt{rek2}(|\sigma|)-\I\,\sgn(\sigma)\,\texttt{imk2}(|\sigma|). \label{eq:love-spectra}
  \end{equation}
\end{itemize}

\section{Force convention}
The tidal-raising potential is negative: $U=-\G M^\star/\|\vr^\star-\vr\|$. The force on a mass $m_0$ is $-m_0\nabla(\text{potential})$. With $T$ the potential of the bulge (Chapter~\ref{ch:tidal}),
\begin{equation}
  \Fps=-M^\star\nabla T\big|_{\vr=\vr^\star},\qquad \Fsp=+M^\star\nabla T\big|_{\vr=\vr^\star}. \label{eq:force-convention}
\end{equation}

\section{Notation summary}
\begin{longtable}{>{$}l<{$}p{11.5cm}}
\toprule
\text{Symbol} & Meaning\\\midrule\endhead
\G & gravitational constant\\
M^\star,\ M_{\mathrm p},\ R & mass of the star, mass and radius of the planet\\
(r,\theta,\varphi) & planet-centred spherical coordinates (colatitude, body-fixed longitude)\\
\vr,\ \vr^\star & field point, position of the star seen from the planet\\
\Theta & angle between $\vr$ and $\vr^\star$, $\cos\Theta=\hat{\vr}\cdot\hat{\vr}^\star$\\
a,e,i,\omega,\Omega,\M & osculating elements of the star's orbit, referred to the planet's equator\\
n,\ f,\ u & mean motion $\sqrt{\G(M^\star+M_{\mathrm p})/a^3}$, true anomaly, argument of latitude $u=\omega+f$\\
\varpi,\ \lambda,\ u_\M & longitude of pericentre $\Omega+\omega$, mean longitude $\Omega+\omega+\M$, argument of the mean latitude $\omega+\M$\\
\bar n,\ \dot{\bar\lambda},\ \dot{\bar\Omega} & mean motion, rate of the mean longitude, precession rate of the node, averaged over the fast angles (Section~\ref{sec:rates}, Appendix~\ref{app:frequency})\\
k & $m-2+2p$, the integer multiplying $\Omega$ in the frequency at $l=2$ (Appendix~\ref{app:frequency})\\
\thetap,\ \Op & rotation angle and spin rate of the planet\\
P_l^m,\ \PK_{lm},\ Y_l^m,\ N_{lm} & Legendre functions (with / without Condon--Shortley phase), spherical harmonics, normalisation\\
F_{lmp}(i),\ G_{lpq}(e) & Kaula inclination and eccentricity functions\\
\nu_{lmpq} & Kaula's angle $(l-2p)\omega+(l-2p+q)\M+m(\Omega-\thetap)$\\
f_{lm} & parity phase $\big((-1)^{l-m}-1\big)\pi/4$, i.e. $0$ for $l-m$ even and $-\pi/2$ for $l-m$ odd\\
\zeta_{lmpq} & $\nu_{lmpq}+f_{lm}$\\
\sigma_{lmpq},\ \chi_{lmpq} & signed tidal frequency $\dot\nu_{lmpq}$, physical frequency $|\sigma_{lmpq}|$\\
\kb_l,\ \epsilon_l & complex Love number, phase lag\\
U,\ T & tidal-raising potential, tidal potential (potential of the bulge)\\
\mathcal K_{lmpq} & $\G M^\star\frac{R^{2l+1}}{a^{l+1}}\frac{(l-m)!}{(l+m)!}(2-\delta_{0,m})F_{lmp}(i)G_{lpq}(e)$\\
\mathcal A_{lmpq} & $-|\kb_l|\G M^\star\frac{R^{2l+1}}{a^{l+1}}(-1)^m\sqrt{\frac{4\pi}{2l+1}\frac{(l-m)!}{(l+m)!}}(2-\delta_{0,m})F_{lmp}G_{lpq}$\\
c^0_{lm},\ c^\pm_{lm} & ladder constants (Chapter~\ref{ch:gradient})\\
\Lambda_{lmpq} & phase at the field point, $\zeta_{lmpq}-m\varphi$\\
\hat{\mathbf s},\ \hat{\mathbf h},\ \hat{\mathbf N},\ \hat{\mathbf Y} & spin axis, orbit normal, ascending node on the equator, $\hat{\mathbf s}\times\hat{\mathbf N}$ (nodal frame, Section~\ref{sec:geometry})\\
\rho,\ c^\star,\ \lambda_N & $\sin\theta^\star e^{\I\lambda_N}$, $\cos\theta^\star$, longitude of the star from the node (Section~\ref{sec:geometry})\\
\mathcal P_{3m} & $\dd^mP_3/\dd x^m$, with $\PK_{3m}=\sin^m\theta\,\mathcal P_{3m}$ (eq.~\ref{eq:P3m})\\
\Psi_{mpq},\ \mathcal E_{mpq} & regular phase $(2-2p)u_\M+q\M+f_{2m}$ and $\kb_2(\bar\sigma_{2mpq})e^{\I\Psi_{mpq}}$ (Section~\ref{sec:implementable})\\
q_{\max}(e) & truncation of the $q$-series (Section~\ref{sec:truncation}, Appendix~\ref{app:truncation})\\
\Fsp,\ \Fps & force of the star on the planet, of the planet on the star\\
\boldsymbol\Gamma & torque on the spin of the planet, $\dd(I\mathbf s)/\dd t$\\
\bottomrule
\end{longtable}

\chapter{Legendre functions and spherical harmonics}\label{app:harmonics}

This appendix collects the definitions, conventions and identities used in the main text. The main reference is \citet{Varshalovich1988}, Chapters 5 and 7. Kaula's conventions are those of \citet[\S1.2]{Kaula2000}.

\paragraph{Convention used throughout (summary).} The phase convention is stated here once and kept throughout the document:
\begin{itemize}
  \item the associated Legendre functions $P_l^m$ include the Condon--Shortley phase $(-1)^m$, eq.~\eqref{eq:PCS};
  \item the spherical harmonics are $Y_l^m=N_{lm}P_l^me^{\I m\varphi}$, with no further $(-1)^m$, eq.~\eqref{eq:Ylm-def};
  \item Kaula's associated function has no phase factor, $\PK_{lm}=(-1)^mP_l^m$, eq.~\eqref{eq:PK-PCS}. Every formula of Kaula used with $P_l^m$ or $Y_l^m$ therefore carries a factor $(-1)^m$.
\end{itemize}

\section{Legendre polynomials}
The Legendre polynomials are generated by
\begin{equation}
  (1-2xt+t^2)^{-1/2}=\sum_{l\ge0}t^lP_l(x),\qquad |t|<1, \label{eq:legendre-gen}
\end{equation}
or given by Rodrigues' formula $P_l(x)=\frac{1}{2^ll!}\frac{\dd^l}{\dd x^l}(x^2-1)^l$. They obey the recursion
\[
  (l+1)P_{l+1}=(2l+1)xP_l-lP_{l-1}.
\]
For example, $P_0=1$, $P_1=x$, $P_2=\tfrac12(3x^2-1)$, $P_3=\tfrac12(5x^3-3x)$.

\section{Associated Legendre functions: two phase conventions}
\paragraph{Condon--Shortley convention (used here).} For $0\le m\le l$,
\begin{equation}
  P_l^m(x)=(-1)^m(1-x^2)^{m/2}\frac{\dd^mP_l(x)}{\dd x^m}, \label{eq:PCS}
\end{equation}
as in \citet[eq.~5.2(6)]{Varshalovich1988} and \citet{CondonShortley1935}.

\paragraph{Kaula's convention.} Without the phase factor \citep[eq.~1.28--1.30]{Kaula2000},
\begin{equation}
  \PK_{lm}(\sin\phi)=\cos^m\phi\sum_{t=0}^{\lfloor(l-m)/2\rfloor}T_{lmt}\sin^{l-m-2t}\phi,\qquad T_{lmt}=\frac{(-1)^t(2l-2t)!}{2^lt!(l-t)!(l-m-2t)!}, \label{eq:PK}
\end{equation}
where $\phi$ is the latitude, $\sin\phi=\cos\theta$. Hence
\begin{equation}
  \PK_{lm}=(-1)^mP_l^m. \label{eq:PK-PCS}
\end{equation}

The inclination functions $F_{lmp}$ and Kaula's transformation are derived with $\PK_{lm}$. Every formula of Kaula used with $P_l^m$ must therefore carry the factor $(-1)^m$ of \eqref{eq:PK-PCS}.

\paragraph{Low-degree functions} (Kaula's convention, $x=\cos\theta$, $s=\sin\theta$):
\begin{gather*}
  \PK_{20}=\tfrac12(3x^2-1),\qquad \PK_{21}=3xs,\qquad \PK_{22}=3s^2,\\
  \PK_{30}=\tfrac12(5x^3-3x),\qquad \PK_{31}=\tfrac32(5x^2-1)s,\qquad \PK_{32}=15xs^2,\qquad \PK_{33}=15s^3.
\end{gather*}
The degree-3 functions are $\PK_{3m}=s^m\mathcal P_{3m}(x)$, with $\mathcal P_{3m}=\dd^mP_3/\dd x^m$ given in~\eqref{eq:P3m}.

\paragraph{Values on the equator.} For $l-m$ even,
\[
  \PK_{lm}(0)=\frac{(-1)^{(l-m)/2}(l+m)!}{2^l\,\big(\tfrac{l-m}{2}\big)!\,\big(\tfrac{l+m}{2}\big)!},
\]
and $\PK_{lm}(0)=0$ for $l-m$ odd. For example $\PK_{20}(0)=-\tfrac12$, $\PK_{22}(0)=3$, $\PK_{31}(0)=-\tfrac32$, $\PK_{33}(0)=15$.

\section{Derivative and $1/\sin\theta$ identities}
Write $x=\cos\theta$, $s=\sin\theta$ and $P^{(m)}_l=\dd^mP_l/\dd x^m$, so that $\PK_{lm}=s^mP^{(m)}_l$. Two classical facts are used.

\paragraph{(a) Differentiated Legendre equation.} Differentiating $(1-x^2)y''-2xy'+n(n+1)y=0$ (with $y=P_n$) $k$ times gives
\[
  (1-x^2)P^{(k+2)}_n-2(k+1)xP^{(k+1)}_n+[n(n+1)-k(k+1)]P^{(k)}_n=0.
\]

\paragraph{(b) A three-term relation.} Differentiating $xP'_n-P'_{n-1}=nP_n$ $k$ times gives
\[
  xP^{(k+1)}_n+kP^{(k)}_n-P^{(k+1)}_{n-1}=nP^{(k)}_n.
\]

\paragraph{$\theta$-derivative.} Differentiating $\PK_{lm}=s^mP^{(m)}_l$ gives $\dd\PK_{lm}/\dd\theta=m\,s^{m-1}xP^{(m)}_l-s^{m+1}P^{(m+1)}_l$. Fact (a) with $n=l$, $k=m-1$ gives $s^2P^{(m+1)}_l=2mxP^{(m)}_l-(l+m)(l-m+1)P^{(m-1)}_l$. Using it for half of the second term,
\begin{equation}
  \frac{\dd\PK_{lm}}{\dd\theta}=\tfrac12\big[(l+m)(l-m+1)\PK_{l,m-1}-\PK_{l,m+1}\big]\quad(m\ge1),\qquad \frac{\dd\PK_{l0}}{\dd\theta}=-\PK_{l1}. \label{eq:dPK}
\end{equation}

\paragraph{Division by $\sin\theta$.} For $m\ge1$, divide by $s^{m-1}$; the claim
\begin{equation}
  \frac{m\PK_{lm}}{\sin\theta}=\tfrac12\big[\PK_{l+1,m+1}+(l-m+1)(l-m+2)\PK_{l+1,m-1}\big] \label{eq:mPK-sin}
\end{equation}
becomes
\[
  2mP^{(m)}_l=s^2P^{(m+1)}_{l+1}+(l-m+1)(l-m+2)P^{(m-1)}_{l+1}.
\]
Fact (a) with $n=l+1$, $k=m-1$ turns the right-hand side into $2m\big[xP^{(m)}_{l+1}-(l-m+2)P^{(m-1)}_{l+1}\big]$. Fact (b) with $n=l+1$, $k=m-1$ shows that the bracket equals $P^{(m)}_l$, which proves \eqref{eq:mPK-sin}.

The right-hand side contains no $1/\sin\theta$: the azimuthal component of the tidal force is regular at the poles.

\section{Spherical harmonics}
\begin{equation}
  Y_l^m(\theta,\varphi)=N_{lm}P_l^m(\cos\theta)e^{\I m\varphi},\qquad N_{lm}=\sqrt{\frac{2l+1}{4\pi}\frac{(l-m)!}{(l+m)!}}, \label{eq:Ylm-def}
\end{equation}
as in \citet[eq.~5.2(1)]{Varshalovich1988}. There is no further $(-1)^m$: the Condon--Shortley phase is already inside $P_l^m$. Equivalently, $Y_l^m=(-1)^mN_{lm}\PK_{lm}e^{\I m\varphi}$.

\paragraph{Properties} \citep[\S5.1]{Varshalovich1988}:
\begin{itemize}
  \item orthonormality: $\int Y_l^{m*}Y_{l'}^{m'}\dd\Omega=\delta_{ll'}\delta_{mm'}$;
  \item negative orders: $Y_l^{-m}=(-1)^mY_l^{m*}$;
  \item low degrees:
  \begin{gather*}
    Y_1^1=-\sqrt{\tfrac{3}{8\pi}}\sin\theta\,e^{\I\varphi},\qquad Y_2^0=\sqrt{\tfrac{5}{16\pi}}(3\cos^2\theta-1),\\
    Y_2^1=-\sqrt{\tfrac{15}{8\pi}}\sin\theta\cos\theta\,e^{\I\varphi},\qquad Y_2^2=\sqrt{\tfrac{15}{32\pi}}\sin^2\theta\,e^{2\I\varphi}.
  \end{gather*}
\end{itemize}

\paragraph{Addition theorem} \citep[eq.~5.17(9)]{Varshalovich1988}:
\begin{align*}
  P_l(\cos\Theta)&=\frac{4\pi}{2l+1}\sum_{m=-l}^lY_l^{m*}(\theta,\varphi)Y_l^m(\theta',\varphi')\\
  &=\sum_{m=0}^l\frac{(l-m)!}{(l+m)!}(2-\delta_{0,m})P_l^m(\cos\theta)P_l^m(\cos\theta')\cos m(\varphi-\varphi').
\end{align*}
The real form is independent of the phase convention.

\section{Recursion and differential relations}
From \citet[\S5.7--5.8]{Varshalovich1988}, with this convention:
\begin{align}
  -2m\cot\theta\,Y_l^m&=\sqrt{l(l+1)-m(m+1)}\,e^{-\I\varphi}Y_l^{m+1}+\sqrt{l(l+1)-m(m-1)}\,e^{\I\varphi}Y_l^{m-1},\\
  \partial_\varphi Y_l^m&=\I mY_l^m,\\
  \partial_\theta Y_l^m&=\tfrac12\sqrt{l(l+1)-m(m+1)}\,Y_l^{m+1}e^{-\I\varphi}-\tfrac12\sqrt{l(l+1)-m(m-1)}\,Y_l^{m-1}e^{\I\varphi},\\
  \partial_\theta Y_l^0&=\sqrt{l(l+1)}\,Y_l^1e^{-\I\varphi},\\
  \frac{m}{\sin\theta}Y_l^{m*}&=-\tfrac12\Big[c^+_{lm}Y_{l+1}^{m+1\,*}e^{\I\varphi}+c^-_{lm}Y_{l+1}^{m-1\,*}e^{-\I\varphi}\Big]\quad(m\ge1),\label{eq:Y-over-sin}
\end{align}
with $c^\pm_{lm}$ from \eqref{eq:c-constants}. The last relation is \eqref{eq:mPK-sin} multiplied by $(-1)^mN_{lm}e^{-\I m\varphi}$.

\section{Gradient of a solid harmonic}
\paragraph{Vector spherical harmonics} $\mathbf Y^{(\lambda)}_{LM}$ \citep[\S7.3, eq.~25--27]{Varshalovich1988}:
\begin{itemize}
  \item $\mathbf Y^{(1)}_{LM}=r\nabla Y_{LM}/\sqrt{L(L+1)}$, tangential;
  \item $\mathbf Y^{(0)}_{LM}=\I\vr\times\nabla Y_{LM}/\sqrt{L(L+1)}$, tangential;
  \item $\mathbf Y^{(-1)}_{LM}=\hat{\vr}Y_{LM}$, radial.
\end{itemize}

\paragraph{Gradient formulas} \citep[eq.~7.3(51)]{Varshalovich1988}:
\[
  \nabla\{fY_{LM}\}=f'\mathbf Y^{(-1)}_{LM}+\sqrt{L(L+1)}\frac fr\mathbf Y^{(1)}_{LM}
  =\sqrt{\tfrac{L}{2L+1}}\Big(\tfrac{\dd}{\dd r}+\tfrac{L+1}{r}\Big)f\,\mathbf Y^{L-1}_{LM}-\sqrt{\tfrac{L+1}{2L+1}}\Big(\tfrac{\dd}{\dd r}-\tfrac{L}{r}\Big)f\,\mathbf Y^{L+1}_{LM}.
\]
For $f=r^{-(L+1)}$ only the $\mathbf Y^{L+1}_{LM}$ term survives. Its Cartesian (cyclic) components, on $\mathbf e_{\pm1}=\mp(\mathbf e_x\pm\I\mathbf e_y)/\sqrt2$ and $\mathbf e_0=\mathbf e_z$, give the ladder relations \eqref{eq:ladder}.

\chapter{Kaula's inclination and eccentricity functions}\label{app:kaula}

\section{The transformation and its two functions}
Kaula's transformation \citep[\S3.3]{Kaula2000} converts a term of a spherical-harmonic potential of degree $l$ and order $m$, evaluated at the position of an orbiting body, into a double Fourier series in the body's angles. Two functions appear:
\begin{itemize}
  \item the inclination function $F_{lmp}(i)$, which results from rotating the frame from the equator to the orbital plane;
  \item the eccentricity function $G_{lpq}(e)$, which results from replacing $(r,f)$ by $(a,\M)$, i.e.\ the Hansen coefficients.
\end{itemize}
The real form of the transformation is \eqref{eq:kaula-real-CS}, the complex form \eqref{eq:kaula-complex}, and the spherical-harmonic form, with the factor $(-1)^m$ of the Condon--Shortley convention, is \eqref{eq:kaula-Ylm}. The angle is $\nu_{lmpq}=(l-2p)\omega+(l-2p+q)\M+m(\Omega-\thetap)$; the parity phase is $f_{lm}$.

\paragraph{Indices.} $l\ge2$ is the degree, $0\le m\le l$ the order, $0\le p\le l$ and $q\in\mathbb Z$ the indices of the inclination and eccentricity expansions. A second set $(k,j)$ is used when the transformation is applied at the field point (Chapter~\ref{ch:second}). Negative orders $m$ never appear in the results.

\section{Inclination function $F_{lmp}(i)$}
\paragraph{Definition} \citep[eq.~3.62]{Kaula2000}:
\begin{multline}
  F_{lmp}(i)=\sum_t\frac{(2l-2t)!}{t!(l-t)!(l-m-2t)!\,2^{2l-2t}}\sin^{l-m-2t}i\\
  \times\sum_{s=0}^m\binom ms\cos^si\sum_c\binom{l-m-2t+s}{c}\binom{m-s}{p-t-c}(-1)^{c-k}. \label{eq:F-def}
\end{multline}
Here $k=\lfloor(l-m)/2\rfloor$, $t$ runs from 0 to $\min(p,k)$, and $c$ over all values that make the binomials nonzero. The function is defined with Kaula's $\PK_{lm}$ (no Condon--Shortley phase).

\paragraph{Properties.}
\begin{itemize}
  \item Symmetry: $F_{l,m,l-p}(i)=(-1)^{l-m}F_{lmp}(\pi-i)$.
  \item Coplanar value: $F_{lmp}(0)=\PK_{lm}(0)$ if $l-m$ is even and $p=(l-m)/2$, and $0$ otherwise.
\end{itemize}
Tables~\ref{tab:F2mp} and \ref{tab:F3mp} give the exact closed forms, generated symbolically from \eqref{eq:F-def}. They agree with Kaula's Table~1.

% Generated from the exact series of the Kaula functions -- do not edit by hand
\begin{table}[ht]\centering\small
\caption{Kaula inclination functions $F_{2mp}(i)$ (Kaula 2000, eq.~3.62; exact). Rows with $F_{2mp}(0)\neq0$ are the only ones in the coplanar case.}
\label{tab:F2mp}
{\renewcommand{\arraystretch}{2.6}\begin{tabular}{ccc|>{$\displaystyle}l<{$}|>{$\displaystyle}l<{$}|>{$\displaystyle}c<{$}}\hline
$l$ & $m$ & $p$ & F_{lmp}(i) & \mathrm dF_{lmp}/\mathrm di & F_{lmp}(0)\\\hline
2 & 0 & 0 & - \frac{3 \sin^{2}{i}}{8} & - \frac{3 \sin{\left(2 i \right)}}{8} & 0\\
2 & 0 & 1 & \frac{3 \sin^{2}{i} - 2}{4} & \frac{3 \sin{\left(2 i \right)}}{4} & - \frac{1}{2}\\
2 & 0 & 2 & - \frac{3 \sin^{2}{i}}{8} & - \frac{3 \sin{\left(2 i \right)}}{8} & 0\\
\hline
2 & 1 & 0 & \frac{3 \left(\cos{i} + 1\right) \sin{i}}{4} & \frac{3 \left(\cos{i} + \cos{\left(2 i \right)}\right)}{4} & 0\\
2 & 1 & 1 & - \frac{3 \sin{\left(2 i \right)}}{4} & - \frac{3 \cos{\left(2 i \right)}}{2} & 0\\
2 & 1 & 2 & \frac{3 \left(\cos{i} - 1\right) \sin{i}}{4} & \frac{3 \left(- \cos{i} + \cos{\left(2 i \right)}\right)}{4} & 0\\
\hline
2 & 2 & 0 & \frac{3 \left(\cos{i} + 1\right)^{2}}{4} & - \frac{3 \left(\cos{i} + 1\right) \sin{i}}{2} & 3\\
2 & 2 & 1 & \frac{3 \sin^{2}{i}}{2} & \frac{3 \sin{\left(2 i \right)}}{2} & 0\\
2 & 2 & 2 & \frac{3 \left(\cos{i} - 1\right)^{2}}{4} & - \frac{3 \left(\cos{i} - 1\right) \sin{i}}{2} & 0\\
\hline
\end{tabular}}\end{table}

\begin{table}[ht]\centering\small
\caption{Kaula inclination functions $F_{3mp}(i)$ (exact; reference only, not needed by the single form).}
\label{tab:F3mp}
{\renewcommand{\arraystretch}{2.6}\begin{tabular}{ccc|>{$\displaystyle}l<{$}|>{$\displaystyle}c<{$}}\hline
$l$ & $m$ & $p$ & F_{lmp}(i) & F_{lmp}(0)\\\hline
3 & 0 & 0 & - \frac{5 \sin^{3}{i}}{16} & 0\\
3 & 0 & 1 & \frac{3 \left(5 \sin^{2}{i} - 4\right) \sin{i}}{16} & 0\\
3 & 0 & 2 & - \frac{3 \left(5 \sin^{2}{i} - 4\right) \sin{i}}{16} & 0\\
3 & 0 & 3 & \frac{5 \sin^{3}{i}}{16} & 0\\
\hline
3 & 1 & 0 & - \frac{15 \left(\cos{i} + 1\right) \sin^{2}{i}}{16} & 0\\
3 & 1 & 1 & - \frac{3 \left(\cos{i} + 10 \cos{\left(2 i \right)} + 15 \cos{\left(3 i \right)} + 6\right)}{64} & - \frac{3}{2}\\
3 & 1 & 2 & \frac{3 \left(20 \sin^{2}{i} + \cos{i} + 15 \cos{\left(3 i \right)} - 16\right)}{64} & 0\\
3 & 1 & 3 & \frac{15 \left(\cos{i} - 1\right) \sin^{2}{i}}{16} & 0\\
\hline
3 & 2 & 0 & \frac{15 \left(\cos{i} + 1\right)^{2} \sin{i}}{8} & 0\\
3 & 2 & 1 & - \frac{15 \left(\cos{i} + 1\right) \left(3 \cos{i} - 1\right) \sin{i}}{8} & 0\\
3 & 2 & 2 & \frac{15 \left(\cos{i} - 1\right) \left(3 \cos{i} + 1\right) \sin{i}}{8} & 0\\
3 & 2 & 3 & - \frac{15 \left(\cos{i} - 1\right)^{2} \sin{i}}{8} & 0\\
\hline
3 & 3 & 0 & \frac{15 \left(\cos{i} + 1\right)^{3}}{8} & 15\\
3 & 3 & 1 & \frac{45 \left(\cos{i} + 1\right) \sin^{2}{i}}{8} & 0\\
3 & 3 & 2 & - \frac{45 \left(\cos{i} - 1\right) \sin^{2}{i}}{8} & 0\\
3 & 3 & 3 & - \frac{15 \left(\cos{i} - 1\right)^{3}}{8} & 0\\
\hline
\end{tabular}}\end{table}


\section{Eccentricity function $G_{lpq}(e)$}
\paragraph{Definition.} $G_{lpq}$ is defined by Kaula's replacement \citep[p.~35]{Kaula2000}:
\begin{equation}
  \Big(\frac ar\Big)^{l+1}e^{\I(l-2p)f}=\sum_{q\in\mathbb Z}G_{lpq}(e)\,e^{\I(l-2p+q)\M}, \label{eq:G-def}
\end{equation}
i.e.\ $G_{lpq}$ is a Hansen coefficient:
\[
  G_{lpq}(e)=\frac1{2\pi}\int_0^{2\pi}\Big(\frac ar\Big)^{l+1}e^{\I(l-2p)f}e^{-\I(l-2p+q)\M}\,\dd\M.
\]

\paragraph{Properties.}
\begin{itemize}
  \item $G_{lpq}$ is real.
  \item Symmetry: $G_{l,l-p,-q}=G_{lpq}$.
  \item $G_{lpq}=\mathcal O(e^{|q|})$.
\end{itemize}

\paragraph{Tisserand's formula} \citep[eq.~3.67--3.69]{Kaula2000}, after \citet[p.~256]{Tisserand1889}. It holds for any $q$:
\begin{equation}
  G_{lpq}(e)=(-1)^{|q|}(1+\beta^2)^l\beta^{|q|}\sum_{k\ge0}P_{lpqk}Q_{lpqk}\beta^{2k},\qquad \beta=\frac{e}{1+\sqrt{1-e^2}}, \label{eq:tisserand}
\end{equation}
\[
  P_{lpqk}=\sum_{r=0}^{h}\binom{2p'-2l}{h-r}\frac{(-1)^r}{r!}\Big(\frac{(l-2p'+q')e}{2\beta}\Big)^r,\qquad h=\begin{cases}k+q',&q'>0,\\k,&q'\le0,\end{cases}
\]
\[
  Q_{lpqk}=\sum_{r=0}^{h}\binom{-2p'}{h-r}\frac{1}{r!}\Big(\frac{(l-2p'+q')e}{2\beta}\Big)^r,\qquad h=\begin{cases}k,&q'>0,\\k-q',&q'\le0,\end{cases}
\]
where:
\begin{itemize}
  \item $p'=p$, $q'=q$ for $p\le l/2$, and $p'=l-p$, $q'=-q$ for $p>l/2$;
  \item the binomials are generalised, $\binom nj=n(n-1)\cdots(n-j+1)/j!$ for any integer $n$.
\end{itemize}

\paragraph{Convergence.} The series converges geometrically in $\beta^2\approx e^2/4$. Number of $k$-terms for an absolute error below $10^{-13}$ on all $G_{2pq}$ needed by the truncation table of Chapter~\ref{ch:nbody}:
\begin{center}
\begin{tabular}{l|ccccccc}\toprule
$e$ & 0.01 & 0.05 & 0.1 & 0.2 & 0.3 & 0.4 & 0.5\\\midrule
$K$ & 4 & 5 & 6 & 8 & 11 & 14 & 22\\\bottomrule
\end{tabular}
\end{center}

\paragraph{Closed forms.} For $l-2p+q=0$ \citep[eq.~3.66]{Kaula2000}, for example $G_{210}=(1-e^2)^{-3/2}$ and $G_{31,-1}=e(1-e^2)^{-5/2}$. Also $G_{20,-2}=G_{222}=0$ exactly.

\paragraph{Power series.} Tables~\ref{tab:G20q}--\ref{tab:G21q} give the exact power series in $e$ for $l=2$, $|q|\le24$, obtained by expanding \eqref{eq:tisserand} in exact rational arithmetic. They extend Kaula's Table~2 and the tables of \citet{Cayley1861}. The series to $e^{40}$, with exact rational coefficients, are available as a data file accompanying this document.

Accuracy of the truncated series, as the maximum over the $G_{2pq}$ needed by the table:
\begin{center}
\begin{tabular}{l|cccccc}\toprule
order $N$ $\backslash$ $e$ & 0.05 & 0.1 & 0.2 & 0.3 & 0.4 & 0.5\\\midrule
7 & $3\times10^{-9}$ & $3\times10^{-6}$ & $8\times10^{-4}$ & $2\times10^{-2}$ & $0.2$ & $0.9$\\
15 & $3\times10^{-16}$ & $4\times10^{-15}$ & $9\times10^{-8}$ & $7\times10^{-5}$ & $6\times10^{-3}$ & $0.2$\\
20 & -- & -- & $8\times10^{-12}$ & $2\times10^{-6}$ & $6\times10^{-4}$ & $6\times10^{-2}$\\
30 & -- & -- & -- & $7\times10^{-12}$ & $6\times10^{-6}$ & $5\times10^{-3}$\\\bottomrule
\end{tabular}
\end{center}
A series truncated at $e^{30}$ is accurate to $7\times10^{-12}$ up to $e=0.3$ (for the reference row $10^{-5}$ of the truncation table of Chapter~\ref{ch:nbody}; the data file stops at $|q|=30$, and the series is taken as zero beyond). Beyond that, Tisserand's formula should be evaluated directly.

\paragraph{Exact values.} In Tables~\ref{tab:G20q} and~\ref{tab:G21q}, fractions longer than 12 digits are printed as decimals with 8 significant digits (relative rounding below $5\times10^{-8}$). The exact fractions are provided with this document:
\begin{itemize}
  \item \texttt{G\_exact\_fractions.pdf}: the same terms as Tables~\ref{tab:G20q}--\ref{tab:G21q}, each coefficient as an exact fraction and in double precision;
  \item \texttt{data/G\_exact\_l2.txt}: the complete series of $G_{20q}$ and $G_{21q}$, $|q|\le30$, to $e^{40}$, as exact fractions;
  \item \texttt{data/G\_series.json}: the same in machine-readable form, also for $l=3$.
\end{itemize}

% Generated from the exact series of the Kaula functions -- do not edit by hand
% Requires \usepackage{longtable,array}. Exact rational coefficients; full series to e^40 in G_series.json
{\renewcommand{\arraystretch}{2.2}\begin{longtable}{ccc|>{$\displaystyle}l<{$}}
\caption{Kaula eccentricity functions $G_{20q}(e)$, exact series to $e^{10}$, $|q|\le24$ (Tisserand, Kaula 2000 eq.~3.67--3.69); for $|q|>6$ the three leading terms. Long fractions are printed with 8 significant digits; the exact fractions of all these terms are given in the separate annex \texttt{annex\_G\_exact/G\_exact\_fractions.pdf}. $G_{22q}=G_{20,-q}$.}\label{tab:G20q}\\
\hline $l$ & $p$ & $q$ & G_{lpq}(e)\\\hline\endfirsthead
\hline $l$ & $p$ & $q$ & G_{lpq}(e)\\\hline\endhead
2 & 0 & -24 & 1.5875184\!\times\!10^{1}\,e^{24} - 5.4928137\!\times\!10^{1}\,e^{26} + 9.8184961\!\times\!10^{1}\,e^{28} + \mathcal O(e^{29})\\
2 & 0 & -23 & 1.1881177\!\times\!10^{1}\,e^{23} - 3.8242537\!\times\!10^{1}\,e^{25} + 6.4178774\!\times\!10^{1}\,e^{27} + \mathcal O(e^{28})\\
2 & 0 & -22 & 8.8967914\!\times\!10^{0}\,e^{22} - 2.6496966\!\times\!10^{1}\,e^{24} + 4.1623123\!\times\!10^{1}\,e^{26} + \mathcal O(e^{27})\\
2 & 0 & -21 & 6.6657926\!\times\!10^{0}\,e^{21} - 1.8255182\!\times\!10^{1}\,e^{23} + 2.6760737\!\times\!10^{1}\,e^{25} + \mathcal O(e^{26})\\
2 & 0 & -20 & 4.9971867\!\times\!10^{0}\,e^{20} - 1.2492967\!\times\!10^{1}\,e^{22} + 1.7039920\!\times\!10^{1}\,e^{24} + \mathcal O(e^{25})\\
2 & 0 & -19 & 3.7485643\!\times\!10^{0}\,e^{19} - 8.4811268\!\times\!10^{0}\,e^{21} + 1.0735007\!\times\!10^{1}\,e^{23} + \mathcal O(e^{24})\\
2 & 0 & -18 & 2.8137062\!\times\!10^{0}\,e^{18} - 5.7014573\!\times\!10^{0}\,e^{20} + 6.6844033\!\times\!10^{0}\,e^{22} + \mathcal O(e^{23})\\
2 & 0 & -17 & 2.1133575\!\times\!10^{0}\,e^{17} - 3.7864321\!\times\!10^{0}\,e^{19} + 4.1102717\!\times\!10^{0}\,e^{21} + \mathcal O(e^{22})\\
2 & 0 & -16 & 1.5883604\!\times\!10^{0}\,e^{16} - 2.4759735\!\times\!10^{0}\,e^{18} + 2.4947899\!\times\!10^{0}\,e^{20} + \mathcal O(e^{21})\\
2 & 0 & -15 & 1.1945401\!\times\!10^{0}\,e^{15} - 1.5864986\!\times\!10^{0}\,e^{17} + 1.4955082\!\times\!10^{0}\,e^{19} + \mathcal O(e^{20})\\
2 & 0 & -14 & \frac{157464}{175175}e^{14} - \frac{78732}{79625}e^{16} + 8.8765920\!\times\!10^{-1}\,e^{18} + \mathcal O(e^{19})\\
2 & 0 & -13 & 6.7675995\!\times\!10^{-1}\,e^{13} - 5.9216495\!\times\!10^{-1}\,e^{15} + 5.2499250\!\times\!10^{-1}\,e^{17} + \mathcal O(e^{18})\\
2 & 0 & -12 & 5.0968645\!\times\!10^{-1}\,e^{12} - 3.3325652\!\times\!10^{-1}\,e^{14} + 3.1330314\!\times\!10^{-1}\,e^{16} + \mathcal O(e^{17})\\
2 & 0 & -11 & 3.8386802\!\times\!10^{-1}\,e^{11} - 1.6794226\!\times\!10^{-1}\,e^{13} + 1.9262608\!\times\!10^{-1}\,e^{15} + \mathcal O(e^{16})\\
2 & 0 & -10 & \frac{4096}{14175}e^{10} - \frac{2048}{31185}e^{12} + \frac{58624}{467775}e^{14} + \mathcal O(e^{15})\\
2 & 0 & -9 & 2.1719477\!\times\!10^{-1}\,e^{9} - 5.4298693\!\times\!10^{-3}\,e^{11} + 8.8173673\!\times\!10^{-2}\,e^{13} + \mathcal O(e^{14})\\
2 & 0 & -8 & \frac{729}{4480}e^{8} + \frac{243}{8960}e^{10} + \frac{24057}{358400}e^{12} + \mathcal O(e^{13})\\
2 & 0 & -7 & \frac{15625}{129024}e^{7} + 4.1628641\!\times\!10^{-2}\,e^{9} + 5.3875451\!\times\!10^{-2}\,e^{11} + \mathcal O(e^{12})\\
2 & 0 & -6 & \frac{4}{45}e^{6} + \frac{2}{45}e^{8} + \frac{11}{252}e^{10} + \mathcal O(e^{11})\\
2 & 0 & -5 & \frac{81}{1280}e^{5} + \frac{81}{2048}e^{7} + \frac{19359}{573440}e^{9} + \mathcal O(e^{11})\\
2 & 0 & -4 & \frac{1}{24}e^{4} + \frac{7}{240}e^{6} + \frac{131}{5760}e^{8} + \frac{4303}{241920}e^{10} + \mathcal O(e^{11})\\
2 & 0 & -3 & \frac{1}{48}e^{3} + \frac{11}{768}e^{5} + \frac{313}{30720}e^{7} + \frac{3355}{442368}e^{9} + \mathcal O(e^{11})\\
2 & 0 & -2 & 0\\
2 & 0 & -1 & -\frac{1}{2}e + \frac{1}{16}e^{3} - \frac{5}{384}e^{5} - \frac{143}{18432}e^{7} - \frac{9097}{1474560}e^{9} + \mathcal O(e^{11})\\
2 & 0 & 0 & 1 - \frac{5}{2}e^{2} + \frac{13}{16}e^{4} - \frac{35}{288}e^{6} - \frac{5}{576}e^{8} - \frac{49}{3600}e^{10} + \mathcal O(e^{11})\\
2 & 0 & 1 & \frac{7}{2}e - \frac{123}{16}e^{3} + \frac{489}{128}e^{5} - \frac{1763}{2048}e^{7} + \frac{13527}{163840}e^{9} + \mathcal O(e^{11})\\
2 & 0 & 2 & \frac{17}{2}e^{2} - \frac{115}{6}e^{4} + \frac{601}{48}e^{6} - \frac{1423}{360}e^{8} + \frac{48619}{69120}e^{10} + \mathcal O(e^{11})\\
2 & 0 & 3 & \frac{845}{48}e^{3} - \frac{32525}{768}e^{5} + \frac{208225}{6144}e^{7} - 1.3840795\!\times\!10^{1}\,e^{9} + \mathcal O(e^{11})\\
2 & 0 & 4 & \frac{533}{16}e^{4} - \frac{13827}{160}e^{6} + \frac{104127}{1280}e^{8} - \frac{146017}{3584}e^{10} + \mathcal O(e^{11})\\
2 & 0 & 5 & \frac{228347}{3840}e^{5} - \frac{3071075}{18432}e^{7} + 1.7969544\!\times\!10^{2}\,e^{9} + \mathcal O(e^{11})\\
2 & 0 & 6 & \frac{73369}{720}e^{6} - \frac{775727}{2520}e^{8} + \frac{7527241}{20160}e^{10} + \mathcal O(e^{11})\\
2 & 0 & 7 & 1.6942345\!\times\!10^{2}\,e^{7} - 5.5010469\!\times\!10^{2}\,e^{9} + 7.4011937\!\times\!10^{2}\,e^{11} + \mathcal O(e^{12})\\
2 & 0 & 8 & \frac{634085}{2304}e^{8} - 9.5714070\!\times\!10^{2}\,e^{10} + 1.4131763\!\times\!10^{3}\,e^{12} + \mathcal O(e^{13})\\
2 & 0 & 9 & 4.3887349\!\times\!10^{2}\,e^{9} - 1.6292003\!\times\!10^{3}\,e^{11} + 2.6171604\!\times\!10^{3}\,e^{13} + \mathcal O(e^{14})\\
2 & 0 & 10 & 6.8943970\!\times\!10^{2}\,e^{10} - 2.7227149\!\times\!10^{3}\,e^{12} + 4.7253268\!\times\!10^{3}\,e^{14} + \mathcal O(e^{15})\\
2 & 0 & 11 & 1.0696726\!\times\!10^{3}\,e^{11} - 4.4798495\!\times\!10^{3}\,e^{13} + 8.3503215\!\times\!10^{3}\,e^{15} + \mathcal O(e^{16})\\
2 & 0 & 12 & 1.6423080\!\times\!10^{3}\,e^{12} - 7.2729041\!\times\!10^{3}\,e^{14} + 1.4486915\!\times\!10^{4}\,e^{16} + \mathcal O(e^{17})\\
2 & 0 & 13 & 2.4990369\!\times\!10^{3}\,e^{13} - 1.1670721\!\times\!10^{4}\,e^{15} + 2.4735012\!\times\!10^{4}\,e^{17} + \mathcal O(e^{18})\\
2 & 0 & 14 & 3.7734131\!\times\!10^{3}\,e^{14} - 1.8537629\!\times\!10^{4}\,e^{16} + 4.1646006\!\times\!10^{4}\,e^{18} + \mathcal O(e^{19})\\
2 & 0 & 15 & 5.6593667\!\times\!10^{3}\,e^{15} - 2.9180345\!\times\!10^{4}\,e^{17} + 6.9257766\!\times\!10^{4}\,e^{19} + \mathcal O(e^{20})\\
2 & 0 & 16 & 8.4377060\!\times\!10^{3}\,e^{16} - 4.5565401\!\times\!10^{4}\,e^{18} + 1.1391686\!\times\!10^{5}\,e^{20} + \mathcal O(e^{21})\\
2 & 0 & 17 & 1.2513987\!\times\!10^{4}\,e^{17} - 7.0639971\!\times\!10^{4}\,e^{19} + 1.8553645\!\times\!10^{5}\,e^{21} + \mathcal O(e^{22})\\
2 & 0 & 18 & 1.8472543\!\times\!10^{4}\,e^{18} - 1.0880399\!\times\!10^{5}\,e^{20} + 2.9951264\!\times\!10^{5}\,e^{22} + \mathcal O(e^{23})\\
2 & 0 & 19 & 2.7153443\!\times\!10^{4}\,e^{19} - 1.6660308\!\times\!10^{5}\,e^{21} + 4.7963238\!\times\!10^{5}\,e^{23} + \mathcal O(e^{24})\\
2 & 0 & 20 & 3.9761965\!\times\!10^{4}\,e^{20} - 2.5374321\!\times\!10^{5}\,e^{22} + 7.6247006\!\times\!10^{5}\,e^{24} + \mathcal O(e^{25})\\
2 & 0 & 21 & 5.8024101\!\times\!10^{4}\,e^{21} - 3.8457302\!\times\!10^{5}\,e^{23} + 1.2040123\!\times\!10^{6}\,e^{25} + \mathcal O(e^{26})\\
2 & 0 & 22 & 8.4407170\!\times\!10^{4}\,e^{22} - 5.8024478\!\times\!10^{5}\,e^{24} + 1.8896090\!\times\!10^{6}\,e^{26} + \mathcal O(e^{27})\\
2 & 0 & 23 & 1.2243244\!\times\!10^{5}\,e^{23} - 8.7185858\!\times\!10^{5}\,e^{25} + 2.9488753\!\times\!10^{6}\,e^{27} + \mathcal O(e^{28})\\
2 & 0 & 24 & 1.7711759\!\times\!10^{5}\,e^{24} - 1.3050284\!\times\!10^{6}\,e^{26} + 4.5779468\!\times\!10^{6}\,e^{28} + \mathcal O(e^{29})\\
\hline\end{longtable}}

{\renewcommand{\arraystretch}{2.2}\begin{longtable}{ccc|>{$\displaystyle}l<{$}}
\caption{Kaula eccentricity functions $G_{21q}(e)$, exact series to $e^{10}$, $|q|\le24$ (Tisserand, Kaula 2000 eq.~3.67--3.69); for $|q|>6$ the three leading terms. Long fractions are printed with 8 significant digits; the exact fractions of all these terms are given in the separate annex \texttt{annex\_G\_exact/G\_exact\_fractions.pdf}. $G_{22q}=G_{20,-q}$.}\label{tab:G21q}\\
\hline $l$ & $p$ & $q$ & G_{lpq}(e)\\\hline\endfirsthead
\hline $l$ & $p$ & $q$ & G_{lpq}(e)\\\hline\endhead
2 & 1 & -24 & 3.9497719\!\times\!10^{3}\,e^{24} - 1.9988026\!\times\!10^{4}\,e^{26} + 4.9741788\!\times\!10^{4}\,e^{28} + \mathcal O(e^{29})\\
2 & 1 & -23 & 2.8594427\!\times\!10^{3}\,e^{23} - 1.3744333\!\times\!10^{4}\,e^{25} + 3.2527219\!\times\!10^{4}\,e^{27} + \mathcal O(e^{28})\\
2 & 1 & -22 & 2.0688306\!\times\!10^{3}\,e^{22} - 9.4183356\!\times\!10^{3}\,e^{24} + 2.1144842\!\times\!10^{4}\,e^{26} + \mathcal O(e^{27})\\
2 & 1 & -21 & 1.4958257\!\times\!10^{3}\,e^{21} - 6.4291475\!\times\!10^{3}\,e^{23} + 1.3656270\!\times\!10^{4}\,e^{25} + \mathcal O(e^{26})\\
2 & 1 & -20 & 1.0807498\!\times\!10^{3}\,e^{20} - 4.3698127\!\times\!10^{3}\,e^{22} + 8.7564646\!\times\!10^{3}\,e^{24} + \mathcal O(e^{25})\\
2 & 1 & -19 & 7.8024284\!\times\!10^{2}\,e^{19} - 2.9557342\!\times\!10^{3}\,e^{21} + 5.5698964\!\times\!10^{3}\,e^{23} + \mathcal O(e^{24})\\
2 & 1 & -18 & 5.6281179\!\times\!10^{2}\,e^{18} - 1.9882666\!\times\!10^{3}\,e^{20} + 3.5114386\!\times\!10^{3}\,e^{22} + \mathcal O(e^{23})\\
2 & 1 & -17 & 4.0559164\!\times\!10^{2}\,e^{17} - 1.3290452\!\times\!10^{3}\,e^{19} + 2.1916798\!\times\!10^{3}\,e^{21} + \mathcal O(e^{22})\\
2 & 1 & -16 & 2.9198810\!\times\!10^{2}\,e^{16} - 8.8191037\!\times\!10^{2}\,e^{18} + 1.3526485\!\times\!10^{3}\,e^{20} + \mathcal O(e^{21})\\
2 & 1 & -15 & 2.0996289\!\times\!10^{2}\,e^{15} - 5.8019774\!\times\!10^{2}\,e^{17} + 8.2431081\!\times\!10^{2}\,e^{19} + \mathcal O(e^{20})\\
2 & 1 & -14 & 1.5078685\!\times\!10^{2}\,e^{14} - 3.7781499\!\times\!10^{2}\,e^{16} + 4.9522083\!\times\!10^{2}\,e^{18} + \mathcal O(e^{19})\\
2 & 1 & -13 & 1.0813325\!\times\!10^{2}\,e^{13} - 2.4298749\!\times\!10^{2}\,e^{15} + 2.9278831\!\times\!10^{2}\,e^{17} + \mathcal O(e^{18})\\
2 & 1 & -12 & 7.7418999\!\times\!10^{1}\,e^{12} - 1.5388210\!\times\!10^{2}\,e^{14} + 1.7006222\!\times\!10^{2}\,e^{16} + \mathcal O(e^{17})\\
2 & 1 & -11 & 5.5325854\!\times\!10^{1}\,e^{11} - 9.5550955\!\times\!10^{1}\,e^{13} + 9.6915004\!\times\!10^{1}\,e^{15} + \mathcal O(e^{16})\\
2 & 1 & -10 & 3.9452758\!\times\!10^{1}\,e^{10} - 5.7801470\!\times\!10^{1}\,e^{12} + 5.4188878\!\times\!10^{1}\,e^{14} + \mathcal O(e^{15})\\
2 & 1 & -9 & 2.8063401\!\times\!10^{1}\,e^{9} - 3.3715196\!\times\!10^{1}\,e^{11} + 2.9828532\!\times\!10^{1}\,e^{13} + \mathcal O(e^{14})\\
2 & 1 & -8 & \frac{178331}{8960}e^{8} - 1.8620290\!\times\!10^{1}\,e^{10} + 1.6340336\!\times\!10^{1}\,e^{12} + \mathcal O(e^{13})\\
2 & 1 & -7 & \frac{432091}{30720}e^{7} - 9.3806244\!\times\!10^{0}\,e^{9} + 9.1335431\!\times\!10^{0}\,e^{11} + \mathcal O(e^{12})\\
2 & 1 & -6 & \frac{3167}{320}e^{6} - \frac{2187}{560}e^{8} + \frac{195147}{35840}e^{10} + \mathcal O(e^{11})\\
2 & 1 & -5 & \frac{1773}{256}e^{5} - \frac{4987}{6144}e^{7} + \frac{418315}{114688}e^{9} + \mathcal O(e^{11})\\
2 & 1 & -4 & \frac{77}{16}e^{4} + \frac{129}{160}e^{6} + \frac{899}{320}e^{8} + \frac{107279}{40320}e^{10} + \mathcal O(e^{11})\\
2 & 1 & -3 & \frac{53}{16}e^{3} + \frac{393}{256}e^{5} + \frac{24753}{10240}e^{7} + \frac{211557}{81920}e^{9} + \mathcal O(e^{11})\\
2 & 1 & -2 & \frac{9}{4}e^{2} + \frac{7}{4}e^{4} + \frac{141}{64}e^{6} + \frac{197}{80}e^{8} + \frac{62401}{23040}e^{10} + \mathcal O(e^{11})\\
2 & 1 & -1 & \frac{3}{2}e + \frac{27}{16}e^{3} + \frac{261}{128}e^{5} + \frac{14309}{6144}e^{7} + \frac{423907}{163840}e^{9} + \mathcal O(e^{11})\\
2 & 1 & 0 & 1 + \frac{3}{2}e^{2} + \frac{15}{8}e^{4} + \frac{35}{16}e^{6} + \frac{315}{128}e^{8} + \frac{693}{256}e^{10} + \mathcal O(e^{11})\\
\hline\end{longtable}}

\noindent Symmetries: $G_{l,l-p,-q}(e)=G_{lpq}(e)$; for $l=2$, $G_{21q}=G_{21,-q}$ and $G_{22q}=G_{20,-q}$.
Leading order: $G_{2pq}=\mathcal O(e^{|q|})$. Closed forms (Kaula eq.~3.66, $l-2p+q=0$): $G_{210}=(1-e^2)^{-3/2}$, $G_{20,-2}=G_{222}=0$.


\section{Erratum of earlier tables}
The exact tables reveal the following errors in earlier working tables. They are listed so that they are not reused.
\begin{center}\small
\begin{tabular}{lll}\toprule
Function & written & exact\\\midrule
$G_{20,+2}$, coefficient of $e^4$ & $-115/16$ & $-115/6$ (as in Kaula, Table~2)\\
$G_{22,-2}=G_{20,+2}$ & $\tfrac{17}{2}e^2-\tfrac{115}{16}e^3$ & $\tfrac{17}{2}e^2-\tfrac{115}{6}e^4$\\
$G_{31,-6}$ & $\tfrac{48203}{9240}e^6$ & $\tfrac{48203}{9216}e^6$\\
$G_{310}$ & $1+\tfrac12e^2+\tfrac{239}{64}e^4-\tfrac{3323}{576}e^6$ & $1+2e^2+\tfrac{239}{64}e^4+\tfrac{3323}{576}e^6$\\
$G_{31,q}$, $q=3..7$ & copied from $G_{30,q}$ & e.g.\ $G_{313}=\tfrac{77}{6}e^3-\tfrac{25}{48}e^5+\dots$\\
$F_{311}$ & $\tfrac{15}{16}\sin^2i(1+3\cos i)-\tfrac34(1+3\cos i)$ & $\tfrac{15}{16}\sin^2i(1+3\cos i)-\tfrac34(1+\cos i)$\\
$F_{312}$ & $\tfrac{15}{16}\sin^2i(1-3\cos i)-\tfrac34(1-3\cos i)$ & $\tfrac{15}{16}\sin^2i(1-3\cos i)-\tfrac34(1-\cos i)$\\
$F_{331}$, $F_{332}$ & $\tfrac{48}{8}\sin^2i(1\pm\cos i)$ & $\tfrac{45}{8}\sin^2i(1\pm\cos i)$\\
$\dd F_{201}/\dd i$ & $\tfrac32\sin i\cos i-\tfrac12$ & $\tfrac32\sin i\cos i$\\\bottomrule
\end{tabular}
\end{center}

\chapter{The double transformation as an implementable form}\label{app:double}

Chapter~\ref{ch:nbody} uses the \emph{single} transformation: Kaula's expansion is applied to the tidal-raising potential only, and the gradient is evaluated at the true position of the star. The \emph{double} transformation of Chapter~\ref{ch:second} also expands the field point in the star's elements, then identifies the elements, eq.~\eqref{eq:grad-node}. It is kept here as an analytical reference.

\section{Form in the nodal frame}
At $l=2$, with $\mathcal K_{2mpq}$ of~\eqref{eq:K2}, $\mathcal E^{D}_{mpq,kj}=\kb_2(\bar\sigma_{2mpq})\,e^{\I D_{pqkj}}$ and $\mathcal E^{S}_{pq,kj}=\kb_2(\bar\sigma_{20pq})\,e^{\I S_{pqkj}}$, where $D$ and $S$ are given by~\eqref{eq:D}--\eqref{eq:S}, eq.~\eqref{eq:grad-node} reads
\begin{equation}
  \begin{aligned}
  V_Z&=\frac{1}{a^4}\sum_{m,p,q}\sum_{k=0}^{3}\sum_{j}(-1)^{3-m}(3-m)\,\mathcal K_{2mpq}\,F_{3mk}(i)\,G_{3kj}(e)\,\Im\,\mathcal E^{D}_{mpq,kj},\\
  V_X+\I V_Y&=\frac{1}{2a^4}\sum_{m=1}^{2}\sum_{p,q,k,j}\mathcal K_{2mpq}\,G_{3kj}\Big[F_{3,m+1,k}\,\mathcal E^{D*}_{mpq,kj}-(3-m)(4-m)\,F_{3,m-1,k}\,\mathcal E^{D}_{mpq,kj}\Big]\\
  &\quad+\frac{1}{2a^4}\sum_{p,q,k,j}\mathcal K_{20pq}\,F_{31k}\,G_{3kj}\Big[\mathcal E^{D*}_{0pq,kj}+\mathcal E^{S}_{pq,kj}\Big],
  \end{aligned} \label{eq:impl-double}
\end{equation}
with $F_{34k}\equiv0$. The force follows from~\eqref{eq:impl-force}.

\section{Comparison with the single form}
\begin{itemize}
  \item \textbf{Agreement.} With the same truncation in $q$ and $|j|\le40$, the two forms agree to $7\times10^{-15}$.
  \item \textbf{Inputs.} The double form uses the elements only: $a$, $e$, $i$, $\omega$, $\M$. No geometric coordinate of the star remains. It also needs $F_{3mk}(i)$ (Table~\ref{tab:F3mp}) and $G_{3kj}(e)$.
  \item \textbf{Cost.} $36\,(2q_{\max}+1)(2j_{\max}+1)$ terms, against $9\,(2q_{\max}+1)$ for the single form: 4\,356 against 99 at $e\le0.1$ with $j_{\max}=q_{\max}$. The ratio grows like $16q_{\max}$.
  \item \textbf{Accuracy at equal truncation.} The double form truncates both $q$ and $j$. At equal truncation it is 2 to 3 times less accurate than the single form: for $|q|,|j|\le Q$, the relative error is $5.5\times10^{-3}$ against $1.7\times10^{-3}$ at $e=0.05$ ($Q=2$), and $7.8\times10^{-2}$ against $2.3\times10^{-2}$ at $e=0.25$ ($Q=5$).
\end{itemize}
The single form is therefore the one to implement. The double form is the natural starting point of a secular theory: averaging over $\M$ selects $j=q$ (and $k=p$ for the average over $\omega$) in~\eqref{eq:impl-double}.

\chapter{Truncation tables}\label{app:truncation}

The method is described in Section~\ref{sec:truncation}. For each tolerance tol, the table gives the smallest $q_{\max}$ such that the truncation error of the single form satisfies, for any Love-number spectrum with $|\kb_2|\le\kappa_{\max}$,
\[
  \|\delta\vF\|\le\mathrm{tol}\cdot\kappa_{\max}F_0(r^\star),\qquad \|\delta\boldsymbol\Gamma\|\le\mathrm{tol}\cdot\kappa_{\max}\,r^\star F_0(r^\star),
\]
at every point of the orbit, where $F_0(r^\star)=3\G M^{\star2}R^5/(r^\star)^7$ is the instantaneous elastic force per unit Love number. Each entry is the worst case over the points of a 35-point grid in $e$ inside the bin, with 714 states per point: 7 inclinations, 3 arguments of pericentre, and 34 positions on the orbit including apocentre and pericentre. The worst state is always near apocentre, where the force is weakest: a table built with a force scale averaged over the orbit would not hold there for $e\ge0.3$. The row $10^{-5}$ is the one used in Chapter~\ref{ch:nbody}.
\begin{center}
\renewcommand{\arraystretch}{1.2}\setlength{\tabcolsep}{4pt}\small
\begin{tabular}{l|cccccccccccc}\toprule
$e\le$ & 0.001 & 0.01 & 0.02 & 0.05 & 0.10 & 0.15 & 0.20 & 0.25 & 0.30 & 0.35 & 0.40 & 0.50\\\midrule
tol $10^{-2}$ & 0 & 1 & 1 & 2 & 3 & 5 & 6 & 8 & 10 & 13 & 16 & 26\\
tol $10^{-3}$ & 1 & 2 & 2 & 3 & 5 & 6 & 8 & 10 & 13 & 16 & 20 & 31\\
tol $10^{-4}$ & 1 & 2 & 3 & 4 & 6 & 8 & 10 & 13 & 16 & 19 & 24 & 37\\
\textbf{tol $10^{-5}$} & \textbf{2} & \textbf{3} & \textbf{3} & \textbf{5} & \textbf{7} & \textbf{10} & \textbf{12} & \textbf{15} & \textbf{18} & \textbf{23} & \textbf{28} & \textbf{42}\\
tol $10^{-6}$ & 2 & 3 & 4 & 6 & 8 & 11 & 14 & 17 & 21 & 26 & 32 & 48\\\bottomrule
\end{tabular}
\end{center}

\paragraph{Choice of the tolerance.} The scale $\kappa_{\max}F_0$ is set mostly by the elastic part of the response. The secular evolution (migration, spin evolution) is driven by the dissipative part, $\Im\kb_2$, often $10^2$ to $10^3$ times smaller. To resolve it to $10^{-3}$, choose tol $\approx10^{-3}|\Im\kb_2|/\kappa_{\max}$, i.e.\ one of the rows $10^{-4}$ to $10^{-6}$; the reference row $10^{-5}$ corresponds to $|\Im\kb_2|/\kappa_{\max}\approx10^{-2}$. The cost is $9(2q_{\max}+1)$ terms: 45 to 765 for the row $10^{-5}$, against 27 to 567 for the row $10^{-3}$. This bound is a worst case over all spectra; for a given rheology the actual error is smaller.

\paragraph{Number of terms of Tisserand's series.} For an absolute error below $10^{-13}$ on all $G_{2pq}$ of the row $10^{-5}$: 4 terms at $e=0.01$, 6 at $e=0.1$, 8 at $e=0.2$, 11 at $e=0.3$, 14 at $e=0.4$ and 22 at $e=0.5$; for the row $10^{-3}$ the same, except 21 at $e=0.5$.

\chapter{The tidal frequency in three equivalent forms}\label{app:frequency}

The frequency of the term $(l,m,p,q)$ is the time derivative of Kaula's angle. This appendix writes it in three equivalent forms, explains why they are equal, and why Chapter~\ref{ch:nbody} uses the third one. It is restricted to $l=2$; the general case follows by replacing $2$ by $l$.

\section{Variables}
All angles refer to the orbit of the star seen from the planet, measured in the planet's equatorial plane (Appendix~\ref{app:conventions}).
\begin{center}\small
\renewcommand{\arraystretch}{1.25}
\begin{tabular}{>{$}l<{$}p{12.3cm}}\toprule
\text{Symbol} & Meaning\\\midrule
\omega & argument of pericentre: angle from the ascending node to the pericentre, in the orbital plane\\
\M & mean anomaly; $\dot\M\simeq n$\\
\Omega & longitude of the ascending node of the star's orbit on the planet's equator\\
\thetap & rotation angle of the planet; $\dot\theta_{\mathrm p}=\Op$, the spin rate\\
\varpi=\Omega+\omega & longitude of pericentre (a "broken" angle: $\Omega$ in the equator, $\omega$ in the orbital plane)\\
\lambda=\Omega+\omega+\M & mean longitude of the star\\
n & mean motion, $\sqrt{\G(M^\star+M_{\mathrm p})/a^3}$\\
k=m-2+2p & integer that multiplies $\Omega$ in the second and third forms\\\bottomrule
\end{tabular}
\end{center}

\section{The three forms}
\paragraph{Kaula's form.} Kaula's angle at $l=2$, eq.~\eqref{eq:nu}, and its time derivative, eq.~\eqref{eq:sigma}, are
\[
  \nu_{2mpq}=(2-2p)\,\omega+(2-2p+q)\,\M+m\,(\Omega-\thetap),
\]
\begin{equation}
  \boxed{\sigma_{2mpq}=(2-2p+q)\,n-m\,\Op+(2-2p)\,\dot\omega+m\,\dot\Omega.} \label{eq:sigma-form1}
\end{equation}

\paragraph{Form with the longitude of pericentre.} Write $\omega=\varpi-\Omega$ in $\nu_{2mpq}$:
\[
  \nu_{2mpq}=(2-2p)\,\varpi+(2-2p+q)\,\M+(m-2+2p)\,\Omega-m\,\thetap .
\]
Hence, with $k=m-2+2p$,
\begin{equation}
  \boxed{\sigma_{2mpq}=(2-2p+q)\,n-m\,\Op+(2-2p)\,\dot\varpi+k\,\dot\Omega.} \label{eq:sigma-form2}
\end{equation}

\paragraph{Form with the mean longitude.} Group $(2-2p+q)\M+(2-2p)\varpi=(2-2p)(\varpi+\M)+q\M=(2-2p)\lambda+q\M$:
\[
  \nu_{2mpq}=(2-2p)\,\lambda+q\,\M+k\,\Omega-m\,\thetap .
\]
Hence
\begin{equation}
  \boxed{\sigma_{2mpq}=(2-2p)\,\dot\lambda+q\,n+k\,\dot\Omega-m\,\Op.} \label{eq:sigma-form3}
\end{equation}
Expanding $\dot\lambda=\dot\Omega+\dot\omega+n$ in~\eqref{eq:sigma-form3}, or $\dot\varpi=\dot\Omega+\dot\omega$ in~\eqref{eq:sigma-form2}, gives back~\eqref{eq:sigma-form1}: the three forms are identical, for all $(m,p,q)$. Equation~\eqref{eq:sigma-form3} is equation~\eqref{eq:sigma-regular} of Chapter~\ref{ch:nbody}.

\section{The coplanar form}
\paragraph{Which terms exist at $i=0$.} The integer $k$ and the order of $F_{2mp}(i)$ at small $i$ are given below:
\begin{center}\small
\setlength{\tabcolsep}{3pt}
\begin{tabular}{l|ccc|ccc|ccc}\toprule
$(m,p)$ & $(0,0)$ & $(0,1)$ & $(0,2)$ & $(1,0)$ & $(1,1)$ & $(1,2)$ & $(2,0)$ & $(2,1)$ & $(2,2)$\\\midrule
$k$ & $-2$ & $0$ & $2$ & $-1$ & $1$ & $3$ & $0$ & $2$ & $4$\\
$F_{2mp}(i)$ & $\mathcal O(i^2)$ & $-\tfrac12$ & $\mathcal O(i^2)$ & $\mathcal O(i)$ & $\mathcal O(i)$ & $\mathcal O(i^3)$ & $3$ & $\mathcal O(i^2)$ & $\mathcal O(i^4)$\\\bottomrule
\end{tabular}

\smallskip
($-\tfrac12$ and $3$ are the values at $i=0$; the corrections are $\mathcal O(i^2)$.)
\end{center}
In every case $F_{2mp}(i)=\mathcal O(i^{|k|})$. At $i=0$ only $(m,p)=(0,1)$ and $(2,0)$ remain, and for both $k=0$.

\paragraph{Coplanar frequencies.} With $k=0$, \eqref{eq:sigma-form2} becomes
\begin{equation}
  \sigma_{2mpq}=(2-2p+q)\,n-m\,\Op+(2-2p)\,\dot\varpi=(2-2p+q)\,n-m\,\Op+(2-2p)\,(\dot\Omega+\dot\omega), \label{eq:sigma-coplanar}
\end{equation}
that is, $\sigma_{201q}=q\,n$ and $\sigma_{220q}=(2+q)\,n-2\Op+2\dot\varpi$, as in~\eqref{eq:nu-coplanar}. This is the form commonly used in the coplanar case. It coincides with Kaula's form for the two coplanar terms.

\paragraph{Oblique planet.} When $i\neq0$, the seven other pairs $(m,p)$ appear, with $k\neq0$. The coplanar form~\eqref{eq:sigma-coplanar} then misses the term $k\,\dot\Omega$. For example, for $(m,p)=(1,1)$, $k=1$:
\[
  \text{Kaula:}\quad\sigma_{211q}=q\,n+\dot\Omega-\Op,\qquad \text{coplanar form:}\quad q\,n-\Op .
\]
The term $k\,\dot\Omega$ is small compared with $n$ and $\Op$ ($|\dot\Omega|\sim10^{-6}n$ in the 3:2 test of Section~\ref{sec:rates}). It matters when the other terms of $\sigma$ nearly cancel. This happens for an oblique planet in synchronous rotation: for example, with $\Op=n$ and $\dot{\bar\lambda}=n$, the term $(m,p,q)=(2,1,2)$ has $\sigma=2n+2\dot\Omega-2\Op=2\dot\Omega$, entirely given by the precession of the node. Its Love number is then evaluated at the precession frequency instead of $0$; this is the regime of Cassini states.

\section{Why the third form is used in an $N$-body code}
In a simulation, the elements are osculating and are recomputed from the state vectors at each step.
\begin{itemize}
  \item $\omega$ and $\varpi$ are undefined at $e=0$. Near $e=0$, the perturbations displace the pericentre by large angles: the osculating $\dot\omega$ and $\dot\varpi$ behave like $1/e$.
  \item $\Omega$ is undefined at $i=0$, and $\dot\Omega$ behaves like $1/\sin i$.
  \item $\lambda=\Omega+\omega+\M$ is defined for any $e$ and $i$: it differs from the true longitude of the star, $\Omega+\omega+f$, by $f-\M$, a periodic term of amplitude about $2e$. Its rate $\dot\lambda$ stays close to $n$.
\end{itemize}
Form~\eqref{eq:sigma-form3} contains $\dot\lambda$, $n$ and $\Op$, which are always defined, and $k\,\dot\Omega$, which is multiplied by $F_{2mp}(i)=\mathcal O(i^{|k|})$ in the force and therefore does not contribute as $i\to0$. It needs no special treatment at $e=0$ or $i=0$. Form~\eqref{eq:sigma-form1} needs $\dot\omega$, which cannot be estimated at small $e$; form~\eqref{eq:sigma-form2} needs $\dot\varpi$, which has the same problem.

\section{Mean rates}
The Love number is evaluated at the frequency of a harmonic forcing (H1). The rates in~\eqref{eq:sigma-form3} are therefore averaged over the fast oscillations (orbital and synodic periods) but not over the slow ones (precessions, resonant angles). A bar denotes this average, over a window $\Delta T$ computed by the code from the mean motions of the other bodies, eq.~\eqref{eq:window}: the longest synodic period of the planet, between one and ten orbital periods (Section~\ref{sec:rates}):
\begin{center}\small
\renewcommand{\arraystretch}{1.25}
\begin{tabular}{>{$}l<{$}p{7.2cm}p{5.6cm}}\toprule
\text{Symbol} & Meaning & How it is obtained\\\midrule
n & osculating mean motion $\sqrt{\G(M^\star+M_{\mathrm p})/a^3}$ & at each step, from $a$ (eq.~\ref{eq:elements-a-e})\\
\bar n & mean of $n$ over the window; replaces $\dot\M$ in $q\,n$ & mean of the stored values of $n$\\
\dot{\bar\lambda} & mean rate of the mean longitude $\lambda$; close to $n$, it contains the precession of the pericentre and of the node & $(\lambda_{\rm last}-\lambda_{\rm first})/\Delta T$, $\lambda$ accumulated step by step\\
\dot{\bar\Omega} & mean precession rate of the node of the star's orbit on the planet's equator, due to the other planets and to the precession of the spin axis; zero for an isolated Kepler orbit & $(\Omega_{\rm last}-\Omega_{\rm first})/\Delta T$, $\Omega$ accumulated from the rotation of $\hat{\mathbf N}$ about $\hat{\mathbf s}$\\
\Op & spin rate of the planet, $\|\mathbf s\|$ & at each step\\\bottomrule
\end{tabular}
\end{center}
With these mean rates,
\[
  \bar\sigma_{2mpq}=(2-2p)\,\dot{\bar\lambda}+q\,\bar n+k\,\dot{\bar\Omega}-m\,\Op,\qquad k=m-2+2p .
\]


\bibliographystyle{plainnat}
\bibliography{references}

\end{document}
